% ============================================================
%  structure-adv.tex -- THE advanced volume's set sequence
%
%  One list, read by main-pl-a4-adv.tex, so a set cannot exist in the tree and
%  be missing from the book. The block lists are themselves generated
%  (sets-adv/generated/_blok-N.tex) and carry the number each set must come up
%  under, so this file and tools/gen_plan.py cannot disagree about numbering.
% ============================================================

\chapter{Blok I — Systemy i reszty}

Dziesiątkowy jest jednym z wielu systemów i najmniej ciekawym. Tu liczysz
w dwójkowym i szesnastkowym, bierzesz reszty z dzielenia jako działanie samo
w sobie i poznajesz cyfrę kontrolną — czyli sposób sprawdzenia całego mnożenia
jednym dodawaniem.

Zestaw robisz w całości, na stoper, i zapisujesz czas pod spodem. Dwa wiersze
są celowo: do zestawu się wraca.

% GENERATED by tools/gen_sets.py -- do not edit.
\btexpect{1}% ======================================================
%  GENERATED by tools/gen_sets.py -- do not edit.
%  Seed 2001; regenerate with `make sets`.
% ======================================================

\begin{zestaw}{System dwójkowy}{2}{6:40}{3}
  \z{$155$ dwójkowo}{10011011}
  \z{$210$ dwójkowo}{11010010}
  \z{$10101011_2$}{171}
  \z{$11010010_2$}{210}
  \z{$10010110_2$}{150}
  \z{$10111000_2$}{184}
  \z{$163$ dwójkowo}{10100011}
  \z{$123$ dwójkowo}{1111011}
  \z{$169$ dwójkowo}{10101001}
  \z{$1011111_2$}{95}
  \z{$116$ dwójkowo}{1110100}
  \z{$222$ dwójkowo}{11011110}
  \z{$101$ dwójkowo}{1100101}
  \z{$46$ dwójkowo}{101110}
  \btnc
  \z{$1000110_2$}{70}
  \z{$15$ dwójkowo}{1111}
  \z{$1000100_2$}{68}
  \z{$126$ dwójkowo}{1111110}
  \z{$208$ dwójkowo}{11010000}
  \z{$138$ dwójkowo}{10001010}
  \z{$60$ dwójkowo}{111100}
  \z{$137$ dwójkowo}{10001001}
  \z{$239$ dwójkowo}{11101111}
  \z{$181$ dwójkowo}{10110101}
  \z{$111$ dwójkowo}{1101111}
  \z{$11010111_2$}{215}
  \z{$10110011_2$}{179}
  \btnc
  \z{$1100000_2$}{96}
  \z{$1100101_2$}{101}
  \z{$11101111_2$}{239}
  \z{$141$ dwójkowo}{10001101}
  \z{$29$ dwójkowo}{11101}
  \z{$198$ dwójkowo}{11000110}
  \z{$1001100_2$}{76}
  \z{$209$ dwójkowo}{11010001}
  \z{$110001_2$}{49}
  \z{$11001_2$}{25}
  \z{$10010100_2$}{148}
  \z{$119$ dwójkowo}{1110111}
  \z{$10110000_2$}{176}
\end{zestaw}

\btexpect{2}% ======================================================
%  GENERATED by tools/gen_sets.py -- do not edit.
%  Seed 2002; regenerate with `make sets`.
% ======================================================

\begin{zestaw}{Reszta modulo}{2}{5:20}{3}
  \z{$562 \bmod 3$}{1}
  \z{$846 \bmod 9$}{0}
  \z{$951 \bmod 4$}{3}
  \z{$956 \bmod 8$}{4}
  \z{$662 \bmod 8$}{6}
  \z{$165 \bmod 3$}{0}
  \z{$574 \bmod 8$}{6}
  \z{$127 \bmod 8$}{7}
  \z{$399 \bmod 3$}{0}
  \z{$948 \bmod 9$}{3}
  \z{$77 \bmod 9$}{5}
  \z{$65 \bmod 12$}{5}
  \z{$661 \bmod 6$}{1}
  \z{$848 \bmod 11$}{1}
  \btnc
  \z{$671 \bmod 9$}{5}
  \z{$581 \bmod 13$}{9}
  \z{$781 \bmod 11$}{0}
  \z{$501 \bmod 3$}{0}
  \z{$64 \bmod 13$}{12}
  \z{$408 \bmod 8$}{0}
  \z{$918 \bmod 13$}{8}
  \z{$660 \bmod 12$}{0}
  \z{$856 \bmod 3$}{1}
  \z{$317 \bmod 9$}{2}
  \z{$415 \bmod 11$}{8}
  \z{$689 \bmod 12$}{5}
  \z{$167 \bmod 3$}{2}
  \btnc
  \z{$773 \bmod 8$}{5}
  \z{$471 \bmod 11$}{9}
  \z{$495 \bmod 12$}{3}
  \z{$569 \bmod 4$}{1}
  \z{$250 \bmod 6$}{4}
  \z{$600 \bmod 6$}{0}
  \z{$588 \bmod 3$}{0}
  \z{$199 \bmod 8$}{7}
  \z{$288 \bmod 6$}{0}
  \z{$701 \bmod 13$}{12}
  \z{$857 \bmod 9$}{2}
  \z{$75 \bmod 8$}{3}
  \z{$91 \bmod 7$}{0}
\end{zestaw}

\btexpect{3}% ======================================================
%  GENERATED by tools/gen_sets.py -- do not edit.
%  Seed 2003; regenerate with `make sets`.
% ======================================================

\begin{zestaw}{System szesnastkowy}{2}{6:40}{3}
  \z{$\mathrm{2C}_{16}$}{44}
  \z{$\mathrm{9E}_{16}$}{158}
  \z{$88$ szesnastkowo}{58}
  \z{$\mathrm{76}_{16}$}{118}
  \z{$\mathrm{DD}_{16}$}{221}
  \z{$\mathrm{E6}_{16}$}{230}
  \z{$\mathrm{29}_{16}$}{41}
  \z{$102$ szesnastkowo}{66}
  \z{$214$ szesnastkowo}{D6}
  \z{$173$ szesnastkowo}{AD}
  \z{$135$ szesnastkowo}{87}
  \z{$23$ szesnastkowo}{17}
  \z{$\mathrm{85}_{16}$}{133}
  \z{$\mathrm{57}_{16}$}{87}
  \btnc
  \z{$147$ szesnastkowo}{93}
  \z{$227$ szesnastkowo}{E3}
  \z{$\mathrm{FE}_{16}$}{254}
  \z{$179$ szesnastkowo}{B3}
  \z{$232$ szesnastkowo}{E8}
  \z{$140$ szesnastkowo}{8C}
  \z{$251$ szesnastkowo}{FB}
  \z{$\mathrm{2E}_{16}$}{46}
  \z{$78$ szesnastkowo}{4E}
  \z{$\mathrm{13}_{16}$}{19}
  \z{$\mathrm{44}_{16}$}{68}
  \z{$183$ szesnastkowo}{B7}
  \z{$\mathrm{C2}_{16}$}{194}
  \btnc
  \z{$\mathrm{19}_{16}$}{25}
  \z{$162$ szesnastkowo}{A2}
  \z{$\mathrm{78}_{16}$}{120}
  \z{$\mathrm{C1}_{16}$}{193}
  \z{$43$ szesnastkowo}{2B}
  \z{$\mathrm{89}_{16}$}{137}
  \z{$178$ szesnastkowo}{B2}
  \z{$\mathrm{6A}_{16}$}{106}
  \z{$\mathrm{16}_{16}$}{22}
  \z{$184$ szesnastkowo}{B8}
  \z{$53$ szesnastkowo}{35}
  \z{$\mathrm{F8}_{16}$}{248}
  \z{$\mathrm{6D}_{16}$}{109}
\end{zestaw}

\btexpect{4}% ======================================================
%  GENERATED by tools/gen_sets.py -- do not edit.
%  Seed 2004; regenerate with `make sets`.
% ======================================================

\begin{zestaw}{Cyfra kontrolna}{2}{4:40}{2}
  \z{$551\,025$ --- cyfra kontrolna}{9}
  \z{$6\,413$ --- cyfra kontrolna}{5}
  \z{$828\,997$ --- cyfra kontrolna}{7}
  \z{$632\,399$ --- cyfra kontrolna}{5}
  \z{$217\,336$ --- cyfra kontrolna}{4}
  \z{$988\,607$ --- cyfra kontrolna}{2}
  \z{$389\,296$ --- cyfra kontrolna}{1}
  \z{$110\,486$ --- cyfra kontrolna}{2}
  \z{$701\,129$ --- cyfra kontrolna}{2}
  \z{$865\,428$ --- cyfra kontrolna}{6}
  \z{$641\,762$ --- cyfra kontrolna}{8}
  \z{$552\,762$ --- cyfra kontrolna}{9}
  \z{$631\,768$ --- cyfra kontrolna}{4}
  \z{$163\,036$ --- cyfra kontrolna}{1}
  \z{$949\,878$ --- cyfra kontrolna}{9}
  \z{$465\,232$ --- cyfra kontrolna}{4}
  \z{$167\,485$ --- cyfra kontrolna}{4}
  \z{$141\,968$ --- cyfra kontrolna}{2}
  \z{$545\,002$ --- cyfra kontrolna}{7}
  \z{$38\,276$ --- cyfra kontrolna}{8}
  \btnc
  \z{$939\,070$ --- cyfra kontrolna}{1}
  \z{$413\,005$ --- cyfra kontrolna}{4}
  \z{$837\,332$ --- cyfra kontrolna}{8}
  \z{$574\,056$ --- cyfra kontrolna}{9}
  \z{$696\,469$ --- cyfra kontrolna}{4}
  \z{$385\,023$ --- cyfra kontrolna}{3}
  \z{$933\,467$ --- cyfra kontrolna}{5}
  \z{$917\,528$ --- cyfra kontrolna}{5}
  \z{$453\,315$ --- cyfra kontrolna}{3}
  \z{$177\,726$ --- cyfra kontrolna}{3}
  \z{$646\,348$ --- cyfra kontrolna}{4}
  \z{$334\,695$ --- cyfra kontrolna}{3}
  \z{$62\,541$ --- cyfra kontrolna}{9}
  \z{$313\,364$ --- cyfra kontrolna}{2}
  \z{$658\,716$ --- cyfra kontrolna}{6}
  \z{$476\,275$ --- cyfra kontrolna}{4}
  \z{$602\,785$ --- cyfra kontrolna}{1}
  \z{$107\,646$ --- cyfra kontrolna}{6}
  \z{$363\,689$ --- cyfra kontrolna}{8}
  \z{$395\,272$ --- cyfra kontrolna}{1}
\end{zestaw}

\btexpect{5}% ======================================================
%  GENERATED by tools/gen_sets.py -- do not edit.
%  Seed 2005; regenerate with `make sets`.
% ======================================================

\begin{zestaw}{Podzielność przez 7, 11, 13}{2}{6:00}{2}
  \zz{czy $8\,806$ dzieli się przez $17$?}{tak}
  \zz{czy $5\,082$ dzieli się przez $11$?}{tak}
  \zz{czy $2\,555$ dzieli się przez $7$?}{tak}
  \zz{czy $2\,078$ dzieli się przez $17$?}{nie}
  \zz{czy $286$ dzieli się przez $13$?}{tak}
  \zz{czy $3\,536$ dzieli się przez $13$?}{tak}
  \zz{czy $3\,570$ dzieli się przez $7$?}{tak}
  \zz{czy $894$ dzieli się przez $11$?}{nie}
  \zz{czy $2\,366$ dzieli się przez $7$?}{tak}
  \zz{czy $147$ dzieli się przez $7$?}{tak}
  \zz{czy $4\,564$ dzieli się przez $17$?}{nie}
  \zz{czy $8\,602$ dzieli się przez $17$?}{tak}
  \zz{czy $7\,523$ dzieli się przez $7$?}{nie}
  \zz{czy $7\,912$ dzieli się przez $11$?}{nie}
  \zz{czy $5\,516$ dzieli się przez $11$?}{nie}
  \zz{czy $2\,774$ dzieli się przez $11$?}{nie}
  \zz{czy $6\,226$ dzieli się przez $11$?}{tak}
  \zz{czy $3\,096$ dzieli się przez $11$?}{nie}
  \zz{czy $5\,741$ dzieli się przez $13$?}{nie}
  \zz{czy $2\,095$ dzieli się przez $13$?}{nie}
  \btnc
  \zz{czy $168$ dzieli się przez $7$?}{tak}
  \zz{czy $1\,533$ dzieli się przez $7$?}{tak}
  \zz{czy $8\,174$ dzieli się przez $7$?}{nie}
  \zz{czy $1\,749$ dzieli się przez $11$?}{tak}
  \zz{czy $1\,586$ dzieli się przez $13$?}{tak}
  \zz{czy $3\,822$ dzieli się przez $11$?}{nie}
  \zz{czy $876$ dzieli się przez $17$?}{nie}
  \zz{czy $9\,174$ dzieli się przez $17$?}{nie}
  \zz{czy $3\,729$ dzieli się przez $11$?}{tak}
  \zz{czy $4\,012$ dzieli się przez $17$?}{tak}
  \zz{czy $4\,033$ dzieli się przez $11$?}{nie}
  \zz{czy $9\,928$ dzieli się przez $17$?}{tak}
  \zz{czy $9\,314$ dzieli się przez $13$?}{nie}
  \zz{czy $4\,536$ dzieli się przez $7$?}{tak}
  \zz{czy $949$ dzieli się przez $13$?}{tak}
  \zz{czy $958$ dzieli się przez $13$?}{nie}
  \zz{czy $9\,632$ dzieli się przez $13$?}{nie}
  \zz{czy $9\,911$ dzieli się przez $7$?}{nie}
  \zz{czy $7\,058$ dzieli się przez $7$?}{nie}
  \zz{czy $566$ dzieli się przez $17$?}{nie}
\end{zestaw}

\btexpect{6}% ======================================================
%  GENERATED by tools/gen_sets.py -- do not edit.
%  Seed 2006; regenerate with `make sets`.
% ======================================================

\begin{zestaw}{Dodawanie dwójkowe}{2}{7:20}{2}
  \z{$100110_2 + 11100_2$}{1000010}
  \z{$111_2 + 101100_2$}{110011}
  \z{$110_2 + 10111_2$}{11101}
  \z{$101100_2 + 100110_2$}{1010010}
  \z{$101100_2 + 111_2$}{110011}
  \z{$101100_2 + 10010_2$}{111110}
  \z{$10000_2 + 11111_2$}{101111}
  \z{$10101_2 + 1111_2$}{100100}
  \z{$1111_2 + 100110_2$}{110101}
  \z{$110100_2 + 11111_2$}{1010011}
  \z{$1111_2 + 110110_2$}{1000101}
  \z{$101110_2 + 111_2$}{110101}
  \z{$10000_2 + 11011_2$}{101011}
  \z{$10111_2 + 100101_2$}{111100}
  \z{$11010_2 + 10101_2$}{101111}
  \z{$100_2 + 100000_2$}{100100}
  \z{$10110_2 + 11100_2$}{110010}
  \z{$1011_2 + 110110_2$}{1000001}
  \z{$101110_2 + 1110_2$}{111100}
  \z{$11101_2 + 101100_2$}{1001001}
  \btnc
  \z{$111001_2 + 110011_2$}{1101100}
  \z{$11110_2 + 101101_2$}{1001011}
  \z{$110111_2 + 101001_2$}{1100000}
  \z{$101011_2 + 1111_2$}{111010}
  \z{$111011_2 + 1001_2$}{1000100}
  \z{$101101_2 + 1001_2$}{110110}
  \z{$100001_2 + 1110_2$}{101111}
  \z{$100101_2 + 110110_2$}{1011011}
  \z{$11100_2 + 101011_2$}{1000111}
  \z{$101001_2 + 101110_2$}{1010111}
  \z{$100100_2 + 111100_2$}{1100000}
  \z{$110_2 + 101110_2$}{110100}
  \z{$100101_2 + 100000_2$}{1000101}
  \z{$111001_2 + 101001_2$}{1100010}
  \z{$101010_2 + 101010_2$}{1010100}
  \z{$100001_2 + 101011_2$}{1001100}
  \z{$110_2 + 1111_2$}{10101}
  \z{$100100_2 + 111_2$}{101011}
  \z{$100000_2 + 10110_2$}{110110}
  \z{$11110_2 + 111000_2$}{1010110}
\end{zestaw}

\btexpect{7}% ======================================================
%  GENERATED by tools/gen_sets.py -- do not edit.
%  Seed 2007; regenerate with `make sets`.
% ======================================================

\begin{zestaw}{Potęga modulo}{2}{8:00}{3}
  \z{$5^{2} \bmod 11$}{3}
  \z{$7^{3} \bmod 13$}{5}
  \z{$4^{2} \bmod 5$}{1}
  \z{$9^{5} \bmod 9$}{0}
  \z{$6^{2} \bmod 5$}{1}
  \z{$6^{6} \bmod 9$}{0}
  \z{$7^{5} \bmod 11$}{10}
  \z{$5^{2} \bmod 5$}{0}
  \z{$4^{5} \bmod 7$}{2}
  \z{$9^{4} \bmod 13$}{9}
  \z{$2^{3} \bmod 5$}{3}
  \z{$4^{3} \bmod 7$}{1}
  \z{$8^{6} \bmod 5$}{4}
  \z{$8^{2} \bmod 7$}{1}
  \btnc
  \z{$2^{5} \bmod 9$}{5}
  \z{$5^{5} \bmod 11$}{1}
  \z{$6^{6} \bmod 5$}{1}
  \z{$5^{3} \bmod 13$}{8}
  \z{$4^{3} \bmod 11$}{9}
  \z{$7^{3} \bmod 5$}{3}
  \z{$3^{4} \bmod 5$}{1}
  \z{$3^{5} \bmod 7$}{5}
  \z{$6^{4} \bmod 7$}{1}
  \z{$7^{6} \bmod 13$}{12}
  \z{$8^{4} \bmod 13$}{1}
  \z{$7^{2} \bmod 9$}{4}
  \z{$3^{4} \bmod 7$}{4}
  \btnc
  \z{$5^{6} \bmod 11$}{5}
  \z{$6^{5} \bmod 11$}{10}
  \z{$6^{5} \bmod 7$}{6}
  \z{$7^{6} \bmod 11$}{4}
  \z{$5^{6} \bmod 7$}{1}
  \z{$7^{3} \bmod 7$}{0}
  \z{$2^{4} \bmod 7$}{2}
  \z{$4^{5} \bmod 13$}{10}
  \z{$4^{4} \bmod 13$}{9}
  \z{$5^{4} \bmod 7$}{2}
  \z{$4^{4} \bmod 7$}{4}
  \z{$9^{2} \bmod 13$}{3}
  \z{$6^{3} \bmod 13$}{8}
\end{zestaw}

\btexpect{8}% ======================================================
%  GENERATED by tools/gen_sets.py -- do not edit.
%  Seed 2008; regenerate with `make sets`.
% ======================================================

\begin{zestaw}{System dwójkowy II}{2}{6:40}{3}
  \z{$232$ dwójkowo}{11101000}
  \z{$228$ dwójkowo}{11100100}
  \z{$10001_2$}{17}
  \z{$106$ dwójkowo}{1101010}
  \z{$11101110_2$}{238}
  \z{$101011_2$}{43}
  \z{$11101011_2$}{235}
  \z{$64$ dwójkowo}{1000000}
  \z{$11000010_2$}{194}
  \z{$1100000_2$}{96}
  \z{$205$ dwójkowo}{11001101}
  \z{$43$ dwójkowo}{101011}
  \z{$10010110_2$}{150}
  \z{$10111000_2$}{184}
  \btnc
  \z{$10100100_2$}{164}
  \z{$11010001_2$}{209}
  \z{$247$ dwójkowo}{11110111}
  \z{$11011011_2$}{219}
  \z{$1000111_2$}{71}
  \z{$11111101_2$}{253}
  \z{$210$ dwójkowo}{11010010}
  \z{$231$ dwójkowo}{11100111}
  \z{$150$ dwójkowo}{10010110}
  \z{$1011111_2$}{95}
  \z{$1110010_2$}{114}
  \z{$36$ dwójkowo}{100100}
  \z{$1100011_2$}{99}
  \btnc
  \z{$78$ dwójkowo}{1001110}
  \z{$101000_2$}{40}
  \z{$86$ dwójkowo}{1010110}
  \z{$1110100_2$}{116}
  \z{$1000000_2$}{64}
  \z{$253$ dwójkowo}{11111101}
  \z{$11100100_2$}{228}
  \z{$179$ dwójkowo}{10110011}
  \z{$10110011_2$}{179}
  \z{$1010000_2$}{80}
  \z{$1111010_2$}{122}
  \z{$10101010_2$}{170}
  \z{$105$ dwójkowo}{1101001}
\end{zestaw}

\btexpect{9}% ======================================================
%  GENERATED by tools/gen_sets.py -- do not edit.
%  Seed 2009; regenerate with `make sets`.
% ======================================================

\begin{zestaw}{Reszta modulo II}{2}{5:20}{3}
  \z{$694 \bmod 8$}{6}
  \z{$473 \bmod 6$}{5}
  \z{$784 \bmod 3$}{1}
  \z{$230 \bmod 9$}{5}
  \z{$683 \bmod 8$}{3}
  \z{$208 \bmod 4$}{0}
  \z{$265 \bmod 7$}{6}
  \z{$978 \bmod 6$}{0}
  \z{$947 \bmod 3$}{2}
  \z{$956 \bmod 13$}{7}
  \z{$554 \bmod 3$}{2}
  \z{$667 \bmod 13$}{4}
  \z{$690 \bmod 12$}{6}
  \z{$685 \bmod 4$}{1}
  \btnc
  \z{$286 \bmod 13$}{0}
  \z{$806 \bmod 11$}{3}
  \z{$379 \bmod 4$}{3}
  \z{$507 \bmod 6$}{3}
  \z{$629 \bmod 13$}{5}
  \z{$114 \bmod 7$}{2}
  \z{$55 \bmod 3$}{1}
  \z{$622 \bmod 7$}{6}
  \z{$288 \bmod 11$}{2}
  \z{$465 \bmod 13$}{10}
  \z{$732 \bmod 9$}{3}
  \z{$889 \bmod 4$}{1}
  \z{$805 \bmod 11$}{2}
  \btnc
  \z{$662 \bmod 9$}{5}
  \z{$352 \bmod 8$}{0}
  \z{$380 \bmod 7$}{2}
  \z{$762 \bmod 3$}{0}
  \z{$233 \bmod 8$}{1}
  \z{$370 \bmod 13$}{6}
  \z{$704 \bmod 12$}{8}
  \z{$475 \bmod 4$}{3}
  \z{$728 \bmod 3$}{2}
  \z{$501 \bmod 13$}{7}
  \z{$601 \bmod 8$}{1}
  \z{$406 \bmod 11$}{10}
  \z{$790 \bmod 8$}{6}
\end{zestaw}

\btexpect{10}% ======================================================
%  GENERATED by tools/gen_sets.py -- do not edit.
%  Seed 2010; regenerate with `make sets`.
% ======================================================

\begin{zestaw}{System szesnastkowy II}{2}{6:40}{3}
  \z{$\mathrm{33}_{16}$}{51}
  \z{$132$ szesnastkowo}{84}
  \z{$155$ szesnastkowo}{9B}
  \z{$\mathrm{ED}_{16}$}{237}
  \z{$\mathrm{51}_{16}$}{81}
  \z{$70$ szesnastkowo}{46}
  \z{$209$ szesnastkowo}{D1}
  \z{$\mathrm{EB}_{16}$}{235}
  \z{$\mathrm{76}_{16}$}{118}
  \z{$\mathrm{30}_{16}$}{48}
  \z{$142$ szesnastkowo}{8E}
  \z{$78$ szesnastkowo}{4E}
  \z{$24$ szesnastkowo}{18}
  \z{$120$ szesnastkowo}{78}
  \btnc
  \z{$\mathrm{B7}_{16}$}{183}
  \z{$153$ szesnastkowo}{99}
  \z{$\mathrm{3E}_{16}$}{62}
  \z{$\mathrm{71}_{16}$}{113}
  \z{$36$ szesnastkowo}{24}
  \z{$43$ szesnastkowo}{2B}
  \z{$\mathrm{24}_{16}$}{36}
  \z{$104$ szesnastkowo}{68}
  \z{$96$ szesnastkowo}{60}
  \z{$\mathrm{FF}_{16}$}{255}
  \z{$\mathrm{6D}_{16}$}{109}
  \z{$\mathrm{91}_{16}$}{145}
  \z{$\mathrm{14}_{16}$}{20}
  \btnc
  \z{$223$ szesnastkowo}{DF}
  \z{$\mathrm{26}_{16}$}{38}
  \z{$42$ szesnastkowo}{2A}
  \z{$\mathrm{A4}_{16}$}{164}
  \z{$\mathrm{1A}_{16}$}{26}
  \z{$118$ szesnastkowo}{76}
  \z{$22$ szesnastkowo}{16}
  \z{$149$ szesnastkowo}{95}
  \z{$\mathrm{46}_{16}$}{70}
  \z{$102$ szesnastkowo}{66}
  \z{$\mathrm{3B}_{16}$}{59}
  \z{$\mathrm{45}_{16}$}{69}
  \z{$\mathrm{27}_{16}$}{39}
\end{zestaw}

\btexpect{11}% ======================================================
%  GENERATED by tools/gen_sets.py -- do not edit.
%  Seed 2011; regenerate with `make sets`.
% ======================================================

\begin{zestaw}{Cyfra kontrolna II}{2}{4:40}{2}
  \z{$544\,550$ --- cyfra kontrolna}{5}
  \z{$570\,007$ --- cyfra kontrolna}{1}
  \z{$953\,414$ --- cyfra kontrolna}{8}
  \z{$656\,750$ --- cyfra kontrolna}{2}
  \z{$261\,586$ --- cyfra kontrolna}{1}
  \z{$847\,280$ --- cyfra kontrolna}{2}
  \z{$324\,872$ --- cyfra kontrolna}{8}
  \z{$176\,947$ --- cyfra kontrolna}{7}
  \z{$311\,581$ --- cyfra kontrolna}{1}
  \z{$951\,896$ --- cyfra kontrolna}{2}
  \z{$59\,834$ --- cyfra kontrolna}{2}
  \z{$688\,502$ --- cyfra kontrolna}{2}
  \z{$377\,856$ --- cyfra kontrolna}{9}
  \z{$19\,434$ --- cyfra kontrolna}{3}
  \z{$308\,319$ --- cyfra kontrolna}{6}
  \z{$952\,061$ --- cyfra kontrolna}{5}
  \z{$227\,871$ --- cyfra kontrolna}{9}
  \z{$881\,834$ --- cyfra kontrolna}{5}
  \z{$233\,904$ --- cyfra kontrolna}{3}
  \z{$725\,286$ --- cyfra kontrolna}{3}
  \btnc
  \z{$41\,047$ --- cyfra kontrolna}{7}
  \z{$953\,673$ --- cyfra kontrolna}{6}
  \z{$50\,808$ --- cyfra kontrolna}{3}
  \z{$182\,663$ --- cyfra kontrolna}{8}
  \z{$340\,213$ --- cyfra kontrolna}{4}
  \z{$426\,795$ --- cyfra kontrolna}{6}
  \z{$988\,999$ --- cyfra kontrolna}{7}
  \z{$537\,441$ --- cyfra kontrolna}{6}
  \z{$955\,842$ --- cyfra kontrolna}{6}
  \z{$483\,623$ --- cyfra kontrolna}{8}
  \z{$635\,855$ --- cyfra kontrolna}{5}
  \z{$21\,420$ --- cyfra kontrolna}{9}
  \z{$117\,868$ --- cyfra kontrolna}{4}
  \z{$696\,307$ --- cyfra kontrolna}{4}
  \z{$305\,501$ --- cyfra kontrolna}{5}
  \z{$972\,421$ --- cyfra kontrolna}{7}
  \z{$662\,705$ --- cyfra kontrolna}{8}
  \z{$400\,559$ --- cyfra kontrolna}{5}
  \z{$638\,849$ --- cyfra kontrolna}{2}
  \z{$872\,496$ --- cyfra kontrolna}{9}
\end{zestaw}

\btexpect{12}% ======================================================
%  GENERATED by tools/gen_sets.py -- do not edit.
%  Seed 2012; regenerate with `make sets`.
% ======================================================

\begin{zestaw}{Podzielność przez 7, 11, 13 II}{2}{6:00}{2}
  \zz{czy $1\,533$ dzieli się przez $7$?}{tak}
  \zz{czy $8\,619$ dzieli się przez $13$?}{tak}
  \zz{czy $1\,677$ dzieli się przez $13$?}{tak}
  \zz{czy $3\,609$ dzieli się przez $13$?}{nie}
  \zz{czy $3\,196$ dzieli się przez $17$?}{tak}
  \zz{czy $9\,859$ dzieli się przez $17$?}{nie}
  \zz{czy $9\,508$ dzieli się przez $13$?}{nie}
  \zz{czy $200$ dzieli się przez $13$?}{nie}
  \zz{czy $7\,216$ dzieli się przez $11$?}{tak}
  \zz{czy $1\,372$ dzieli się przez $7$?}{tak}
  \zz{czy $4\,966$ dzieli się przez $13$?}{tak}
  \zz{czy $5\,775$ dzieli się przez $11$?}{tak}
  \zz{czy $4\,395$ dzieli się przez $7$?}{nie}
  \zz{czy $5\,951$ dzieli się przez $11$?}{tak}
  \zz{czy $2\,624$ dzieli się przez $13$?}{nie}
  \zz{czy $5\,929$ dzieli się przez $17$?}{nie}
  \zz{czy $3\,465$ dzieli się przez $7$?}{tak}
  \zz{czy $2\,349$ dzieli się przez $13$?}{nie}
  \zz{czy $8\,500$ dzieli się przez $17$?}{tak}
  \zz{czy $1\,394$ dzieli się przez $17$?}{tak}
  \btnc
  \zz{czy $5\,657$ dzieli się przez $7$?}{nie}
  \zz{czy $3\,955$ dzieli się przez $7$?}{tak}
  \zz{czy $2\,124$ dzieli się przez $7$?}{nie}
  \zz{czy $11\,441$ dzieli się przez $17$?}{tak}
  \zz{czy $4\,655$ dzieli się przez $7$?}{tak}
  \zz{czy $6\,876$ dzieli się przez $13$?}{nie}
  \zz{czy $3\,322$ dzieli się przez $11$?}{tak}
  \zz{czy $592$ dzieli się przez $11$?}{nie}
  \zz{czy $1\,291$ dzieli się przez $13$?}{nie}
  \zz{czy $2\,499$ dzieli się przez $17$?}{tak}
  \zz{czy $597$ dzieli się przez $13$?}{nie}
  \zz{czy $203$ dzieli się przez $11$?}{nie}
  \zz{czy $1\,693$ dzieli się przez $13$?}{nie}
  \zz{czy $3\,895$ dzieli się przez $7$?}{nie}
  \zz{czy $1\,727$ dzieli się przez $11$?}{tak}
  \zz{czy $3\,570$ dzieli się przez $17$?}{tak}
  \zz{czy $9\,044$ dzieli się przez $17$?}{tak}
  \zz{czy $9\,238$ dzieli się przez $17$?}{nie}
  \zz{czy $2\,024$ dzieli się przez $11$?}{tak}
  \zz{czy $353$ dzieli się przez $17$?}{nie}
\end{zestaw}

\btexpect{13}% ======================================================
%  GENERATED by tools/gen_sets.py -- do not edit.
%  Seed 2013; regenerate with `make sets`.
% ======================================================

\begin{zestaw}{Dodawanie dwójkowe II}{2}{7:20}{2}
  \z{$11110_2 + 111000_2$}{1010110}
  \z{$110_2 + 111001_2$}{111111}
  \z{$11100_2 + 110110_2$}{1010010}
  \z{$111100_2 + 11000_2$}{1010100}
  \z{$110011_2 + 111000_2$}{1101011}
  \z{$11011_2 + 100111_2$}{1000010}
  \z{$100011_2 + 11111_2$}{1000010}
  \z{$11101_2 + 110100_2$}{1010001}
  \z{$100110_2 + 110010_2$}{1011000}
  \z{$100111_2 + 110100_2$}{1011011}
  \z{$10100_2 + 111100_2$}{1010000}
  \z{$101101_2 + 110_2$}{110011}
  \z{$1001_2 + 11100_2$}{100101}
  \z{$10111_2 + 1000_2$}{11111}
  \z{$110100_2 + 110111_2$}{1101011}
  \z{$10010_2 + 101111_2$}{1000001}
  \z{$100000_2 + 100110_2$}{1000110}
  \z{$1011_2 + 111011_2$}{1000110}
  \z{$101111_2 + 10110_2$}{1000101}
  \z{$110010_2 + 11111_2$}{1010001}
  \btnc
  \z{$100010_2 + 10110_2$}{111000}
  \z{$10100_2 + 101111_2$}{1000011}
  \z{$1110_2 + 111_2$}{10101}
  \z{$111000_2 + 110110_2$}{1101110}
  \z{$101000_2 + 10101_2$}{111101}
  \z{$110001_2 + 101000_2$}{1011001}
  \z{$10000_2 + 1010_2$}{11010}
  \z{$1011_2 + 101000_2$}{110011}
  \z{$11000_2 + 100110_2$}{111110}
  \z{$11011_2 + 100000_2$}{111011}
  \z{$110110_2 + 11000_2$}{1001110}
  \z{$101011_2 + 101011_2$}{1010110}
  \z{$101100_2 + 11001_2$}{1000101}
  \z{$110100_2 + 11011_2$}{1001111}
  \z{$10010_2 + 110001_2$}{1000011}
  \z{$101111_2 + 111011_2$}{1101010}
  \z{$11_2 + 100111_2$}{101010}
  \z{$10010_2 + 100011_2$}{110101}
  \z{$100100_2 + 1110_2$}{110010}
  \z{$1101_2 + 110111_2$}{1000100}
\end{zestaw}

\btexpect{14}% ======================================================
%  GENERATED by tools/gen_sets.py -- do not edit.
%  Seed 2014; regenerate with `make sets`.
% ======================================================

\begin{zestaw}{Potęga modulo II}{2}{8:00}{3}
  \z{$7^{6} \bmod 11$}{4}
  \z{$5^{4} \bmod 7$}{2}
  \z{$8^{5} \bmod 11$}{10}
  \z{$2^{6} \bmod 11$}{9}
  \z{$7^{3} \bmod 13$}{5}
  \z{$9^{3} \bmod 11$}{3}
  \z{$5^{4} \bmod 5$}{0}
  \z{$3^{3} \bmod 7$}{6}
  \z{$5^{4} \bmod 11$}{9}
  \z{$8^{5} \bmod 9$}{8}
  \z{$3^{3} \bmod 9$}{0}
  \z{$2^{5} \bmod 11$}{10}
  \z{$6^{6} \bmod 5$}{1}
  \z{$3^{6} \bmod 5$}{4}
  \btnc
  \z{$5^{3} \bmod 5$}{0}
  \z{$5^{4} \bmod 9$}{4}
  \z{$4^{6} \bmod 7$}{1}
  \z{$7^{3} \bmod 5$}{3}
  \z{$2^{3} \bmod 5$}{3}
  \z{$3^{2} \bmod 9$}{0}
  \z{$5^{6} \bmod 7$}{1}
  \z{$8^{2} \bmod 11$}{9}
  \z{$2^{5} \bmod 7$}{4}
  \z{$8^{6} \bmod 11$}{3}
  \z{$6^{5} \bmod 9$}{0}
  \z{$3^{4} \bmod 9$}{0}
  \z{$2^{3} \bmod 9$}{8}
  \btnc
  \z{$7^{4} \bmod 11$}{3}
  \z{$9^{5} \bmod 9$}{0}
  \z{$2^{5} \bmod 9$}{5}
  \z{$8^{2} \bmod 7$}{1}
  \z{$4^{6} \bmod 5$}{1}
  \z{$4^{4} \bmod 7$}{4}
  \z{$3^{5} \bmod 9$}{0}
  \z{$4^{3} \bmod 5$}{4}
  \z{$6^{4} \bmod 9$}{0}
  \z{$5^{2} \bmod 5$}{0}
  \z{$7^{6} \bmod 7$}{0}
  \z{$6^{4} \bmod 5$}{1}
  \z{$4^{2} \bmod 7$}{2}
\end{zestaw}

\btexpect{15}% ======================================================
%  GENERATED by tools/gen_sets.py -- do not edit.
%  Seed 2015; regenerate with `make sets`.
% ======================================================

\begin{zestaw}{System dwójkowy III}{2}{6:40}{3}
  \z{$191$ dwójkowo}{10111111}
  \z{$10101011_2$}{171}
  \z{$100110_2$}{38}
  \z{$10101100_2$}{172}
  \z{$11001001_2$}{201}
  \z{$10111100_2$}{188}
  \z{$1110011_2$}{115}
  \z{$1101_2$}{13}
  \z{$136$ dwójkowo}{10001000}
  \z{$80$ dwójkowo}{1010000}
  \z{$10101110_2$}{174}
  \z{$10100111_2$}{167}
  \z{$240$ dwójkowo}{11110000}
  \z{$105$ dwójkowo}{1101001}
  \btnc
  \z{$10010001_2$}{145}
  \z{$231$ dwójkowo}{11100111}
  \z{$43$ dwójkowo}{101011}
  \z{$178$ dwójkowo}{10110010}
  \z{$127$ dwójkowo}{1111111}
  \z{$52$ dwójkowo}{110100}
  \z{$247$ dwójkowo}{11110111}
  \z{$100$ dwójkowo}{1100100}
  \z{$10110100_2$}{180}
  \z{$32$ dwójkowo}{100000}
  \z{$11111001_2$}{249}
  \z{$1000_2$}{8}
  \z{$11101011_2$}{235}
  \btnc
  \z{$142$ dwójkowo}{10001110}
  \z{$47$ dwójkowo}{101111}
  \z{$133$ dwójkowo}{10000101}
  \z{$216$ dwójkowo}{11011000}
  \z{$11000001_2$}{193}
  \z{$192$ dwójkowo}{11000000}
  \z{$138$ dwójkowo}{10001010}
  \z{$11101100_2$}{236}
  \z{$11011000_2$}{216}
  \z{$1101001_2$}{105}
  \z{$11101111_2$}{239}
  \z{$255$ dwójkowo}{11111111}
  \z{$64$ dwójkowo}{1000000}
\end{zestaw}

\btexpect{16}% ======================================================
%  GENERATED by tools/gen_sets.py -- do not edit.
%  Seed 2016; regenerate with `make sets`.
% ======================================================

\begin{zestaw}{Reszta modulo III}{2}{5:20}{3}
  \z{$963 \bmod 12$}{3}
  \z{$946 \bmod 13$}{10}
  \z{$956 \bmod 8$}{4}
  \z{$886 \bmod 4$}{2}
  \z{$180 \bmod 8$}{4}
  \z{$548 \bmod 4$}{0}
  \z{$147 \bmod 12$}{3}
  \z{$187 \bmod 8$}{3}
  \z{$379 \bmod 6$}{1}
  \z{$409 \bmod 7$}{3}
  \z{$294 \bmod 7$}{0}
  \z{$557 \bmod 6$}{5}
  \z{$996 \bmod 12$}{0}
  \z{$133 \bmod 9$}{7}
  \btnc
  \z{$298 \bmod 13$}{12}
  \z{$672 \bmod 8$}{0}
  \z{$355 \bmod 3$}{1}
  \z{$486 \bmod 7$}{3}
  \z{$66 \bmod 9$}{3}
  \z{$124 \bmod 11$}{3}
  \z{$832 \bmod 8$}{0}
  \z{$149 \bmod 12$}{5}
  \z{$592 \bmod 6$}{4}
  \z{$438 \bmod 3$}{0}
  \z{$934 \bmod 7$}{3}
  \z{$858 \bmod 9$}{3}
  \z{$415 \bmod 7$}{2}
  \btnc
  \z{$449 \bmod 3$}{2}
  \z{$113 \bmod 13$}{9}
  \z{$549 \bmod 8$}{5}
  \z{$748 \bmod 9$}{1}
  \z{$237 \bmod 9$}{3}
  \z{$766 \bmod 8$}{6}
  \z{$684 \bmod 3$}{0}
  \z{$650 \bmod 3$}{2}
  \z{$609 \bmod 12$}{9}
  \z{$315 \bmod 6$}{3}
  \z{$619 \bmod 7$}{3}
  \z{$582 \bmod 12$}{6}
  \z{$976 \bmod 3$}{1}
\end{zestaw}

\btexpect{17}% ======================================================
%  GENERATED by tools/gen_sets.py -- do not edit.
%  Seed 2017; regenerate with `make sets`.
% ======================================================

\begin{zestaw}{System szesnastkowy III}{2}{6:40}{3}
  \z{$66$ szesnastkowo}{42}
  \z{$218$ szesnastkowo}{DA}
  \z{$146$ szesnastkowo}{92}
  \z{$\mathrm{21}_{16}$}{33}
  \z{$\mathrm{85}_{16}$}{133}
  \z{$45$ szesnastkowo}{2D}
  \z{$169$ szesnastkowo}{A9}
  \z{$\mathrm{96}_{16}$}{150}
  \z{$\mathrm{11}_{16}$}{17}
  \z{$234$ szesnastkowo}{EA}
  \z{$\mathrm{DB}_{16}$}{219}
  \z{$235$ szesnastkowo}{EB}
  \z{$\mathrm{56}_{16}$}{86}
  \z{$\mathrm{C4}_{16}$}{196}
  \btnc
  \z{$\mathrm{FD}_{16}$}{253}
  \z{$\mathrm{73}_{16}$}{115}
  \z{$\mathrm{E7}_{16}$}{231}
  \z{$208$ szesnastkowo}{D0}
  \z{$\mathrm{BE}_{16}$}{190}
  \z{$197$ szesnastkowo}{C5}
  \z{$\mathrm{DD}_{16}$}{221}
  \z{$\mathrm{6A}_{16}$}{106}
  \z{$\mathrm{20}_{16}$}{32}
  \z{$\mathrm{2D}_{16}$}{45}
  \z{$\mathrm{1D}_{16}$}{29}
  \z{$\mathrm{4D}_{16}$}{77}
  \z{$\mathrm{CB}_{16}$}{203}
  \btnc
  \z{$33$ szesnastkowo}{21}
  \z{$155$ szesnastkowo}{9B}
  \z{$\mathrm{7B}_{16}$}{123}
  \z{$86$ szesnastkowo}{56}
  \z{$\mathrm{E3}_{16}$}{227}
  \z{$\mathrm{23}_{16}$}{35}
  \z{$151$ szesnastkowo}{97}
  \z{$223$ szesnastkowo}{DF}
  \z{$55$ szesnastkowo}{37}
  \z{$\mathrm{71}_{16}$}{113}
  \z{$\mathrm{A5}_{16}$}{165}
  \z{$193$ szesnastkowo}{C1}
  \z{$\mathrm{19}_{16}$}{25}
\end{zestaw}

\btexpect{18}% ======================================================
%  GENERATED by tools/gen_sets.py -- do not edit.
%  Seed 2018; regenerate with `make sets`.
% ======================================================

\begin{zestaw}{Cyfra kontrolna III}{2}{4:40}{2}
  \z{$560\,056$ --- cyfra kontrolna}{4}
  \z{$134\,610$ --- cyfra kontrolna}{6}
  \z{$30\,034$ --- cyfra kontrolna}{1}
  \z{$424\,372$ --- cyfra kontrolna}{4}
  \z{$658\,973$ --- cyfra kontrolna}{2}
  \z{$703\,952$ --- cyfra kontrolna}{8}
  \z{$532\,838$ --- cyfra kontrolna}{2}
  \z{$313\,533$ --- cyfra kontrolna}{9}
  \z{$127\,285$ --- cyfra kontrolna}{7}
  \z{$456\,542$ --- cyfra kontrolna}{8}
  \z{$561\,104$ --- cyfra kontrolna}{8}
  \z{$576\,111$ --- cyfra kontrolna}{3}
  \z{$248\,098$ --- cyfra kontrolna}{4}
  \z{$247\,812$ --- cyfra kontrolna}{6}
  \z{$37\,645$ --- cyfra kontrolna}{7}
  \z{$44\,074$ --- cyfra kontrolna}{1}
  \z{$905\,754$ --- cyfra kontrolna}{3}
  \z{$68\,150$ --- cyfra kontrolna}{2}
  \z{$456\,452$ --- cyfra kontrolna}{8}
  \z{$670\,316$ --- cyfra kontrolna}{5}
  \btnc
  \z{$554\,845$ --- cyfra kontrolna}{4}
  \z{$36\,737$ --- cyfra kontrolna}{8}
  \z{$482\,179$ --- cyfra kontrolna}{4}
  \z{$810\,363$ --- cyfra kontrolna}{3}
  \z{$670\,975$ --- cyfra kontrolna}{7}
  \z{$634\,609$ --- cyfra kontrolna}{1}
  \z{$804\,197$ --- cyfra kontrolna}{2}
  \z{$24\,851$ --- cyfra kontrolna}{2}
  \z{$83\,813$ --- cyfra kontrolna}{5}
  \z{$786\,746$ --- cyfra kontrolna}{2}
  \z{$173\,658$ --- cyfra kontrolna}{3}
  \z{$238\,556$ --- cyfra kontrolna}{2}
  \z{$212\,848$ --- cyfra kontrolna}{7}
  \z{$990\,831$ --- cyfra kontrolna}{3}
  \z{$451\,868$ --- cyfra kontrolna}{5}
  \z{$712\,277$ --- cyfra kontrolna}{8}
  \z{$617\,669$ --- cyfra kontrolna}{8}
  \z{$646\,122$ --- cyfra kontrolna}{3}
  \z{$355\,252$ --- cyfra kontrolna}{4}
  \z{$640\,965$ --- cyfra kontrolna}{3}
\end{zestaw}

\btexpect{19}% ======================================================
%  GENERATED by tools/gen_sets.py -- do not edit.
%  Seed 2019; regenerate with `make sets`.
% ======================================================

\begin{zestaw}{Mieszane — systemy}{2}{6:40}{3}
  \z{$218$ dwójkowo}{11011010}
  \z{$1000011_2$}{67}
  \z{$46$ dwójkowo}{101110}
  \z{$1000100_2$}{68}
  \z{$1010110_2$}{86}
  \z{$11010000_2$}{208}
  \z{$10101000_2$}{168}
  \z{$10011110_2$}{158}
  \z{$1101110_2$}{110}
  \z{$11100011_2$}{227}
  \z{$133$ szesnastkowo}{85}
  \z{$186$ szesnastkowo}{BA}
  \z{$24$ szesnastkowo}{18}
  \z{$102$ szesnastkowo}{66}
  \btnc
  \z{$\mathrm{44}_{16}$}{68}
  \z{$160$ szesnastkowo}{A0}
  \z{$159$ szesnastkowo}{9F}
  \z{$129$ szesnastkowo}{81}
  \z{$\mathrm{3F}_{16}$}{63}
  \z{$\mathrm{1B}_{16}$}{27}
  \z{$114 \bmod 3$}{0}
  \z{$381 \bmod 11$}{7}
  \z{$672 \bmod 12$}{0}
  \z{$351 \bmod 11$}{10}
  \z{$237 \bmod 11$}{6}
  \z{$771 \bmod 8$}{3}
  \z{$811 \bmod 7$}{6}
  \btnc
  \z{$836 \bmod 8$}{4}
  \z{$642 \bmod 4$}{2}
  \z{$562 \bmod 8$}{2}
  \z{$198\,972$ --- cyfra kontrolna}{9}
  \z{$356\,581$ --- cyfra kontrolna}{1}
  \z{$24\,184$ --- cyfra kontrolna}{1}
  \z{$255\,691$ --- cyfra kontrolna}{1}
  \z{$312\,146$ --- cyfra kontrolna}{8}
  \z{$219\,813$ --- cyfra kontrolna}{6}
  \z{$787\,312$ --- cyfra kontrolna}{1}
  \z{$610\,470$ --- cyfra kontrolna}{9}
  \z{$355\,316$ --- cyfra kontrolna}{5}
  \z{$3\,891$ --- cyfra kontrolna}{3}
\end{zestaw}

\btexpect{20}% ======================================================
%  GENERATED by tools/gen_sets.py -- do not edit.
%  Seed 2020; regenerate with `make sets`.
% ======================================================

\begin{zestaw}{Podzielność przez 7, 11, 13 III}{2}{6:00}{2}
  \zz{czy $7\,753$ dzieli się przez $11$?}{nie}
  \zz{czy $7\,633$ dzieli się przez $17$?}{tak}
  \zz{czy $7\,700$ dzieli się przez $11$?}{tak}
  \zz{czy $6\,913$ dzieli się przez $11$?}{nie}
  \zz{czy $7\,924$ dzieli się przez $11$?}{nie}
  \zz{czy $3\,145$ dzieli się przez $17$?}{tak}
  \zz{czy $9\,005$ dzieli się przez $17$?}{nie}
  \zz{czy $7\,454$ dzieli się przez $17$?}{nie}
  \zz{czy $5\,851$ dzieli się przez $7$?}{nie}
  \zz{czy $6\,204$ dzieli się przez $11$?}{tak}
  \zz{czy $4\,508$ dzieli się przez $7$?}{tak}
  \zz{czy $6\,292$ dzieli się przez $11$?}{tak}
  \zz{czy $4\,342$ dzieli się przez $13$?}{tak}
  \zz{czy $1\,037$ dzieli się przez $17$?}{tak}
  \zz{czy $9\,524$ dzieli się przez $11$?}{nie}
  \zz{czy $2\,299$ dzieli się przez $13$?}{nie}
  \zz{czy $5\,389$ dzieli się przez $17$?}{tak}
  \zz{czy $5\,922$ dzieli się przez $11$?}{nie}
  \zz{czy $1\,820$ dzieli się przez $7$?}{tak}
  \zz{czy $9\,076$ dzieli się przez $11$?}{nie}
  \btnc
  \zz{czy $2\,478$ dzieli się przez $7$?}{tak}
  \zz{czy $4\,186$ dzieli się przez $7$?}{tak}
  \zz{czy $1\,898$ dzieli się przez $17$?}{nie}
  \zz{czy $3\,935$ dzieli się przez $17$?}{nie}
  \zz{czy $5\,763$ dzieli się przez $17$?}{tak}
  \zz{czy $5\,797$ dzieli się przez $17$?}{tak}
  \zz{czy $6\,604$ dzieli się przez $13$?}{tak}
  \zz{czy $4\,348$ dzieli się przez $17$?}{nie}
  \zz{czy $3\,570$ dzieli się przez $7$?}{tak}
  \zz{czy $3\,915$ dzieli się przez $17$?}{nie}
  \zz{czy $4\,900$ dzieli się przez $7$?}{tak}
  \zz{czy $2\,506$ dzieli się przez $7$?}{tak}
  \zz{czy $9\,571$ dzieli się przez $17$?}{tak}
  \zz{czy $6\,947$ dzieli się przez $11$?}{nie}
  \zz{czy $1\,148$ dzieli się przez $7$?}{tak}
  \zz{czy $1\,513$ dzieli się przez $17$?}{tak}
  \zz{czy $5\,863$ dzieli się przez $11$?}{tak}
  \zz{czy $3\,213$ dzieli się przez $17$?}{tak}
  \zz{czy $391$ dzieli się przez $17$?}{tak}
  \zz{czy $1\,610$ dzieli się przez $17$?}{nie}
\end{zestaw}

\btexpect{21}% ======================================================
%  GENERATED by tools/gen_sets.py -- do not edit.
%  Seed 2021; regenerate with `make sets`.
% ======================================================

\begin{zestaw}{Dodawanie dwójkowe III}{2}{7:20}{2}
  \z{$111000_2 + 11100_2$}{1010100}
  \z{$111001_2 + 101011_2$}{1100100}
  \z{$100101_2 + 10100_2$}{111001}
  \z{$10010_2 + 111011_2$}{1001101}
  \z{$101011_2 + 101_2$}{110000}
  \z{$11111_2 + 100001_2$}{1000000}
  \z{$100111_2 + 111_2$}{101110}
  \z{$10111_2 + 10100_2$}{101011}
  \z{$10101_2 + 100001_2$}{110110}
  \z{$111_2 + 111011_2$}{1000010}
  \z{$1101_2 + 1101_2$}{11010}
  \z{$111000_2 + 110_2$}{111110}
  \z{$110010_2 + 1001_2$}{111011}
  \z{$101110_2 + 110000_2$}{1011110}
  \z{$100110_2 + 10101_2$}{111011}
  \z{$110_2 + 100101_2$}{101011}
  \z{$110111_2 + 101110_2$}{1100101}
  \z{$100000_2 + 1010_2$}{101010}
  \z{$1111_2 + 101001_2$}{111000}
  \z{$101011_2 + 11111_2$}{1001010}
  \btnc
  \z{$101011_2 + 111000_2$}{1100011}
  \z{$10100_2 + 100100_2$}{111000}
  \z{$11000_2 + 101101_2$}{1000101}
  \z{$11010_2 + 110110_2$}{1010000}
  \z{$11100_2 + 101111_2$}{1001011}
  \z{$1010_2 + 10001_2$}{11011}
  \z{$10010_2 + 100000_2$}{110010}
  \z{$1010_2 + 101100_2$}{110110}
  \z{$100101_2 + 10001_2$}{110110}
  \z{$100000_2 + 100001_2$}{1000001}
  \z{$10101_2 + 111100_2$}{1010001}
  \z{$1011_2 + 10100_2$}{11111}
  \z{$11_2 + 100_2$}{111}
  \z{$100101_2 + 110011_2$}{1011000}
  \z{$111011_2 + 10001_2$}{1001100}
  \z{$1111_2 + 110011_2$}{1000010}
  \z{$1101_2 + 111000_2$}{1000101}
  \z{$110110_2 + 10010_2$}{1001000}
  \z{$11100_2 + 11110_2$}{111010}
  \z{$101000_2 + 101100_2$}{1010100}
\end{zestaw}

\btexpect{22}% ======================================================
%  GENERATED by tools/gen_sets.py -- do not edit.
%  Seed 2022; regenerate with `make sets`.
% ======================================================

\begin{zestaw}{Potęga modulo III}{2}{8:00}{3}
  \z{$6^{5} \bmod 13$}{2}
  \z{$6^{6} \bmod 5$}{1}
  \z{$8^{4} \bmod 5$}{1}
  \z{$8^{4} \bmod 11$}{4}
  \z{$2^{4} \bmod 5$}{1}
  \z{$7^{4} \bmod 11$}{3}
  \z{$3^{5} \bmod 9$}{0}
  \z{$9^{3} \bmod 5$}{4}
  \z{$2^{3} \bmod 11$}{8}
  \z{$8^{2} \bmod 5$}{4}
  \z{$3^{6} \bmod 11$}{3}
  \z{$9^{4} \bmod 5$}{1}
  \z{$7^{5} \bmod 9$}{4}
  \z{$2^{5} \bmod 7$}{4}
  \btnc
  \z{$2^{2} \bmod 11$}{4}
  \z{$7^{6} \bmod 11$}{4}
  \z{$6^{2} \bmod 5$}{1}
  \z{$4^{2} \bmod 11$}{5}
  \z{$8^{5} \bmod 13$}{8}
  \z{$4^{4} \bmod 13$}{9}
  \z{$8^{3} \bmod 5$}{2}
  \z{$5^{3} \bmod 5$}{0}
  \z{$9^{3} \bmod 11$}{3}
  \z{$3^{3} \bmod 13$}{1}
  \z{$7^{3} \bmod 9$}{1}
  \z{$7^{4} \bmod 9$}{7}
  \z{$3^{5} \bmod 11$}{1}
  \btnc
  \z{$8^{5} \bmod 11$}{10}
  \z{$5^{6} \bmod 11$}{5}
  \z{$9^{3} \bmod 7$}{1}
  \z{$4^{5} \bmod 11$}{1}
  \z{$6^{2} \bmod 7$}{1}
  \z{$6^{6} \bmod 9$}{0}
  \z{$4^{4} \bmod 5$}{1}
  \z{$6^{3} \bmod 7$}{6}
  \z{$5^{5} \bmod 9$}{2}
  \z{$5^{2} \bmod 7$}{4}
  \z{$7^{2} \bmod 9$}{4}
  \z{$8^{2} \bmod 13$}{12}
  \z{$4^{5} \bmod 7$}{2}
\end{zestaw}

\btexpect{23}% ======================================================
%  GENERATED by tools/gen_sets.py -- do not edit.
%  Seed 2023; regenerate with `make sets`.
% ======================================================

\begin{zestaw}{System dwójkowy IV}{2}{6:40}{3}
  \z{$1100111_2$}{103}
  \z{$1110111_2$}{119}
  \z{$1010110_2$}{86}
  \z{$243$ dwójkowo}{11110011}
  \z{$11110_2$}{30}
  \z{$171$ dwójkowo}{10101011}
  \z{$10110111_2$}{183}
  \z{$140$ dwójkowo}{10001100}
  \z{$1000000_2$}{64}
  \z{$242$ dwójkowo}{11110010}
  \z{$9$ dwójkowo}{1001}
  \z{$10110_2$}{22}
  \z{$10111001_2$}{185}
  \z{$14$ dwójkowo}{1110}
  \btnc
  \z{$253$ dwójkowo}{11111101}
  \z{$133$ dwójkowo}{10000101}
  \z{$10001110_2$}{142}
  \z{$132$ dwójkowo}{10000100}
  \z{$100100_2$}{36}
  \z{$11000010_2$}{194}
  \z{$10111101_2$}{189}
  \z{$153$ dwójkowo}{10011001}
  \z{$101000_2$}{40}
  \z{$1111111_2$}{127}
  \z{$10011111_2$}{159}
  \z{$141$ dwójkowo}{10001101}
  \z{$1001111_2$}{79}
  \btnc
  \z{$11010110_2$}{214}
  \z{$100011_2$}{35}
  \z{$19$ dwójkowo}{10011}
  \z{$101100_2$}{44}
  \z{$11011_2$}{27}
  \z{$11110001_2$}{241}
  \z{$10011000_2$}{152}
  \z{$75$ dwójkowo}{1001011}
  \z{$213$ dwójkowo}{11010101}
  \z{$11011001_2$}{217}
  \z{$11000_2$}{24}
  \z{$10010100_2$}{148}
  \z{$101010_2$}{42}
\end{zestaw}

\btexpect{24}% ======================================================
%  GENERATED by tools/gen_sets.py -- do not edit.
%  Seed 2024; regenerate with `make sets`.
% ======================================================

\begin{zestaw}{Reszta modulo IV}{2}{5:20}{3}
  \z{$236 \bmod 12$}{8}
  \z{$255 \bmod 8$}{7}
  \z{$825 \bmod 11$}{0}
  \z{$595 \bmod 8$}{3}
  \z{$701 \bmod 7$}{1}
  \z{$412 \bmod 12$}{4}
  \z{$589 \bmod 11$}{6}
  \z{$366 \bmod 7$}{2}
  \z{$770 \bmod 13$}{3}
  \z{$581 \bmod 9$}{5}
  \z{$798 \bmod 4$}{2}
  \z{$755 \bmod 7$}{6}
  \z{$996 \bmod 12$}{0}
  \z{$594 \bmod 6$}{0}
  \btnc
  \z{$938 \bmod 7$}{0}
  \z{$109 \bmod 11$}{10}
  \z{$693 \bmod 9$}{0}
  \z{$527 \bmod 11$}{10}
  \z{$792 \bmod 4$}{0}
  \z{$831 \bmod 6$}{3}
  \z{$449 \bmod 9$}{8}
  \z{$403 \bmod 9$}{7}
  \z{$383 \bmod 7$}{5}
  \z{$474 \bmod 11$}{1}
  \z{$631 \bmod 9$}{1}
  \z{$976 \bmod 7$}{3}
  \z{$284 \bmod 11$}{9}
  \btnc
  \z{$91 \bmod 7$}{0}
  \z{$830 \bmod 7$}{4}
  \z{$936 \bmod 3$}{0}
  \z{$914 \bmod 8$}{2}
  \z{$376 \bmod 13$}{12}
  \z{$678 \bmod 11$}{7}
  \z{$389 \bmod 4$}{1}
  \z{$291 \bmod 7$}{4}
  \z{$424 \bmod 12$}{4}
  \z{$257 \bmod 6$}{5}
  \z{$555 \bmod 9$}{6}
  \z{$318 \bmod 6$}{0}
  \z{$688 \bmod 11$}{6}
\end{zestaw}



\chapter{Blok II — Struktura liczby}

Liczba to nie tylko jej wartość. Rozkład na czynniki pierwsze, liczba
dzielników, algorytm Euklidesa dla par, w których wspólnego dzielnika nie widać
gołym okiem, i dwa pytania, na które odpowiedź brzmi tak albo nie: czy to jest
liczba pierwsza i czy to jest kwadrat.

% GENERATED by tools/gen_sets.py -- do not edit.
\btexpect{25}% ======================================================
%  GENERATED by tools/gen_sets.py -- do not edit.
%  Seed 2025; regenerate with `make sets`.
% ======================================================

\begin{zestaw}{Rozkład na czynniki}{2}{8:00}{3}
  \z{$309$ na czynniki}{$3 \cdot 103$}
  \z{$66$ na czynniki}{$2 \cdot 3 \cdot 11$}
  \z{$354$ na czynniki}{$2 \cdot 3 \cdot 59$}
  \z{$268$ na czynniki}{$2^{2} \cdot 67$}
  \z{$112$ na czynniki}{$2^{4} \cdot 7$}
  \z{$294$ na czynniki}{$2 \cdot 3 \cdot 7^{2}$}
  \z{$24$ na czynniki}{$2^{3} \cdot 3$}
  \z{$214$ na czynniki}{$2 \cdot 107$}
  \z{$216$ na czynniki}{$2^{3} \cdot 3^{3}$}
  \z{$315$ na czynniki}{$3^{2} \cdot 5 \cdot 7$}
  \z{$142$ na czynniki}{$2 \cdot 71$}
  \z{$57$ na czynniki}{$3 \cdot 19$}
  \z{$219$ na czynniki}{$3 \cdot 73$}
  \z{$74$ na czynniki}{$2 \cdot 37$}
  \btnc
  \z{$46$ na czynniki}{$2 \cdot 23$}
  \z{$84$ na czynniki}{$2^{2} \cdot 3 \cdot 7$}
  \z{$132$ na czynniki}{$2^{2} \cdot 3 \cdot 11$}
  \z{$234$ na czynniki}{$2 \cdot 3^{2} \cdot 13$}
  \z{$38$ na czynniki}{$2 \cdot 19$}
  \z{$285$ na czynniki}{$3 \cdot 5 \cdot 19$}
  \z{$58$ na czynniki}{$2 \cdot 29$}
  \z{$123$ na czynniki}{$3 \cdot 41$}
  \z{$120$ na czynniki}{$2^{3} \cdot 3 \cdot 5$}
  \z{$386$ na czynniki}{$2 \cdot 193$}
  \z{$140$ na czynniki}{$2^{2} \cdot 5 \cdot 7$}
  \z{$351$ na czynniki}{$3^{3} \cdot 13$}
  \z{$259$ na czynniki}{$7 \cdot 37$}
  \btnc
  \z{$205$ na czynniki}{$5 \cdot 41$}
  \z{$290$ na czynniki}{$2 \cdot 5 \cdot 29$}
  \z{$39$ na czynniki}{$3 \cdot 13$}
  \z{$328$ na czynniki}{$2^{3} \cdot 41$}
  \z{$80$ na czynniki}{$2^{4} \cdot 5$}
  \z{$303$ na czynniki}{$3 \cdot 101$}
  \z{$130$ na czynniki}{$2 \cdot 5 \cdot 13$}
  \z{$86$ na czynniki}{$2 \cdot 43$}
  \z{$148$ na czynniki}{$2^{2} \cdot 37$}
  \z{$334$ na czynniki}{$2 \cdot 167$}
  \z{$40$ na czynniki}{$2^{3} \cdot 5$}
  \z{$93$ na czynniki}{$3 \cdot 31$}
  \z{$306$ na czynniki}{$2 \cdot 3^{2} \cdot 17$}
\end{zestaw}

\btexpect{26}% ======================================================
%  GENERATED by tools/gen_sets.py -- do not edit.
%  Seed 2026; regenerate with `make sets`.
% ======================================================

\begin{zestaw}{Czy pierwsza}{2}{6:40}{2}
  \z{czy $223$ jest pierwsza?}{tak}
  \z{czy $161$ jest pierwsza?}{nie}
  \z{czy $361$ jest pierwsza?}{nie}
  \z{czy $323$ jest pierwsza?}{nie}
  \z{czy $53$ jest pierwsza?}{tak}
  \z{czy $91$ jest pierwsza?}{nie}
  \z{czy $119$ jest pierwsza?}{nie}
  \z{czy $391$ jest pierwsza?}{nie}
  \z{czy $187$ jest pierwsza?}{nie}
  \z{czy $397$ jest pierwsza?}{tak}
  \z{czy $221$ jest pierwsza?}{nie}
  \z{czy $133$ jest pierwsza?}{nie}
  \z{czy $269$ jest pierwsza?}{tak}
  \z{czy $247$ jest pierwsza?}{nie}
  \z{czy $283$ jest pierwsza?}{tak}
  \z{czy $197$ jest pierwsza?}{tak}
  \z{czy $179$ jest pierwsza?}{tak}
  \z{czy $227$ jest pierwsza?}{tak}
  \z{czy $181$ jest pierwsza?}{tak}
  \z{czy $89$ jest pierwsza?}{tak}
  \btnc
  \z{czy $437$ jest pierwsza?}{nie}
  \z{czy $97$ jest pierwsza?}{tak}
  \z{czy $113$ jest pierwsza?}{tak}
  \z{czy $289$ jest pierwsza?}{nie}
  \z{czy $347$ jest pierwsza?}{tak}
  \z{czy $61$ jest pierwsza?}{tak}
  \z{czy $359$ jest pierwsza?}{tak}
  \z{czy $143$ jest pierwsza?}{nie}
  \z{czy $211$ jest pierwsza?}{tak}
  \z{czy $313$ jest pierwsza?}{tak}
  \z{czy $373$ jest pierwsza?}{tak}
  \z{czy $379$ jest pierwsza?}{tak}
  \z{czy $299$ jest pierwsza?}{nie}
  \z{czy $349$ jest pierwsza?}{tak}
  \z{czy $193$ jest pierwsza?}{tak}
  \z{czy $49$ jest pierwsza?}{nie}
  \z{czy $229$ jest pierwsza?}{tak}
  \z{czy $77$ jest pierwsza?}{nie}
  \z{czy $151$ jest pierwsza?}{tak}
  \z{czy $209$ jest pierwsza?}{nie}
\end{zestaw}

\btexpect{27}% ======================================================
%  GENERATED by tools/gen_sets.py -- do not edit.
%  Seed 2027; regenerate with `make sets`.
% ======================================================

\begin{zestaw}{Algorytm Euklidesa}{2}{8:40}{3}
  \z{NWD$(330, 522)$}{6}
  \z{NWD$(322, 329)$}{7}
  \z{NWD$(1\,869, 420)$}{21}
  \z{NWD$(1\,134, 1\,029)$}{21}
  \z{NWD$(240, 186)$}{6}
  \z{NWD$(660, 468)$}{12}
  \z{NWD$(897, 767)$}{13}
  \z{NWD$(357, 623)$}{7}
  \z{NWD$(210, 108)$}{6}
  \z{NWD$(852, 528)$}{12}
  \z{NWD$(1\,411, 374)$}{17}
  \z{NWD$(315, 602)$}{7}
  \z{NWD$(126, 534)$}{6}
  \z{NWD$(1\,275, 1\,088)$}{17}
  \btnc
  \z{NWD$(438, 240)$}{6}
  \z{NWD$(441, 774)$}{9}
  \z{NWD$(1\,575, 462)$}{21}
  \z{NWD$(238, 1\,445)$}{17}
  \z{NWD$(546, 1\,869)$}{21}
  \z{NWD$(98, 77)$}{7}
  \z{NWD$(1\,113, 1\,218)$}{21}
  \z{NWD$(799, 1\,275)$}{17}
  \z{NWD$(481, 767)$}{13}
  \z{NWD$(630, 574)$}{14}
  \z{NWD$(760, 437)$}{19}
  \z{NWD$(390, 168)$}{6}
  \z{NWD$(903, 1\,407)$}{21}
  \btnc
  \z{NWD$(301, 616)$}{7}
  \z{NWD$(585, 728)$}{13}
  \z{NWD$(270, 549)$}{9}
  \z{NWD$(931, 703)$}{19}
  \z{NWD$(616, 935)$}{11}
  \z{NWD$(801, 540)$}{9}
  \z{NWD$(385, 264)$}{11}
  \z{NWD$(234, 258)$}{6}
  \z{NWD$(806, 1\,105)$}{13}
  \z{NWD$(405, 711)$}{9}
  \z{NWD$(338, 273)$}{13}
  \z{NWD$(1\,088, 731)$}{17}
  \z{NWD$(468, 498)$}{6}
\end{zestaw}

\btexpect{28}% ======================================================
%  GENERATED by tools/gen_sets.py -- do not edit.
%  Seed 2028; regenerate with `make sets`.
% ======================================================

\begin{zestaw}{Czy kwadrat}{2}{6:00}{2}
  \z{czy $1\,156$ jest kwadratem?}{tak}
  \z{czy $6\,249$ jest kwadratem?}{nie}
  \z{czy $1\,024$ jest kwadratem?}{tak}
  \z{czy $3\,225$ jest kwadratem?}{nie}
  \z{czy $8\,100$ jest kwadratem?}{tak}
  \z{czy $144$ jest kwadratem?}{tak}
  \z{czy $8\,669$ jest kwadratem?}{nie}
  \z{czy $9\,212$ jest kwadratem?}{nie}
  \z{czy $377$ jest kwadratem?}{nie}
  \z{czy $5\,184$ jest kwadratem?}{tak}
  \z{czy $4\,977$ jest kwadratem?}{nie}
  \z{czy $3\,901$ jest kwadratem?}{nie}
  \z{czy $8\,361$ jest kwadratem?}{nie}
  \z{czy $7\,813$ jest kwadratem?}{nie}
  \z{czy $7\,866$ jest kwadratem?}{nie}
  \z{czy $6\,891$ jest kwadratem?}{nie}
  \z{czy $8\,649$ jest kwadratem?}{tak}
  \z{czy $7\,056$ jest kwadratem?}{tak}
  \z{czy $9\,246$ jest kwadratem?}{nie}
  \z{czy $1\,764$ jest kwadratem?}{tak}
  \btnc
  \z{czy $196$ jest kwadratem?}{tak}
  \z{czy $8\,073$ jest kwadratem?}{nie}
  \z{czy $8\,828$ jest kwadratem?}{nie}
  \z{czy $1\,600$ jest kwadratem?}{tak}
  \z{czy $5\,476$ jest kwadratem?}{tak}
  \z{czy $5\,041$ jest kwadratem?}{tak}
  \z{czy $3\,745$ jest kwadratem?}{nie}
  \z{czy $1\,131$ jest kwadratem?}{nie}
  \z{czy $8\,175$ jest kwadratem?}{nie}
  \z{czy $3\,569$ jest kwadratem?}{nie}
  \z{czy $4\,374$ jest kwadratem?}{nie}
  \z{czy $7\,412$ jest kwadratem?}{nie}
  \z{czy $6\,201$ jest kwadratem?}{nie}
  \z{czy $8\,859$ jest kwadratem?}{nie}
  \z{czy $4\,671$ jest kwadratem?}{nie}
  \z{czy $6\,825$ jest kwadratem?}{nie}
  \z{czy $2\,193$ jest kwadratem?}{nie}
  \z{czy $9\,764$ jest kwadratem?}{nie}
  \z{czy $6\,637$ jest kwadratem?}{nie}
  \z{czy $9\,588$ jest kwadratem?}{nie}
\end{zestaw}

\btexpect{29}% ======================================================
%  GENERATED by tools/gen_sets.py -- do not edit.
%  Seed 2029; regenerate with `make sets`.
% ======================================================

\begin{zestaw}{Sumy ciągów}{2}{6:00}{3}
  \z{$1 + 3 + \ldots + 103$}{2\,704}
  \z{$1 + 2 + \ldots + 36$}{666}
  \z{$1 + 2 + \ldots + 53$}{1\,431}
  \z{$1 + 3 + \ldots + 99$}{2\,500}
  \z{$2 + 4 + \ldots + 78$}{1\,560}
  \z{$1 + 3 + \ldots + 27$}{196}
  \z{$2 + 4 + \ldots + 52$}{702}
  \z{$1 + 3 + \ldots + 77$}{1\,521}
  \z{$1 + 3 + \ldots + 113$}{3\,249}
  \z{$1 + 3 + \ldots + 61$}{961}
  \z{$1 + 2 + \ldots + 40$}{820}
  \z{$2 + 4 + \ldots + 44$}{506}
  \z{$2 + 4 + \ldots + 14$}{56}
  \z{$1 + 3 + \ldots + 23$}{144}
  \btnc
  \z{$1 + 2 + \ldots + 17$}{153}
  \z{$2 + 4 + \ldots + 16$}{72}
  \z{$1 + 2 + \ldots + 21$}{231}
  \z{$2 + 4 + \ldots + 116$}{3\,422}
  \z{$2 + 4 + \ldots + 24$}{156}
  \z{$2 + 4 + \ldots + 96$}{2\,352}
  \z{$2 + 4 + \ldots + 48$}{600}
  \z{$1 + 2 + \ldots + 10$}{55}
  \z{$1 + 2 + \ldots + 12$}{78}
  \z{$1 + 3 + \ldots + 39$}{400}
  \z{$1 + 3 + \ldots + 65$}{1\,089}
  \z{$1 + 3 + \ldots + 57$}{841}
  \z{$1 + 3 + \ldots + 51$}{676}
  \btnc
  \z{$1 + 2 + \ldots + 18$}{171}
  \z{$1 + 2 + \ldots + 35$}{630}
  \z{$1 + 2 + \ldots + 45$}{1\,035}
  \z{$1 + 3 + \ldots + 55$}{784}
  \z{$1 + 3 + \ldots + 79$}{1\,600}
  \z{$1 + 2 + \ldots + 56$}{1\,596}
  \z{$1 + 2 + \ldots + 50$}{1\,275}
  \z{$1 + 2 + \ldots + 19$}{190}
  \z{$1 + 2 + \ldots + 25$}{325}
  \z{$1 + 3 + \ldots + 29$}{225}
  \z{$1 + 3 + \ldots + 35$}{324}
  \z{$2 + 4 + \ldots + 60$}{930}
  \z{$1 + 3 + \ldots + 109$}{3\,025}
\end{zestaw}

\btexpect{30}% ======================================================
%  GENERATED by tools/gen_sets.py -- do not edit.
%  Seed 2030; regenerate with `make sets`.
% ======================================================

\begin{zestaw}{Ile dzielników}{2}{7:20}{2}
  \z{ile dzielników ma $158$}{4}
  \z{ile dzielników ma $209$}{4}
  \z{ile dzielników ma $202$}{4}
  \z{ile dzielników ma $268$}{6}
  \z{ile dzielników ma $12$}{6}
  \z{ile dzielników ma $252$}{18}
  \z{ile dzielników ma $362$}{4}
  \z{ile dzielników ma $246$}{8}
  \z{ile dzielników ma $235$}{4}
  \z{ile dzielników ma $207$}{6}
  \z{ile dzielników ma $79$}{2}
  \z{ile dzielników ma $384$}{16}
  \z{ile dzielników ma $323$}{4}
  \z{ile dzielników ma $164$}{6}
  \z{ile dzielników ma $59$}{2}
  \z{ile dzielników ma $308$}{12}
  \z{ile dzielników ma $228$}{12}
  \z{ile dzielników ma $41$}{2}
  \z{ile dzielników ma $276$}{12}
  \z{ile dzielników ma $122$}{4}
  \btnc
  \z{ile dzielników ma $250$}{8}
  \z{ile dzielników ma $274$}{4}
  \z{ile dzielników ma $206$}{4}
  \z{ile dzielników ma $191$}{2}
  \z{ile dzielników ma $129$}{4}
  \z{ile dzielników ma $278$}{4}
  \z{ile dzielników ma $290$}{8}
  \z{ile dzielników ma $204$}{12}
  \z{ile dzielników ma $91$}{4}
  \z{ile dzielników ma $128$}{8}
  \z{ile dzielników ma $213$}{4}
  \z{ile dzielników ma $150$}{12}
  \z{ile dzielników ma $65$}{4}
  \z{ile dzielników ma $118$}{4}
  \z{ile dzielników ma $264$}{16}
  \z{ile dzielników ma $265$}{4}
  \z{ile dzielników ma $353$}{2}
  \z{ile dzielników ma $287$}{4}
  \z{ile dzielników ma $285$}{8}
  \z{ile dzielników ma $310$}{8}
\end{zestaw}

\btexpect{31}% ======================================================
%  GENERATED by tools/gen_sets.py -- do not edit.
%  Seed 2031; regenerate with `make sets`.
% ======================================================

\begin{zestaw}{Następna pierwsza}{2}{7:20}{2}
  \zz{pierwsza liczba pierwsza po $360$}{367}
  \zz{pierwsza liczba pierwsza po $111$}{113}
  \zz{pierwsza liczba pierwsza po $189$}{191}
  \zz{pierwsza liczba pierwsza po $208$}{211}
  \zz{pierwsza liczba pierwsza po $161$}{163}
  \zz{pierwsza liczba pierwsza po $255$}{257}
  \zz{pierwsza liczba pierwsza po $246$}{251}
  \zz{pierwsza liczba pierwsza po $316$}{317}
  \zz{pierwsza liczba pierwsza po $120$}{127}
  \zz{pierwsza liczba pierwsza po $133$}{137}
  \zz{pierwsza liczba pierwsza po $265$}{269}
  \zz{pierwsza liczba pierwsza po $165$}{167}
  \zz{pierwsza liczba pierwsza po $225$}{227}
  \zz{pierwsza liczba pierwsza po $54$}{59}
  \zz{pierwsza liczba pierwsza po $300$}{307}
  \zz{pierwsza liczba pierwsza po $105$}{107}
  \zz{pierwsza liczba pierwsza po $119$}{127}
  \zz{pierwsza liczba pierwsza po $364$}{367}
  \zz{pierwsza liczba pierwsza po $230$}{233}
  \zz{pierwsza liczba pierwsza po $48$}{53}
  \btnc
  \zz{pierwsza liczba pierwsza po $159$}{163}
  \zz{pierwsza liczba pierwsza po $121$}{127}
  \zz{pierwsza liczba pierwsza po $311$}{313}
  \zz{pierwsza liczba pierwsza po $209$}{211}
  \zz{pierwsza liczba pierwsza po $40$}{41}
  \zz{pierwsza liczba pierwsza po $234$}{239}
  \zz{pierwsza liczba pierwsza po $240$}{241}
  \zz{pierwsza liczba pierwsza po $138$}{139}
  \zz{pierwsza liczba pierwsza po $239$}{241}
  \zz{pierwsza liczba pierwsza po $163$}{167}
  \zz{pierwsza liczba pierwsza po $392$}{397}
  \zz{pierwsza liczba pierwsza po $81$}{83}
  \zz{pierwsza liczba pierwsza po $193$}{197}
  \zz{pierwsza liczba pierwsza po $130$}{131}
  \zz{pierwsza liczba pierwsza po $389$}{397}
  \zz{pierwsza liczba pierwsza po $217$}{223}
  \zz{pierwsza liczba pierwsza po $24$}{29}
  \zz{pierwsza liczba pierwsza po $65$}{67}
  \zz{pierwsza liczba pierwsza po $28$}{29}
  \zz{pierwsza liczba pierwsza po $88$}{89}
\end{zestaw}

\btexpect{32}% ======================================================
%  GENERATED by tools/gen_sets.py -- do not edit.
%  Seed 2032; regenerate with `make sets`.
% ======================================================

\begin{zestaw}{Rozkład na czynniki II}{2}{8:00}{3}
  \z{$321$ na czynniki}{$3 \cdot 107$}
  \z{$360$ na czynniki}{$2^{3} \cdot 3^{2} \cdot 5$}
  \z{$260$ na czynniki}{$2^{2} \cdot 5 \cdot 13$}
  \z{$85$ na czynniki}{$5 \cdot 17$}
  \z{$261$ na czynniki}{$3^{2} \cdot 29$}
  \z{$140$ na czynniki}{$2^{2} \cdot 5 \cdot 7$}
  \z{$114$ na czynniki}{$2 \cdot 3 \cdot 19$}
  \z{$363$ na czynniki}{$3 \cdot 11^{2}$}
  \z{$105$ na czynniki}{$3 \cdot 5 \cdot 7$}
  \z{$312$ na czynniki}{$2^{3} \cdot 3 \cdot 13$}
  \z{$324$ na czynniki}{$2^{2} \cdot 3^{4}$}
  \z{$358$ na czynniki}{$2 \cdot 179$}
  \z{$206$ na czynniki}{$2 \cdot 103$}
  \z{$44$ na czynniki}{$2^{2} \cdot 11$}
  \btnc
  \z{$400$ na czynniki}{$2^{4} \cdot 5^{2}$}
  \z{$393$ na czynniki}{$3 \cdot 131$}
  \z{$259$ na czynniki}{$7 \cdot 37$}
  \z{$355$ na czynniki}{$5 \cdot 71$}
  \z{$177$ na czynniki}{$3 \cdot 59$}
  \z{$147$ na czynniki}{$3 \cdot 7^{2}$}
  \z{$162$ na czynniki}{$2 \cdot 3^{4}$}
  \z{$369$ na czynniki}{$3^{2} \cdot 41$}
  \z{$323$ na czynniki}{$17 \cdot 19$}
  \z{$381$ na czynniki}{$3 \cdot 127$}
  \z{$335$ na czynniki}{$5 \cdot 67$}
  \z{$95$ na czynniki}{$5 \cdot 19$}
  \z{$302$ na czynniki}{$2 \cdot 151$}
  \btnc
  \z{$86$ na czynniki}{$2 \cdot 43$}
  \z{$205$ na czynniki}{$5 \cdot 41$}
  \z{$93$ na czynniki}{$3 \cdot 31$}
  \z{$195$ na czynniki}{$3 \cdot 5 \cdot 13$}
  \z{$326$ na czynniki}{$2 \cdot 163$}
  \z{$226$ na czynniki}{$2 \cdot 113$}
  \z{$90$ na czynniki}{$2 \cdot 3^{2} \cdot 5$}
  \z{$375$ na czynniki}{$3 \cdot 5^{3}$}
  \z{$40$ na czynniki}{$2^{3} \cdot 5$}
  \z{$362$ na czynniki}{$2 \cdot 181$}
  \z{$98$ na czynniki}{$2 \cdot 7^{2}$}
  \z{$351$ na czynniki}{$3^{3} \cdot 13$}
  \z{$196$ na czynniki}{$2^{2} \cdot 7^{2}$}
\end{zestaw}

\btexpect{33}% ======================================================
%  GENERATED by tools/gen_sets.py -- do not edit.
%  Seed 2033; regenerate with `make sets`.
% ======================================================

\begin{zestaw}{Czy pierwsza II}{2}{6:40}{2}
  \z{czy $383$ jest pierwsza?}{tak}
  \z{czy $239$ jest pierwsza?}{tak}
  \z{czy $187$ jest pierwsza?}{nie}
  \z{czy $391$ jest pierwsza?}{nie}
  \z{czy $233$ jest pierwsza?}{tak}
  \z{czy $209$ jest pierwsza?}{nie}
  \z{czy $127$ jest pierwsza?}{tak}
  \z{czy $91$ jest pierwsza?}{nie}
  \z{czy $161$ jest pierwsza?}{nie}
  \z{czy $307$ jest pierwsza?}{tak}
  \z{czy $173$ jest pierwsza?}{tak}
  \z{czy $53$ jest pierwsza?}{tak}
  \z{czy $389$ jest pierwsza?}{tak}
  \z{czy $311$ jest pierwsza?}{tak}
  \z{czy $143$ jest pierwsza?}{nie}
  \z{czy $179$ jest pierwsza?}{tak}
  \z{czy $367$ jest pierwsza?}{tak}
  \z{czy $221$ jest pierwsza?}{nie}
  \z{czy $253$ jest pierwsza?}{nie}
  \z{czy $119$ jest pierwsza?}{nie}
  \btnc
  \z{czy $101$ jest pierwsza?}{tak}
  \z{czy $299$ jest pierwsza?}{nie}
  \z{czy $199$ jest pierwsza?}{tak}
  \z{czy $133$ jest pierwsza?}{nie}
  \z{czy $167$ jest pierwsza?}{tak}
  \z{czy $257$ jest pierwsza?}{tak}
  \z{czy $191$ jest pierwsza?}{tak}
  \z{czy $131$ jest pierwsza?}{tak}
  \z{czy $181$ jest pierwsza?}{tak}
  \z{czy $89$ jest pierwsza?}{tak}
  \z{czy $151$ jest pierwsza?}{tak}
  \z{czy $373$ jest pierwsza?}{tak}
  \z{czy $323$ jest pierwsza?}{nie}
  \z{czy $361$ jest pierwsza?}{nie}
  \z{czy $49$ jest pierwsza?}{nie}
  \z{czy $73$ jest pierwsza?}{tak}
  \z{czy $397$ jest pierwsza?}{tak}
  \z{czy $277$ jest pierwsza?}{tak}
  \z{czy $337$ jest pierwsza?}{tak}
  \z{czy $289$ jest pierwsza?}{nie}
\end{zestaw}

\btexpect{34}% ======================================================
%  GENERATED by tools/gen_sets.py -- do not edit.
%  Seed 2034; regenerate with `make sets`.
% ======================================================

\begin{zestaw}{Algorytm Euklidesa II}{2}{8:40}{3}
  \z{NWD$(102, 492)$}{6}
  \z{NWD$(210, 534)$}{6}
  \z{NWD$(165, 108)$}{3}
  \z{NWD$(143, 396)$}{11}
  \z{NWD$(216, 237)$}{3}
  \z{NWD$(708, 132)$}{12}
  \z{NWD$(196, 630)$}{14}
  \z{NWD$(469, 336)$}{7}
  \z{NWD$(196, 119)$}{7}
  \z{NWD$(114, 141)$}{3}
  \z{NWD$(437, 494)$}{19}
  \z{NWD$(150, 534)$}{6}
  \z{NWD$(444, 600)$}{12}
  \z{NWD$(495, 621)$}{9}
  \btnc
  \z{NWD$(658, 434)$}{14}
  \z{NWD$(702, 949)$}{13}
  \z{NWD$(1\,027, 1\,040)$}{13}
  \z{NWD$(1\,530, 391)$}{17}
  \z{NWD$(180, 39)$}{3}
  \z{NWD$(408, 697)$}{17}
  \z{NWD$(132, 336)$}{12}
  \z{NWD$(728, 247)$}{13}
  \z{NWD$(522, 282)$}{6}
  \z{NWD$(1\,539, 380)$}{19}
  \z{NWD$(1\,134, 350)$}{14}
  \z{NWD$(1\,869, 735)$}{21}
  \z{NWD$(828, 744)$}{12}
  \btnc
  \z{NWD$(1\,027, 221)$}{13}
  \z{NWD$(637, 1\,079)$}{13}
  \z{NWD$(1\,311, 1\,387)$}{19}
  \z{NWD$(949, 598)$}{13}
  \z{NWD$(595, 126)$}{7}
  \z{NWD$(132, 624)$}{12}
  \z{NWD$(627, 407)$}{11}
  \z{NWD$(406, 119)$}{7}
  \z{NWD$(243, 183)$}{3}
  \z{NWD$(329, 140)$}{7}
  \z{NWD$(767, 676)$}{13}
  \z{NWD$(189, 33)$}{3}
  \z{NWD$(517, 660)$}{11}
\end{zestaw}

\btexpect{35}% ======================================================
%  GENERATED by tools/gen_sets.py -- do not edit.
%  Seed 2035; regenerate with `make sets`.
% ======================================================

\begin{zestaw}{Czy kwadrat II}{2}{6:00}{2}
  \z{czy $2\,934$ jest kwadratem?}{nie}
  \z{czy $702$ jest kwadratem?}{nie}
  \z{czy $2\,601$ jest kwadratem?}{tak}
  \z{czy $1\,814$ jest kwadratem?}{nie}
  \z{czy $529$ jest kwadratem?}{tak}
  \z{czy $5\,392$ jest kwadratem?}{nie}
  \z{czy $1\,708$ jest kwadratem?}{nie}
  \z{czy $6\,388$ jest kwadratem?}{nie}
  \z{czy $4\,470$ jest kwadratem?}{nie}
  \z{czy $2\,304$ jest kwadratem?}{tak}
  \z{czy $8\,569$ jest kwadratem?}{nie}
  \z{czy $625$ jest kwadratem?}{tak}
  \z{czy $4\,356$ jest kwadratem?}{tak}
  \z{czy $5\,306$ jest kwadratem?}{nie}
  \z{czy $1\,936$ jest kwadratem?}{tak}
  \z{czy $3\,481$ jest kwadratem?}{tak}
  \z{czy $641$ jest kwadratem?}{nie}
  \z{czy $2\,256$ jest kwadratem?}{nie}
  \z{czy $5\,476$ jest kwadratem?}{tak}
  \z{czy $608$ jest kwadratem?}{nie}
  \btnc
  \z{czy $3\,685$ jest kwadratem?}{nie}
  \z{czy $289$ jest kwadratem?}{tak}
  \z{czy $2\,401$ jest kwadratem?}{tak}
  \z{czy $2\,637$ jest kwadratem?}{nie}
  \z{czy $4\,898$ jest kwadratem?}{nie}
  \z{czy $784$ jest kwadratem?}{tak}
  \z{czy $6\,602$ jest kwadratem?}{nie}
  \z{czy $441$ jest kwadratem?}{tak}
  \z{czy $389$ jest kwadratem?}{nie}
  \z{czy $4\,552$ jest kwadratem?}{nie}
  \z{czy $7\,921$ jest kwadratem?}{tak}
  \z{czy $196$ jest kwadratem?}{tak}
  \z{czy $7\,285$ jest kwadratem?}{nie}
  \z{czy $8\,464$ jest kwadratem?}{tak}
  \z{czy $4\,225$ jest kwadratem?}{tak}
  \z{czy $1\,225$ jest kwadratem?}{tak}
  \z{czy $9\,409$ jest kwadratem?}{tak}
  \z{czy $4\,926$ jest kwadratem?}{nie}
  \z{czy $5\,776$ jest kwadratem?}{tak}
  \z{czy $900$ jest kwadratem?}{tak}
\end{zestaw}

\btexpect{36}% ======================================================
%  GENERATED by tools/gen_sets.py -- do not edit.
%  Seed 2036; regenerate with `make sets`.
% ======================================================

\begin{zestaw}{Sumy ciągów II}{2}{6:00}{3}
  \z{$2 + 4 + \ldots + 18$}{90}
  \z{$2 + 4 + \ldots + 84$}{1\,806}
  \z{$1 + 3 + \ldots + 119$}{3\,600}
  \z{$1 + 2 + \ldots + 6$}{21}
  \z{$1 + 2 + \ldots + 41$}{861}
  \z{$1 + 3 + \ldots + 61$}{961}
  \z{$2 + 4 + \ldots + 20$}{110}
  \z{$2 + 4 + \ldots + 14$}{56}
  \z{$2 + 4 + \ldots + 66$}{1\,122}
  \z{$1 + 3 + \ldots + 19$}{100}
  \z{$1 + 2 + \ldots + 57$}{1\,653}
  \z{$1 + 2 + \ldots + 32$}{528}
  \z{$1 + 2 + \ldots + 58$}{1\,711}
  \z{$1 + 2 + \ldots + 60$}{1\,830}
  \btnc
  \z{$2 + 4 + \ldots + 32$}{272}
  \z{$1 + 2 + \ldots + 16$}{136}
  \z{$1 + 3 + \ldots + 65$}{1\,089}
  \z{$1 + 2 + \ldots + 10$}{55}
  \z{$2 + 4 + \ldots + 116$}{3\,422}
  \z{$2 + 4 + \ldots + 60$}{930}
  \z{$1 + 3 + \ldots + 31$}{256}
  \z{$1 + 2 + \ldots + 12$}{78}
  \z{$1 + 2 + \ldots + 31$}{496}
  \z{$2 + 4 + \ldots + 12$}{42}
  \z{$2 + 4 + \ldots + 56$}{812}
  \z{$1 + 3 + \ldots + 71$}{1\,296}
  \z{$2 + 4 + \ldots + 48$}{600}
  \btnc
  \z{$1 + 2 + \ldots + 55$}{1\,540}
  \z{$1 + 3 + \ldots + 113$}{3\,249}
  \z{$2 + 4 + \ldots + 100$}{2\,550}
  \z{$1 + 3 + \ldots + 63$}{1\,024}
  \z{$1 + 3 + \ldots + 91$}{2\,116}
  \z{$1 + 3 + \ldots + 89$}{2\,025}
  \z{$1 + 2 + \ldots + 24$}{300}
  \z{$1 + 2 + \ldots + 51$}{1\,326}
  \z{$1 + 3 + \ldots + 23$}{144}
  \z{$1 + 2 + \ldots + 26$}{351}
  \z{$1 + 3 + \ldots + 107$}{2\,916}
  \z{$2 + 4 + \ldots + 112$}{3\,192}
  \z{$1 + 2 + \ldots + 46$}{1\,081}
\end{zestaw}

\btexpect{37}% ======================================================
%  GENERATED by tools/gen_sets.py -- do not edit.
%  Seed 2037; regenerate with `make sets`.
% ======================================================

\begin{zestaw}{Ile dzielników II}{2}{7:20}{2}
  \z{ile dzielników ma $142$}{4}
  \z{ile dzielników ma $296$}{8}
  \z{ile dzielników ma $398$}{4}
  \z{ile dzielników ma $40$}{8}
  \z{ile dzielników ma $170$}{8}
  \z{ile dzielników ma $245$}{6}
  \z{ile dzielników ma $79$}{2}
  \z{ile dzielników ma $34$}{4}
  \z{ile dzielników ma $205$}{4}
  \z{ile dzielników ma $61$}{2}
  \z{ile dzielników ma $140$}{12}
  \z{ile dzielników ma $148$}{6}
  \z{ile dzielników ma $230$}{8}
  \z{ile dzielników ma $147$}{6}
  \z{ile dzielników ma $100$}{9}
  \z{ile dzielników ma $165$}{8}
  \z{ile dzielników ma $278$}{4}
  \z{ile dzielników ma $222$}{8}
  \z{ile dzielników ma $246$}{8}
  \z{ile dzielników ma $306$}{12}
  \btnc
  \z{ile dzielników ma $379$}{2}
  \z{ile dzielników ma $393$}{4}
  \z{ile dzielników ma $365$}{4}
  \z{ile dzielników ma $292$}{6}
  \z{ile dzielników ma $45$}{6}
  \z{ile dzielników ma $382$}{4}
  \z{ile dzielników ma $336$}{20}
  \z{ile dzielników ma $123$}{4}
  \z{ile dzielników ma $260$}{12}
  \z{ile dzielników ma $185$}{4}
  \z{ile dzielników ma $208$}{10}
  \z{ile dzielników ma $214$}{4}
  \z{ile dzielników ma $350$}{12}
  \z{ile dzielników ma $104$}{8}
  \z{ile dzielników ma $135$}{8}
  \z{ile dzielników ma $273$}{8}
  \z{ile dzielników ma $49$}{3}
  \z{ile dzielników ma $284$}{6}
  \z{ile dzielników ma $237$}{4}
  \z{ile dzielników ma $360$}{24}
\end{zestaw}

\btexpect{38}% ======================================================
%  GENERATED by tools/gen_sets.py -- do not edit.
%  Seed 2038; regenerate with `make sets`.
% ======================================================

\begin{zestaw}{Następna pierwsza II}{2}{7:20}{2}
  \zz{pierwsza liczba pierwsza po $65$}{67}
  \zz{pierwsza liczba pierwsza po $263$}{269}
  \zz{pierwsza liczba pierwsza po $56$}{59}
  \zz{pierwsza liczba pierwsza po $75$}{79}
  \zz{pierwsza liczba pierwsza po $292$}{293}
  \zz{pierwsza liczba pierwsza po $304$}{307}
  \zz{pierwsza liczba pierwsza po $328$}{331}
  \zz{pierwsza liczba pierwsza po $27$}{29}
  \zz{pierwsza liczba pierwsza po $133$}{137}
  \zz{pierwsza liczba pierwsza po $105$}{107}
  \zz{pierwsza liczba pierwsza po $196$}{197}
  \zz{pierwsza liczba pierwsza po $49$}{53}
  \zz{pierwsza liczba pierwsza po $175$}{179}
  \zz{pierwsza liczba pierwsza po $70$}{71}
  \zz{pierwsza liczba pierwsza po $184$}{191}
  \zz{pierwsza liczba pierwsza po $242$}{251}
  \zz{pierwsza liczba pierwsza po $280$}{281}
  \zz{pierwsza liczba pierwsza po $24$}{29}
  \zz{pierwsza liczba pierwsza po $35$}{37}
  \zz{pierwsza liczba pierwsza po $275$}{277}
  \btnc
  \zz{pierwsza liczba pierwsza po $274$}{277}
  \zz{pierwsza liczba pierwsza po $255$}{257}
  \zz{pierwsza liczba pierwsza po $189$}{191}
  \zz{pierwsza liczba pierwsza po $301$}{307}
  \zz{pierwsza liczba pierwsza po $210$}{211}
  \zz{pierwsza liczba pierwsza po $368$}{373}
  \zz{pierwsza liczba pierwsza po $62$}{67}
  \zz{pierwsza liczba pierwsza po $199$}{211}
  \zz{pierwsza liczba pierwsza po $381$}{383}
  \zz{pierwsza liczba pierwsza po $309$}{311}
  \zz{pierwsza liczba pierwsza po $88$}{89}
  \zz{pierwsza liczba pierwsza po $373$}{379}
  \zz{pierwsza liczba pierwsza po $370$}{373}
  \zz{pierwsza liczba pierwsza po $169$}{173}
  \zz{pierwsza liczba pierwsza po $267$}{269}
  \zz{pierwsza liczba pierwsza po $300$}{307}
  \zz{pierwsza liczba pierwsza po $81$}{83}
  \zz{pierwsza liczba pierwsza po $207$}{211}
  \zz{pierwsza liczba pierwsza po $38$}{41}
  \zz{pierwsza liczba pierwsza po $71$}{73}
\end{zestaw}

\btexpect{39}% ======================================================
%  GENERATED by tools/gen_sets.py -- do not edit.
%  Seed 2039; regenerate with `make sets`.
% ======================================================

\begin{zestaw}{Rozkład na czynniki III}{2}{8:00}{3}
  \z{$365$ na czynniki}{$5 \cdot 73$}
  \z{$298$ na czynniki}{$2 \cdot 149$}
  \z{$299$ na czynniki}{$13 \cdot 23$}
  \z{$226$ na czynniki}{$2 \cdot 113$}
  \z{$30$ na czynniki}{$2 \cdot 3 \cdot 5$}
  \z{$391$ na czynniki}{$17 \cdot 23$}
  \z{$155$ na czynniki}{$5 \cdot 31$}
  \z{$171$ na czynniki}{$3^{2} \cdot 19$}
  \z{$330$ na czynniki}{$2 \cdot 3 \cdot 5 \cdot 11$}
  \z{$187$ na czynniki}{$11 \cdot 17$}
  \z{$352$ na czynniki}{$2^{5} \cdot 11$}
  \z{$322$ na czynniki}{$2 \cdot 7 \cdot 23$}
  \z{$36$ na czynniki}{$2^{2} \cdot 3^{2}$}
  \z{$124$ na czynniki}{$2^{2} \cdot 31$}
  \btnc
  \z{$182$ na czynniki}{$2 \cdot 7 \cdot 13$}
  \z{$320$ na czynniki}{$2^{6} \cdot 5$}
  \z{$91$ na czynniki}{$7 \cdot 13$}
  \z{$176$ na czynniki}{$2^{4} \cdot 11$}
  \z{$35$ na czynniki}{$5 \cdot 7$}
  \z{$300$ na czynniki}{$2^{2} \cdot 3 \cdot 5^{2}$}
  \z{$377$ na czynniki}{$13 \cdot 29$}
  \z{$316$ na czynniki}{$2^{2} \cdot 79$}
  \z{$240$ na czynniki}{$2^{4} \cdot 3 \cdot 5$}
  \z{$40$ na czynniki}{$2^{3} \cdot 5$}
  \z{$132$ na czynniki}{$2^{2} \cdot 3 \cdot 11$}
  \z{$159$ na czynniki}{$3 \cdot 53$}
  \z{$228$ na czynniki}{$2^{2} \cdot 3 \cdot 19$}
  \btnc
  \z{$160$ na czynniki}{$2^{5} \cdot 5$}
  \z{$395$ na czynniki}{$5 \cdot 79$}
  \z{$213$ na czynniki}{$3 \cdot 71$}
  \z{$270$ na czynniki}{$2 \cdot 3^{3} \cdot 5$}
  \z{$51$ na czynniki}{$3 \cdot 17$}
  \z{$342$ na czynniki}{$2 \cdot 3^{2} \cdot 19$}
  \z{$328$ na czynniki}{$2^{3} \cdot 41$}
  \z{$326$ na czynniki}{$2 \cdot 163$}
  \z{$70$ na czynniki}{$2 \cdot 5 \cdot 7$}
  \z{$148$ na czynniki}{$2^{2} \cdot 37$}
  \z{$371$ na czynniki}{$7 \cdot 53$}
  \z{$195$ na czynniki}{$3 \cdot 5 \cdot 13$}
  \z{$105$ na czynniki}{$3 \cdot 5 \cdot 7$}
\end{zestaw}

\btexpect{40}% ======================================================
%  GENERATED by tools/gen_sets.py -- do not edit.
%  Seed 2040; regenerate with `make sets`.
% ======================================================

\begin{zestaw}{Czy pierwsza III}{2}{6:40}{2}
  \z{czy $233$ jest pierwsza?}{tak}
  \z{czy $331$ jest pierwsza?}{tak}
  \z{czy $397$ jest pierwsza?}{tak}
  \z{czy $359$ jest pierwsza?}{tak}
  \z{czy $121$ jest pierwsza?}{nie}
  \z{czy $283$ jest pierwsza?}{tak}
  \z{czy $119$ jest pierwsza?}{nie}
  \z{czy $169$ jest pierwsza?}{nie}
  \z{czy $293$ jest pierwsza?}{tak}
  \z{czy $143$ jest pierwsza?}{nie}
  \z{czy $49$ jest pierwsza?}{nie}
  \z{czy $299$ jest pierwsza?}{nie}
  \z{czy $59$ jest pierwsza?}{tak}
  \z{czy $317$ jest pierwsza?}{tak}
  \z{czy $133$ jest pierwsza?}{nie}
  \z{czy $101$ jest pierwsza?}{tak}
  \z{czy $247$ jest pierwsza?}{nie}
  \z{czy $389$ jest pierwsza?}{tak}
  \z{czy $131$ jest pierwsza?}{tak}
  \z{czy $379$ jest pierwsza?}{tak}
  \btnc
  \z{czy $391$ jest pierwsza?}{nie}
  \z{czy $323$ jest pierwsza?}{nie}
  \z{czy $79$ jest pierwsza?}{tak}
  \z{czy $253$ jest pierwsza?}{nie}
  \z{czy $187$ jest pierwsza?}{nie}
  \z{czy $71$ jest pierwsza?}{tak}
  \z{czy $91$ jest pierwsza?}{nie}
  \z{czy $289$ jest pierwsza?}{nie}
  \z{czy $211$ jest pierwsza?}{tak}
  \z{czy $193$ jest pierwsza?}{tak}
  \z{czy $349$ jest pierwsza?}{tak}
  \z{czy $257$ jest pierwsza?}{tak}
  \z{czy $383$ jest pierwsza?}{tak}
  \z{czy $269$ jest pierwsza?}{tak}
  \z{czy $251$ jest pierwsza?}{tak}
  \z{czy $77$ jest pierwsza?}{nie}
  \z{czy $89$ jest pierwsza?}{tak}
  \z{czy $209$ jest pierwsza?}{nie}
  \z{czy $241$ jest pierwsza?}{tak}
  \z{czy $191$ jest pierwsza?}{tak}
\end{zestaw}

\btexpect{41}% ======================================================
%  GENERATED by tools/gen_sets.py -- do not edit.
%  Seed 2041; regenerate with `make sets`.
% ======================================================

\begin{zestaw}{Algorytm Euklidesa III}{2}{8:40}{3}
  \z{NWD$(150, 81)$}{3}
  \z{NWD$(602, 868)$}{14}
  \z{NWD$(588, 312)$}{12}
  \z{NWD$(1\,425, 1\,273)$}{19}
  \z{NWD$(913, 539)$}{11}
  \z{NWD$(702, 315)$}{9}
  \z{NWD$(711, 684)$}{9}
  \z{NWD$(377, 1\,170)$}{13}
  \z{NWD$(324, 387)$}{9}
  \z{NWD$(912, 468)$}{12}
  \z{NWD$(242, 187)$}{11}
  \z{NWD$(528, 900)$}{12}
  \z{NWD$(689, 936)$}{13}
  \z{NWD$(869, 308)$}{11}
  \btnc
  \z{NWD$(1\,275, 1\,241)$}{17}
  \z{NWD$(650, 1\,157)$}{13}
  \z{NWD$(1\,045, 266)$}{19}
  \z{NWD$(399, 1\,134)$}{21}
  \z{NWD$(918, 1\,343)$}{17}
  \z{NWD$(1\,235, 684)$}{19}
  \z{NWD$(1\,411, 1\,377)$}{17}
  \z{NWD$(1\,037, 629)$}{17}
  \z{NWD$(189, 490)$}{7}
  \z{NWD$(980, 854)$}{14}
  \z{NWD$(639, 702)$}{9}
  \z{NWD$(638, 869)$}{11}
  \z{NWD$(817, 1\,539)$}{19}
  \btnc
  \z{NWD$(946, 583)$}{11}
  \z{NWD$(425, 1\,207)$}{17}
  \z{NWD$(780, 559)$}{13}
  \z{NWD$(141, 96)$}{3}
  \z{NWD$(407, 770)$}{11}
  \z{NWD$(390, 247)$}{13}
  \z{NWD$(252, 217)$}{7}
  \z{NWD$(483, 469)$}{7}
  \z{NWD$(1\,428, 315)$}{21}
  \z{NWD$(494, 377)$}{13}
  \z{NWD$(559, 273)$}{13}
  \z{NWD$(671, 891)$}{11}
  \z{NWD$(1\,520, 209)$}{19}
\end{zestaw}

\btexpect{42}% ======================================================
%  GENERATED by tools/gen_sets.py -- do not edit.
%  Seed 2042; regenerate with `make sets`.
% ======================================================

\begin{zestaw}{Czy kwadrat III}{2}{6:00}{2}
  \z{czy $3\,980$ jest kwadratem?}{nie}
  \z{czy $5\,189$ jest kwadratem?}{nie}
  \z{czy $4\,225$ jest kwadratem?}{tak}
  \z{czy $484$ jest kwadratem?}{tak}
  \z{czy $5\,476$ jest kwadratem?}{tak}
  \z{czy $8\,830$ jest kwadratem?}{nie}
  \z{czy $5\,830$ jest kwadratem?}{nie}
  \z{czy $4\,582$ jest kwadratem?}{nie}
  \z{czy $2\,601$ jest kwadratem?}{tak}
  \z{czy $5\,625$ jest kwadratem?}{tak}
  \z{czy $6\,638$ jest kwadratem?}{nie}
  \z{czy $1\,911$ jest kwadratem?}{nie}
  \z{czy $9\,025$ jest kwadratem?}{tak}
  \z{czy $6\,561$ jest kwadratem?}{tak}
  \z{czy $4\,761$ jest kwadratem?}{tak}
  \z{czy $3\,249$ jest kwadratem?}{tak}
  \z{czy $7\,056$ jest kwadratem?}{tak}
  \z{czy $1\,849$ jest kwadratem?}{tak}
  \z{czy $3\,844$ jest kwadratem?}{tak}
  \z{czy $9\,381$ jest kwadratem?}{nie}
  \btnc
  \z{czy $6\,621$ jest kwadratem?}{nie}
  \z{czy $676$ jest kwadratem?}{tak}
  \z{czy $9\,510$ jest kwadratem?}{nie}
  \z{czy $2\,401$ jest kwadratem?}{tak}
  \z{czy $2\,966$ jest kwadratem?}{nie}
  \z{czy $6\,539$ jest kwadratem?}{nie}
  \z{czy $4\,624$ jest kwadratem?}{tak}
  \z{czy $1\,245$ jest kwadratem?}{nie}
  \z{czy $8\,100$ jest kwadratem?}{tak}
  \z{czy $8\,416$ jest kwadratem?}{nie}
  \z{czy $5\,381$ jest kwadratem?}{nie}
  \z{czy $5\,286$ jest kwadratem?}{nie}
  \z{czy $5\,329$ jest kwadratem?}{tak}
  \z{czy $7\,744$ jest kwadratem?}{tak}
  \z{czy $5\,942$ jest kwadratem?}{nie}
  \z{czy $9\,631$ jest kwadratem?}{nie}
  \z{czy $9\,582$ jest kwadratem?}{nie}
  \z{czy $8\,836$ jest kwadratem?}{tak}
  \z{czy $3\,865$ jest kwadratem?}{nie}
  \z{czy $2\,916$ jest kwadratem?}{tak}
\end{zestaw}

\btexpect{43}% ======================================================
%  GENERATED by tools/gen_sets.py -- do not edit.
%  Seed 2043; regenerate with `make sets`.
% ======================================================

\begin{zestaw}{Mieszane — struktura}{2}{7:20}{3}
  \z{$340$ na czynniki}{$2^{2} \cdot 5 \cdot 17$}
  \z{$148$ na czynniki}{$2^{2} \cdot 37$}
  \z{$365$ na czynniki}{$5 \cdot 73$}
  \z{$363$ na czynniki}{$3 \cdot 11^{2}$}
  \z{$318$ na czynniki}{$2 \cdot 3 \cdot 53$}
  \z{$309$ na czynniki}{$3 \cdot 103$}
  \z{$385$ na czynniki}{$5 \cdot 7 \cdot 11$}
  \z{$352$ na czynniki}{$2^{5} \cdot 11$}
  \z{$91$ na czynniki}{$7 \cdot 13$}
  \z{$348$ na czynniki}{$2^{2} \cdot 3 \cdot 29$}
  \z{ile dzielników ma $247$}{4}
  \z{ile dzielników ma $147$}{6}
  \z{ile dzielników ma $272$}{10}
  \z{ile dzielników ma $85$}{4}
  \btnc
  \z{ile dzielników ma $369$}{6}
  \z{ile dzielników ma $251$}{2}
  \z{ile dzielników ma $194$}{4}
  \z{ile dzielników ma $234$}{12}
  \z{ile dzielników ma $340$}{12}
  \z{ile dzielników ma $155$}{4}
  \z{NWD$(391, 1\,156)$}{17}
  \z{NWD$(737, 176)$}{11}
  \z{NWD$(312, 474)$}{6}
  \z{NWD$(603, 450)$}{9}
  \z{NWD$(1\,462, 357)$}{17}
  \z{NWD$(671, 605)$}{11}
  \z{NWD$(528, 318)$}{6}
  \btnc
  \z{NWD$(684, 528)$}{12}
  \z{NWD$(329, 126)$}{7}
  \z{NWD$(180, 534)$}{6}
  \z{$1 + 3 + \ldots + 17$}{81}
  \z{$1 + 3 + \ldots + 81$}{1\,681}
  \z{$2 + 4 + \ldots + 66$}{1\,122}
  \z{$2 + 4 + \ldots + 82$}{1\,722}
  \z{$1 + 3 + \ldots + 19$}{100}
  \z{$1 + 2 + \ldots + 22$}{253}
  \z{$1 + 3 + \ldots + 51$}{676}
  \z{$1 + 3 + \ldots + 107$}{2\,916}
  \z{$2 + 4 + \ldots + 70$}{1\,260}
  \z{$1 + 3 + \ldots + 35$}{324}
\end{zestaw}

\btexpect{44}% ======================================================
%  GENERATED by tools/gen_sets.py -- do not edit.
%  Seed 2044; regenerate with `make sets`.
% ======================================================

\begin{zestaw}{Sumy ciągów III}{2}{6:00}{3}
  \z{$2 + 4 + \ldots + 26$}{182}
  \z{$1 + 2 + \ldots + 26$}{351}
  \z{$2 + 4 + \ldots + 78$}{1\,560}
  \z{$1 + 3 + \ldots + 43$}{484}
  \z{$1 + 2 + \ldots + 39$}{780}
  \z{$1 + 3 + \ldots + 25$}{169}
  \z{$2 + 4 + \ldots + 20$}{110}
  \z{$1 + 3 + \ldots + 69$}{1\,225}
  \z{$1 + 2 + \ldots + 49$}{1\,225}
  \z{$2 + 4 + \ldots + 80$}{1\,640}
  \z{$2 + 4 + \ldots + 84$}{1\,806}
  \z{$1 + 3 + \ldots + 35$}{324}
  \z{$2 + 4 + \ldots + 50$}{650}
  \z{$1 + 3 + \ldots + 93$}{2\,209}
  \btnc
  \z{$1 + 3 + \ldots + 89$}{2\,025}
  \z{$1 + 2 + \ldots + 22$}{253}
  \z{$2 + 4 + \ldots + 60$}{930}
  \z{$1 + 3 + \ldots + 11$}{36}
  \z{$1 + 2 + \ldots + 29$}{435}
  \z{$2 + 4 + \ldots + 116$}{3\,422}
  \z{$1 + 3 + \ldots + 87$}{1\,936}
  \z{$2 + 4 + \ldots + 18$}{90}
  \z{$1 + 3 + \ldots + 109$}{3\,025}
  \z{$2 + 4 + \ldots + 58$}{870}
  \z{$1 + 3 + \ldots + 115$}{3\,364}
  \z{$1 + 2 + \ldots + 27$}{378}
  \z{$1 + 3 + \ldots + 21$}{121}
  \btnc
  \z{$2 + 4 + \ldots + 24$}{156}
  \z{$1 + 3 + \ldots + 19$}{100}
  \z{$2 + 4 + \ldots + 70$}{1\,260}
  \z{$1 + 2 + \ldots + 34$}{595}
  \z{$1 + 2 + \ldots + 17$}{153}
  \z{$1 + 2 + \ldots + 57$}{1\,653}
  \z{$1 + 2 + \ldots + 18$}{171}
  \z{$1 + 2 + \ldots + 56$}{1\,596}
  \z{$1 + 2 + \ldots + 38$}{741}
  \z{$1 + 3 + \ldots + 29$}{225}
  \z{$2 + 4 + \ldots + 102$}{2\,652}
  \z{$2 + 4 + \ldots + 54$}{756}
  \z{$1 + 2 + \ldots + 9$}{45}
\end{zestaw}

\btexpect{45}% ======================================================
%  GENERATED by tools/gen_sets.py -- do not edit.
%  Seed 2045; regenerate with `make sets`.
% ======================================================

\begin{zestaw}{Ile dzielników III}{2}{7:20}{2}
  \z{ile dzielników ma $77$}{4}
  \z{ile dzielników ma $226$}{4}
  \z{ile dzielników ma $37$}{2}
  \z{ile dzielników ma $301$}{4}
  \z{ile dzielników ma $342$}{12}
  \z{ile dzielników ma $132$}{12}
  \z{ile dzielników ma $61$}{2}
  \z{ile dzielników ma $166$}{4}
  \z{ile dzielników ma $215$}{4}
  \z{ile dzielników ma $79$}{2}
  \z{ile dzielników ma $364$}{12}
  \z{ile dzielników ma $129$}{4}
  \z{ile dzielników ma $352$}{12}
  \z{ile dzielników ma $370$}{8}
  \z{ile dzielników ma $263$}{2}
  \z{ile dzielników ma $346$}{4}
  \z{ile dzielników ma $115$}{4}
  \z{ile dzielników ma $255$}{8}
  \z{ile dzielników ma $195$}{8}
  \z{ile dzielników ma $291$}{4}
  \btnc
  \z{ile dzielników ma $45$}{6}
  \z{ile dzielników ma $276$}{12}
  \z{ile dzielników ma $323$}{4}
  \z{ile dzielników ma $279$}{6}
  \z{ile dzielników ma $375$}{8}
  \z{ile dzielników ma $339$}{4}
  \z{ile dzielników ma $173$}{2}
  \z{ile dzielników ma $217$}{4}
  \z{ile dzielników ma $344$}{8}
  \z{ile dzielników ma $133$}{4}
  \z{ile dzielników ma $219$}{4}
  \z{ile dzielników ma $198$}{12}
  \z{ile dzielników ma $73$}{2}
  \z{ile dzielników ma $229$}{2}
  \z{ile dzielników ma $126$}{12}
  \z{ile dzielników ma $358$}{4}
  \z{ile dzielników ma $365$}{4}
  \z{ile dzielników ma $275$}{6}
  \z{ile dzielników ma $355$}{4}
  \z{ile dzielników ma $379$}{2}
\end{zestaw}

\btexpect{46}% ======================================================
%  GENERATED by tools/gen_sets.py -- do not edit.
%  Seed 2046; regenerate with `make sets`.
% ======================================================

\begin{zestaw}{Następna pierwsza III}{2}{7:20}{2}
  \zz{pierwsza liczba pierwsza po $304$}{307}
  \zz{pierwsza liczba pierwsza po $227$}{229}
  \zz{pierwsza liczba pierwsza po $346$}{347}
  \zz{pierwsza liczba pierwsza po $188$}{191}
  \zz{pierwsza liczba pierwsza po $228$}{229}
  \zz{pierwsza liczba pierwsza po $175$}{179}
  \zz{pierwsza liczba pierwsza po $121$}{127}
  \zz{pierwsza liczba pierwsza po $82$}{83}
  \zz{pierwsza liczba pierwsza po $272$}{277}
  \zz{pierwsza liczba pierwsza po $108$}{109}
  \zz{pierwsza liczba pierwsza po $159$}{163}
  \zz{pierwsza liczba pierwsza po $341$}{347}
  \zz{pierwsza liczba pierwsza po $259$}{263}
  \zz{pierwsza liczba pierwsza po $186$}{191}
  \zz{pierwsza liczba pierwsza po $393$}{397}
  \zz{pierwsza liczba pierwsza po $24$}{29}
  \zz{pierwsza liczba pierwsza po $172$}{173}
  \zz{pierwsza liczba pierwsza po $111$}{113}
  \zz{pierwsza liczba pierwsza po $35$}{37}
  \zz{pierwsza liczba pierwsza po $221$}{223}
  \btnc
  \zz{pierwsza liczba pierwsza po $290$}{293}
  \zz{pierwsza liczba pierwsza po $294$}{307}
  \zz{pierwsza liczba pierwsza po $189$}{191}
  \zz{pierwsza liczba pierwsza po $181$}{191}
  \zz{pierwsza liczba pierwsza po $287$}{293}
  \zz{pierwsza liczba pierwsza po $21$}{23}
  \zz{pierwsza liczba pierwsza po $213$}{223}
  \zz{pierwsza liczba pierwsza po $52$}{53}
  \zz{pierwsza liczba pierwsza po $146$}{149}
  \zz{pierwsza liczba pierwsza po $25$}{29}
  \zz{pierwsza liczba pierwsza po $69$}{71}
  \zz{pierwsza liczba pierwsza po $120$}{127}
  \zz{pierwsza liczba pierwsza po $299$}{307}
  \zz{pierwsza liczba pierwsza po $322$}{331}
  \zz{pierwsza liczba pierwsza po $79$}{83}
  \zz{pierwsza liczba pierwsza po $255$}{257}
  \zz{pierwsza liczba pierwsza po $133$}{137}
  \zz{pierwsza liczba pierwsza po $43$}{47}
  \zz{pierwsza liczba pierwsza po $298$}{307}
  \zz{pierwsza liczba pierwsza po $399$}{401}
\end{zestaw}

\btexpect{47}% ======================================================
%  GENERATED by tools/gen_sets.py -- do not edit.
%  Seed 2047; regenerate with `make sets`.
% ======================================================

\begin{zestaw}{Rozkład na czynniki IV}{2}{8:00}{3}
  \z{$348$ na czynniki}{$2^{2} \cdot 3 \cdot 29$}
  \z{$136$ na czynniki}{$2^{3} \cdot 17$}
  \z{$265$ na czynniki}{$5 \cdot 53$}
  \z{$92$ na czynniki}{$2^{2} \cdot 23$}
  \z{$237$ na czynniki}{$3 \cdot 79$}
  \z{$50$ na czynniki}{$2 \cdot 5^{2}$}
  \z{$231$ na czynniki}{$3 \cdot 7 \cdot 11$}
  \z{$252$ na czynniki}{$2^{2} \cdot 3^{2} \cdot 7$}
  \z{$42$ na czynniki}{$2 \cdot 3 \cdot 7$}
  \z{$80$ na czynniki}{$2^{4} \cdot 5$}
  \z{$105$ na czynniki}{$3 \cdot 5 \cdot 7$}
  \z{$69$ na czynniki}{$3 \cdot 23$}
  \z{$310$ na czynniki}{$2 \cdot 5 \cdot 31$}
  \z{$298$ na czynniki}{$2 \cdot 149$}
  \btnc
  \z{$262$ na czynniki}{$2 \cdot 131$}
  \z{$58$ na czynniki}{$2 \cdot 29$}
  \z{$213$ na czynniki}{$3 \cdot 71$}
  \z{$261$ na czynniki}{$3^{2} \cdot 29$}
  \z{$336$ na czynniki}{$2^{4} \cdot 3 \cdot 7$}
  \z{$90$ na czynniki}{$2 \cdot 3^{2} \cdot 5$}
  \z{$295$ na czynniki}{$5 \cdot 59$}
  \z{$115$ na czynniki}{$5 \cdot 23$}
  \z{$78$ na czynniki}{$2 \cdot 3 \cdot 13$}
  \z{$147$ na czynniki}{$3 \cdot 7^{2}$}
  \z{$342$ na czynniki}{$2 \cdot 3^{2} \cdot 19$}
  \z{$222$ na czynniki}{$2 \cdot 3 \cdot 37$}
  \z{$266$ na czynniki}{$2 \cdot 7 \cdot 19$}
  \btnc
  \z{$286$ na czynniki}{$2 \cdot 11 \cdot 13$}
  \z{$350$ na czynniki}{$2 \cdot 5^{2} \cdot 7$}
  \z{$196$ na czynniki}{$2^{2} \cdot 7^{2}$}
  \z{$38$ na czynniki}{$2 \cdot 19$}
  \z{$99$ na czynniki}{$3^{2} \cdot 11$}
  \z{$75$ na czynniki}{$3 \cdot 5^{2}$}
  \z{$221$ na czynniki}{$13 \cdot 17$}
  \z{$245$ na czynniki}{$5 \cdot 7^{2}$}
  \z{$380$ na czynniki}{$2^{2} \cdot 5 \cdot 19$}
  \z{$192$ na czynniki}{$2^{6} \cdot 3$}
  \z{$351$ na czynniki}{$3^{3} \cdot 13$}
  \z{$309$ na czynniki}{$3 \cdot 103$}
  \z{$171$ na czynniki}{$3^{2} \cdot 19$}
\end{zestaw}

\btexpect{48}% ======================================================
%  GENERATED by tools/gen_sets.py -- do not edit.
%  Seed 2048; regenerate with `make sets`.
% ======================================================

\begin{zestaw}{Czy pierwsza IV}{2}{6:40}{2}
  \z{czy $227$ jest pierwsza?}{tak}
  \z{czy $143$ jest pierwsza?}{nie}
  \z{czy $373$ jest pierwsza?}{tak}
  \z{czy $133$ jest pierwsza?}{nie}
  \z{czy $77$ jest pierwsza?}{nie}
  \z{czy $121$ jest pierwsza?}{nie}
  \z{czy $323$ jest pierwsza?}{nie}
  \z{czy $299$ jest pierwsza?}{nie}
  \z{czy $193$ jest pierwsza?}{tak}
  \z{czy $49$ jest pierwsza?}{nie}
  \z{czy $127$ jest pierwsza?}{tak}
  \z{czy $239$ jest pierwsza?}{tak}
  \z{czy $247$ jest pierwsza?}{nie}
  \z{czy $101$ jest pierwsza?}{tak}
  \z{czy $161$ jest pierwsza?}{nie}
  \z{czy $209$ jest pierwsza?}{nie}
  \z{czy $131$ jest pierwsza?}{tak}
  \z{czy $119$ jest pierwsza?}{nie}
  \z{czy $53$ jest pierwsza?}{tak}
  \z{czy $67$ jest pierwsza?}{tak}
  \btnc
  \z{czy $253$ jest pierwsza?}{nie}
  \z{czy $277$ jest pierwsza?}{tak}
  \z{czy $361$ jest pierwsza?}{nie}
  \z{czy $349$ jest pierwsza?}{tak}
  \z{czy $139$ jest pierwsza?}{tak}
  \z{czy $353$ jest pierwsza?}{tak}
  \z{czy $187$ jest pierwsza?}{nie}
  \z{czy $347$ jest pierwsza?}{tak}
  \z{czy $229$ jest pierwsza?}{tak}
  \z{czy $73$ jest pierwsza?}{tak}
  \z{czy $137$ jest pierwsza?}{tak}
  \z{czy $151$ jest pierwsza?}{tak}
  \z{czy $307$ jest pierwsza?}{tak}
  \z{czy $383$ jest pierwsza?}{tak}
  \z{czy $293$ jest pierwsza?}{tak}
  \z{czy $379$ jest pierwsza?}{tak}
  \z{czy $167$ jest pierwsza?}{tak}
  \z{czy $59$ jest pierwsza?}{tak}
  \z{czy $257$ jest pierwsza?}{tak}
  \z{czy $389$ jest pierwsza?}{tak}
\end{zestaw}



\chapter{Blok III — Ułamki i potęgi}

Pierwszy tom dodaje ułamki. Ten je mnoży i dzieli — a dzielenie przez ułamek
daje wynik \emph{większy} od tego, od czego zacząłeś, i to musi przestać
zaskakiwać. Do tego wykładnik ujemny, który jest odwrotnością, a nie liczbą
ujemną, wykładnik ułamkowy, logarytmy i notacja naukowa.

% GENERATED by tools/gen_sets.py -- do not edit.
\btexpect{49}% ======================================================
%  GENERATED by tools/gen_sets.py -- do not edit.
%  Seed 2049; regenerate with `make sets`.
% ======================================================

\begin{zestaw}{Mnożenie ułamków}{3}{6:40}{3}
  \z{$2/12 \times 5/6$}{5/36}
  \z{$1/9 \times 2/9$}{2/81}
  \z{$7/9 \times 2/8$}{7/36}
  \z{$8/9 \times 3/4$}{2/3}
  \z{$2/6 \times 1/2$}{1/6}
  \z{$2/12 \times 5/8$}{5/48}
  \z{$3/5 \times 5/9$}{1/3}
  \z{$2/9 \times 1/8$}{1/36}
  \z{$4/5 \times 4/5$}{16/25}
  \z{$2/12 \times 1/6$}{1/36}
  \z{$1/2 \times 4/9$}{2/9}
  \z{$1/8 \times 8/12$}{1/12}
  \z{$7/9 \times 3/4$}{7/12}
  \z{$5/8 \times 5/12$}{25/96}
  \btnc
  \z{$1/3 \times 1/4$}{1/12}
  \z{$2/5 \times 1/3$}{2/15}
  \z{$3/4 \times 4/6$}{1/2}
  \z{$4/12 \times 9/10$}{3/10}
  \z{$1/12 \times 1/2$}{1/24}
  \z{$5/12 \times 1/2$}{5/24}
  \z{$4/5 \times 5/8$}{1/2}
  \z{$1/4 \times 2/4$}{1/8}
  \z{$2/8 \times 1/10$}{1/40}
  \z{$10/12 \times 4/6$}{5/9}
  \z{$7/10 \times 2/3$}{7/15}
  \z{$2/3 \times 2/9$}{4/27}
  \z{$1/12 \times 4/6$}{1/18}
  \btnc
  \z{$1/3 \times 4/5$}{4/15}
  \z{$3/8 \times 5/8$}{15/64}
  \z{$3/4 \times 8/9$}{2/3}
  \z{$1/2 \times 4/6$}{1/3}
  \z{$5/6 \times 1/12$}{5/72}
  \z{$1/9 \times 2/3$}{2/27}
  \z{$1/2 \times 3/12$}{1/8}
  \z{$5/9 \times 2/8$}{5/36}
  \z{$1/12 \times 10/12$}{5/72}
  \z{$1/5 \times 6/9$}{2/15}
  \z{$1/3 \times 4/8$}{1/6}
  \z{$1/6 \times 11/12$}{11/72}
  \z{$4/6 \times 7/8$}{7/12}
\end{zestaw}

\btexpect{50}% ======================================================
%  GENERATED by tools/gen_sets.py -- do not edit.
%  Seed 2050; regenerate with `make sets`.
% ======================================================

\begin{zestaw}{Logarytmy}{3}{6:00}{3}
  \z{$\log_{7} 16\,807$}{5}
  \z{$\log_{5} 78\,125$}{7}
  \z{$\log_{6} 46\,656$}{6}
  \z{$\log_{8} 32\,768$}{5}
  \z{$\log_{8} 512$}{3}
  \z{$\log_{5} 15\,625$}{6}
  \z{$\log_{3} 59\,049$}{10}
  \z{$\log_{9} 59\,049$}{5}
  \z{$\log_{7} 49$}{2}
  \z{$\log_{10} 100\,000$}{5}
  \z{$\log_{5} 25$}{2}
  \z{$\log_{4} 65\,536$}{8}
  \z{$\log_{2} 32$}{5}
  \z{$\log_{2} 65\,536$}{16}
  \btnc
  \z{$\log_{10} 10\,000$}{4}
  \z{$\log_{2} 4$}{2}
  \z{$\log_{3} 9$}{2}
  \z{$\log_{2} 1\,024$}{10}
  \z{$\log_{4} 16\,384$}{7}
  \z{$\log_{3} 19\,683$}{9}
  \z{$\log_{2} 64$}{6}
  \z{$\log_{7} 2\,401$}{4}
  \z{$\log_{5} 625$}{4}
  \z{$\log_{2} 8$}{3}
  \z{$\log_{10} 100$}{2}
  \z{$\log_{9} 6\,561$}{4}
  \z{$\log_{10} 1\,000$}{3}
  \btnc
  \z{$\log_{4} 16$}{2}
  \z{$\log_{3} 729$}{6}
  \z{$\log_{4} 64$}{3}
  \z{$\log_{2} 16$}{4}
  \z{$\log_{2} 32\,768$}{15}
  \z{$\log_{2} 8\,192$}{13}
  \z{$\log_{5} 3\,125$}{5}
  \z{$\log_{2} 2\,048$}{11}
  \z{$\log_{5} 125$}{3}
  \z{$\log_{6} 216$}{3}
  \z{$\log_{3} 243$}{5}
  \z{$\log_{8} 64$}{2}
  \z{$\log_{2} 4\,096$}{12}
\end{zestaw}

\btexpect{51}% ======================================================
%  GENERATED by tools/gen_sets.py -- do not edit.
%  Seed 2051; regenerate with `make sets`.
% ======================================================

\begin{zestaw}{Potęgi ujemne i zero}{3}{5:20}{3}
  \z{$5^{-4}$}{1/625}
  \z{$7^{-2}$}{1/49}
  \z{$10^{-2}$}{1/100}
  \z{$7^{-1}$}{1/7}
  \z{$7^{0}$}{1}
  \z{$8^{-4}$}{1/4096}
  \z{$2^{-1}$}{1/2}
  \z{$9^{-2}$}{1/81}
  \z{$2^{-4}$}{1/16}
  \z{$9^{0}$}{1}
  \z{$8^{-3}$}{1/512}
  \z{$8^{-1}$}{1/8}
  \z{$4^{-4}$}{1/256}
  \z{$9^{-4}$}{1/6561}
  \btnc
  \z{$10^{-4}$}{1/10000}
  \z{$6^{-1}$}{1/6}
  \z{$4^{-3}$}{1/64}
  \z{$10^{0}$}{1}
  \z{$8^{0}$}{1}
  \z{$6^{-2}$}{1/36}
  \z{$6^{0}$}{1}
  \z{$6^{-3}$}{1/216}
  \z{$7^{-3}$}{1/343}
  \z{$5^{-3}$}{1/125}
  \z{$10^{-1}$}{1/10}
  \z{$9^{-1}$}{1/9}
  \z{$9^{-3}$}{1/729}
  \btnc
  \z{$3^{-4}$}{1/81}
  \z{$5^{-2}$}{1/25}
  \z{$3^{0}$}{1}
  \z{$6^{-4}$}{1/1296}
  \z{$2^{-3}$}{1/8}
  \z{$4^{-1}$}{1/4}
  \z{$10^{-3}$}{1/1000}
  \z{$2^{-2}$}{1/4}
  \z{$3^{-1}$}{1/3}
  \z{$4^{-2}$}{1/16}
  \z{$3^{-3}$}{1/27}
  \z{$3^{-2}$}{1/9}
  \z{$4^{0}$}{1}
\end{zestaw}

\btexpect{52}% ======================================================
%  GENERATED by tools/gen_sets.py -- do not edit.
%  Seed 2052; regenerate with `make sets`.
% ======================================================

\begin{zestaw}{Ułamek na dziesiętny}{3}{6:40}{3}
  \z{$17/20$ dziesiętnie}{0,85}
  \z{$22/25$ dziesiętnie}{0,88}
  \z{$42/50$ dziesiętnie}{0,84}
  \z{$1/2$ dziesiętnie}{0,5}
  \z{$2/4$ dziesiętnie}{0,5}
  \z{$3/4$ dziesiętnie}{0,75}
  \z{$4/5$ dziesiętnie}{0,8}
  \z{$7/40$ dziesiętnie}{0,175}
  \z{$1/8$ dziesiętnie}{0,125}
  \z{$7/20$ dziesiętnie}{0,35}
  \z{$4/8$ dziesiętnie}{0,5}
  \z{$9/20$ dziesiętnie}{0,45}
  \z{$27/50$ dziesiętnie}{0,54}
  \z{$7/8$ dziesiętnie}{0,875}
  \btnc
  \z{$36/50$ dziesiętnie}{0,72}
  \z{$12/25$ dziesiętnie}{0,48}
  \z{$23/25$ dziesiętnie}{0,92}
  \z{$19/25$ dziesiętnie}{0,76}
  \z{$38/40$ dziesiętnie}{0,95}
  \z{$1/4$ dziesiętnie}{0,25}
  \z{$7/50$ dziesiętnie}{0,14}
  \z{$3/10$ dziesiętnie}{0,3}
  \z{$19/20$ dziesiętnie}{0,95}
  \z{$1/5$ dziesiętnie}{0,2}
  \z{$4/10$ dziesiętnie}{0,4}
  \z{$3/20$ dziesiętnie}{0,15}
  \z{$47/50$ dziesiętnie}{0,94}
  \btnc
  \z{$21/25$ dziesiętnie}{0,84}
  \z{$20/25$ dziesiętnie}{0,8}
  \z{$12/20$ dziesiętnie}{0,6}
  \z{$24/25$ dziesiętnie}{0,96}
  \z{$23/50$ dziesiętnie}{0,46}
  \z{$17/40$ dziesiętnie}{0,425}
  \z{$2/8$ dziesiętnie}{0,25}
  \z{$6/10$ dziesiętnie}{0,6}
  \z{$6/8$ dziesiętnie}{0,75}
  \z{$2/20$ dziesiętnie}{0,1}
  \z{$1/20$ dziesiętnie}{0,05}
  \z{$2/5$ dziesiętnie}{0,4}
  \z{$3/8$ dziesiętnie}{0,375}
\end{zestaw}

\btexpect{53}% ======================================================
%  GENERATED by tools/gen_sets.py -- do not edit.
%  Seed 2053; regenerate with `make sets`.
% ======================================================

\begin{zestaw}{Dzielenie ułamków}{3}{7:20}{3}
  \z{$1/6 \div 2/6$}{1/2}
  \z{$1/5 \div 5/10$}{2/5}
  \z{$1/10 \div 2/12$}{3/5}
  \z{$2/6 \div 2/5$}{5/6}
  \z{$3/6 \div 2/6$}{3/2}
  \z{$1/3 \div 1/4$}{4/3}
  \z{$5/6 \div 1/2$}{5/3}
  \z{$1/5 \div 9/10$}{2/9}
  \z{$9/10 \div 1/2$}{9/5}
  \z{$2/9 \div 3/5$}{10/27}
  \z{$10/12 \div 1/3$}{5/2}
  \z{$4/12 \div 2/3$}{1/2}
  \z{$4/6 \div 1/3$}{2}
  \z{$10/12 \div 2/6$}{5/2}
  \btnc
  \z{$3/4 \div 1/5$}{15/4}
  \z{$1/2 \div 2/3$}{3/4}
  \z{$4/6 \div 1/5$}{10/3}
  \z{$1/4 \div 7/9$}{9/28}
  \z{$1/10 \div 3/5$}{1/6}
  \z{$2/3 \div 3/5$}{10/9}
  \z{$2/4 \div 5/9$}{9/10}
  \z{$1/4 \div 4/9$}{9/16}
  \z{$2/5 \div 1/9$}{18/5}
  \z{$5/12 \div 4/6$}{5/8}
  \z{$4/8 \div 2/3$}{3/4}
  \z{$3/6 \div 7/9$}{9/14}
  \z{$7/8 \div 1/5$}{35/8}
  \btnc
  \z{$4/6 \div 11/12$}{8/11}
  \z{$4/8 \div 10/12$}{3/5}
  \z{$6/10 \div 3/10$}{2}
  \z{$8/12 \div 2/6$}{2}
  \z{$1/4 \div 3/5$}{5/12}
  \z{$2/3 \div 3/4$}{8/9}
  \z{$3/10 \div 1/12$}{18/5}
  \z{$1/2 \div 8/12$}{3/4}
  \z{$2/4 \div 8/9$}{9/16}
  \z{$3/8 \div 2/9$}{27/16}
  \z{$1/3 \div 8/12$}{1/2}
  \z{$1/3 \div 2/9$}{3/2}
  \z{$11/12 \div 2/9$}{33/8}
\end{zestaw}

\btexpect{54}% ======================================================
%  GENERATED by tools/gen_sets.py -- do not edit.
%  Seed 2054; regenerate with `make sets`.
% ======================================================

\begin{zestaw}{Notacja naukowa}{3}{7:20}{2}
  \z{$5 \cdot 10^{4} \times 3 \cdot 10^{7}$}{$1{,}5 \cdot 10^{12}$}
  \z{$6 \cdot 10^{3} \times 5 \cdot 10^{4}$}{$3 \cdot 10^{8}$}
  \z{$4 \cdot 10^{4} \times 5 \cdot 10^{2}$}{$2 \cdot 10^{7}$}
  \z{$4 \cdot 10^{7} \times 2 \cdot 10^{6}$}{$8 \cdot 10^{13}$}
  \z{$8 \cdot 10^{7} \times 6 \cdot 10^{7}$}{$4{,}8 \cdot 10^{15}$}
  \z{$4 \cdot 10^{5} \times 4 \cdot 10^{7}$}{$1{,}6 \cdot 10^{13}$}
  \z{$2 \cdot 10^{2} \times 8 \cdot 10^{6}$}{$1{,}6 \cdot 10^{9}$}
  \z{$4 \cdot 10^{3} \times 9 \cdot 10^{2}$}{$3{,}6 \cdot 10^{6}$}
  \z{$4 \cdot 10^{5} \times 3 \cdot 10^{2}$}{$1{,}2 \cdot 10^{8}$}
  \z{$8 \cdot 10^{6} \times 6 \cdot 10^{4}$}{$4{,}8 \cdot 10^{11}$}
  \z{$6 \cdot 10^{3} \times 6 \cdot 10^{5}$}{$3{,}6 \cdot 10^{9}$}
  \z{$7 \cdot 10^{6} \times 3 \cdot 10^{8}$}{$2{,}1 \cdot 10^{15}$}
  \z{$4 \cdot 10^{8} \times 6 \cdot 10^{6}$}{$2{,}4 \cdot 10^{15}$}
  \z{$2 \cdot 10^{6} \times 8 \cdot 10^{2}$}{$1{,}6 \cdot 10^{9}$}
  \z{$2 \cdot 10^{5} \times 4 \cdot 10^{7}$}{$8 \cdot 10^{12}$}
  \z{$9 \cdot 10^{3} \times 3 \cdot 10^{5}$}{$2{,}7 \cdot 10^{9}$}
  \z{$3 \cdot 10^{4} \times 2 \cdot 10^{8}$}{$6 \cdot 10^{12}$}
  \z{$9 \cdot 10^{5} \times 9 \cdot 10^{3}$}{$8{,}1 \cdot 10^{9}$}
  \z{$5 \cdot 10^{5} \times 3 \cdot 10^{7}$}{$1{,}5 \cdot 10^{13}$}
  \z{$2 \cdot 10^{3} \times 3 \cdot 10^{6}$}{$6 \cdot 10^{9}$}
  \btnc
  \z{$7 \cdot 10^{2} \times 3 \cdot 10^{4}$}{$2{,}1 \cdot 10^{7}$}
  \z{$2 \cdot 10^{6} \times 3 \cdot 10^{4}$}{$6 \cdot 10^{10}$}
  \z{$5 \cdot 10^{4} \times 9 \cdot 10^{6}$}{$4{,}5 \cdot 10^{11}$}
  \z{$5 \cdot 10^{5} \times 4 \cdot 10^{7}$}{$2 \cdot 10^{13}$}
  \z{$5 \cdot 10^{2} \times 6 \cdot 10^{4}$}{$3 \cdot 10^{7}$}
  \z{$5 \cdot 10^{8} \times 2 \cdot 10^{4}$}{$1 \cdot 10^{13}$}
  \z{$3 \cdot 10^{4} \times 9 \cdot 10^{8}$}{$2{,}7 \cdot 10^{13}$}
  \z{$6 \cdot 10^{4} \times 3 \cdot 10^{2}$}{$1{,}8 \cdot 10^{7}$}
  \z{$9 \cdot 10^{2} \times 4 \cdot 10^{4}$}{$3{,}6 \cdot 10^{7}$}
  \z{$8 \cdot 10^{5} \times 8 \cdot 10^{4}$}{$6{,}4 \cdot 10^{10}$}
  \z{$9 \cdot 10^{8} \times 2 \cdot 10^{5}$}{$1{,}8 \cdot 10^{14}$}
  \z{$6 \cdot 10^{6} \times 4 \cdot 10^{3}$}{$2{,}4 \cdot 10^{10}$}
  \z{$4 \cdot 10^{4} \times 4 \cdot 10^{8}$}{$1{,}6 \cdot 10^{13}$}
  \z{$9 \cdot 10^{2} \times 4 \cdot 10^{8}$}{$3{,}6 \cdot 10^{11}$}
  \z{$7 \cdot 10^{7} \times 6 \cdot 10^{5}$}{$4{,}2 \cdot 10^{13}$}
  \z{$9 \cdot 10^{4} \times 2 \cdot 10^{3}$}{$1{,}8 \cdot 10^{8}$}
  \z{$6 \cdot 10^{5} \times 9 \cdot 10^{5}$}{$5{,}4 \cdot 10^{11}$}
  \z{$6 \cdot 10^{3} \times 9 \cdot 10^{5}$}{$5{,}4 \cdot 10^{9}$}
  \z{$4 \cdot 10^{2} \times 2 \cdot 10^{5}$}{$8 \cdot 10^{7}$}
  \z{$9 \cdot 10^{3} \times 7 \cdot 10^{7}$}{$6{,}3 \cdot 10^{11}$}
\end{zestaw}

\btexpect{55}% ======================================================
%  GENERATED by tools/gen_sets.py -- do not edit.
%  Seed 2055; regenerate with `make sets`.
% ======================================================

\begin{zestaw}{Wykładnik ułamkowy}{3}{7:20}{3}
  \z{$25^{2/2}$}{25}
  \z{$16^{1/2}$}{4}
  \z{$144^{1/2}$}{12}
  \z{$16^{2/4}$}{4}
  \z{$144^{2/2}$}{144}
  \z{$27^{3/3}$}{27}
  \z{$9^{1/2}$}{3}
  \z{$100^{3/2}$}{1\,000}
  \z{$144^{3/2}$}{1\,728}
  \z{$121^{1/2}$}{11}
  \z{$25^{3/2}$}{125}
  \z{$8^{1/3}$}{2}
  \z{$64^{2/6}$}{4}
  \z{$128^{2/7}$}{4}
  \btnc
  \z{$16^{3/4}$}{8}
  \z{$32^{1/5}$}{2}
  \z{$64^{1/2}$}{8}
  \z{$27^{2/3}$}{9}
  \z{$128^{1/7}$}{2}
  \z{$121^{3/2}$}{1\,331}
  \z{$125^{3/3}$}{125}
  \z{$125^{2/3}$}{25}
  \z{$121^{2/2}$}{121}
  \z{$4^{3/2}$}{8}
  \z{$125^{1/3}$}{5}
  \z{$81^{1/2}$}{9}
  \z{$49^{2/2}$}{49}
  \btnc
  \z{$81^{3/2}$}{729}
  \z{$27^{1/3}$}{3}
  \z{$9^{3/2}$}{27}
  \z{$16^{2/2}$}{16}
  \z{$49^{1/2}$}{7}
  \z{$32^{2/5}$}{4}
  \z{$4^{2/2}$}{4}
  \z{$16^{3/2}$}{64}
  \z{$64^{3/3}$}{64}
  \z{$9^{2/2}$}{9}
  \z{$25^{1/2}$}{5}
  \z{$169^{3/2}$}{2\,197}
  \z{$64^{1/3}$}{4}
\end{zestaw}

\btexpect{56}% ======================================================
%  GENERATED by tools/gen_sets.py -- do not edit.
%  Seed 2056; regenerate with `make sets`.
% ======================================================

\begin{zestaw}{Ułamek na procent}{3}{6:40}{3}
  \z{$38/40$ na procent}{95\%}
  \z{$9/10$ na procent}{90\%}
  \z{$7/20$ na procent}{35\%}
  \z{$1/5$ na procent}{20\%}
  \z{$1/2$ na procent}{50\%}
  \z{$1/8$ na procent}{12,5\%}
  \z{$1/20$ na procent}{5\%}
  \z{$17/40$ na procent}{42,5\%}
  \z{$5/10$ na procent}{50\%}
  \z{$13/25$ na procent}{52\%}
  \z{$4/5$ na procent}{80\%}
  \z{$16/50$ na procent}{32\%}
  \z{$7/8$ na procent}{87,5\%}
  \z{$3/8$ na procent}{37,5\%}
  \btnc
  \z{$3/4$ na procent}{75\%}
  \z{$9/25$ na procent}{36\%}
  \z{$1/4$ na procent}{25\%}
  \z{$29/40$ na procent}{72,5\%}
  \z{$43/50$ na procent}{86\%}
  \z{$2/8$ na procent}{25\%}
  \z{$19/20$ na procent}{95\%}
  \z{$1/10$ na procent}{10\%}
  \z{$1/16$ na procent}{6,25\%}
  \z{$10/20$ na procent}{50\%}
  \z{$19/50$ na procent}{38\%}
  \z{$2/4$ na procent}{50\%}
  \z{$12/16$ na procent}{75\%}
  \btnc
  \z{$3/20$ na procent}{15\%}
  \z{$3/10$ na procent}{30\%}
  \z{$3/5$ na procent}{60\%}
  \z{$6/8$ na procent}{75\%}
  \z{$2/10$ na procent}{20\%}
  \z{$14/25$ na procent}{56\%}
  \z{$14/16$ na procent}{87,5\%}
  \z{$4/8$ na procent}{50\%}
  \z{$8/10$ na procent}{80\%}
  \z{$4/16$ na procent}{25\%}
  \z{$5/8$ na procent}{62,5\%}
  \z{$2/5$ na procent}{40\%}
  \z{$46/50$ na procent}{92\%}
\end{zestaw}

\btexpect{57}% ======================================================
%  GENERATED by tools/gen_sets.py -- do not edit.
%  Seed 2057; regenerate with `make sets`.
% ======================================================

\begin{zestaw}{Małe dziesiętne}{3}{6:00}{3}
  \z{$0{,}009 \times 3\,000$}{27}
  \z{$0{,}04 \times 2\,000$}{80}
  \z{$0{,}08 \times 200$}{16}
  \z{$0{,}006 \times 10\,000$}{60}
  \z{$0{,}06 \times 1\,300$}{78}
  \z{$0{,}004 \times 3\,000$}{12}
  \z{$0{,}03 \times 1\,900$}{57}
  \z{$0{,}08 \times 2\,000$}{160}
  \z{$0{,}007 \times 20\,000$}{140}
  \z{$0{,}009 \times 18\,000$}{162}
  \z{$0{,}007 \times 17\,000$}{119}
  \z{$0{,}006 \times 9\,000$}{54}
  \z{$0{,}003 \times 14\,000$}{42}
  \z{$0{,}006 \times 3\,000$}{18}
  \btnc
  \z{$0{,}008 \times 11\,000$}{88}
  \z{$0{,}06 \times 700$}{42}
  \z{$0{,}09 \times 1\,800$}{162}
  \z{$0{,}008 \times 20\,000$}{160}
  \z{$0{,}08 \times 1\,400$}{112}
  \z{$0{,}02 \times 1\,900$}{38}
  \z{$0{,}004 \times 18\,000$}{72}
  \z{$0{,}006 \times 7\,000$}{42}
  \z{$0{,}06 \times 400$}{24}
  \z{$0{,}004 \times 16\,000$}{64}
  \z{$0{,}06 \times 1\,200$}{72}
  \z{$0{,}07 \times 1\,700$}{119}
  \z{$0{,}09 \times 400$}{36}
  \btnc
  \z{$0{,}03 \times 1\,800$}{54}
  \z{$0{,}002 \times 9\,000$}{18}
  \z{$0{,}004 \times 20\,000$}{80}
  \z{$0{,}09 \times 1\,600$}{144}
  \z{$0{,}04 \times 1\,100$}{44}
  \z{$0{,}005 \times 2\,000$}{10}
  \z{$0{,}05 \times 1\,300$}{65}
  \z{$0{,}09 \times 600$}{54}
  \z{$0{,}005 \times 14\,000$}{70}
  \z{$0{,}009 \times 13\,000$}{117}
  \z{$0{,}003 \times 11\,000$}{33}
  \z{$0{,}06 \times 600$}{36}
  \z{$0{,}009 \times 20\,000$}{180}
\end{zestaw}

\btexpect{58}% ======================================================
%  GENERATED by tools/gen_sets.py -- do not edit.
%  Seed 2058; regenerate with `make sets`.
% ======================================================

\begin{zestaw}{Mnożenie ułamków II}{3}{6:40}{3}
  \z{$1/4 \times 1/2$}{1/8}
  \z{$2/5 \times 2/8$}{1/10}
  \z{$7/10 \times 4/6$}{7/15}
  \z{$1/5 \times 5/9$}{1/9}
  \z{$1/2 \times 1/10$}{1/20}
  \z{$7/8 \times 3/5$}{21/40}
  \z{$2/3 \times 4/5$}{8/15}
  \z{$2/8 \times 4/6$}{1/6}
  \z{$2/3 \times 1/4$}{1/6}
  \z{$2/8 \times 2/9$}{1/18}
  \z{$3/5 \times 4/5$}{12/25}
  \z{$4/6 \times 5/9$}{10/27}
  \z{$6/12 \times 5/9$}{5/18}
  \z{$2/4 \times 9/10$}{9/20}
  \btnc
  \z{$4/5 \times 2/3$}{8/15}
  \z{$2/3 \times 3/5$}{2/5}
  \z{$4/12 \times 1/4$}{1/12}
  \z{$5/12 \times 3/4$}{5/16}
  \z{$5/6 \times 1/2$}{5/12}
  \z{$1/2 \times 4/9$}{2/9}
  \z{$6/9 \times 6/8$}{1/2}
  \z{$6/8 \times 1/2$}{3/8}
  \z{$10/12 \times 9/10$}{3/4}
  \z{$1/2 \times 3/10$}{3/20}
  \z{$3/12 \times 4/5$}{1/5}
  \z{$1/4 \times 3/4$}{3/16}
  \z{$5/12 \times 7/12$}{35/144}
  \btnc
  \z{$7/10 \times 1/2$}{7/20}
  \z{$1/2 \times 2/4$}{1/4}
  \z{$3/5 \times 1/3$}{1/5}
  \z{$1/5 \times 6/9$}{2/15}
  \z{$3/8 \times 4/8$}{3/16}
  \z{$1/2 \times 1/4$}{1/8}
  \z{$1/6 \times 3/5$}{1/10}
  \z{$1/2 \times 3/4$}{3/8}
  \z{$1/2 \times 2/3$}{1/3}
  \z{$1/3 \times 1/6$}{1/18}
  \z{$7/12 \times 4/5$}{7/15}
  \z{$3/4 \times 7/10$}{21/40}
  \z{$3/10 \times 7/10$}{21/100}
\end{zestaw}

\btexpect{59}% ======================================================
%  GENERATED by tools/gen_sets.py -- do not edit.
%  Seed 2059; regenerate with `make sets`.
% ======================================================

\begin{zestaw}{Logarytmy II}{3}{6:00}{3}
  \z{$\log_{8} 512$}{3}
  \z{$\log_{9} 59\,049$}{5}
  \z{$\log_{6} 46\,656$}{6}
  \z{$\log_{4} 65\,536$}{8}
  \z{$\log_{3} 81$}{4}
  \z{$\log_{3} 59\,049$}{10}
  \z{$\log_{4} 1\,024$}{5}
  \z{$\log_{4} 256$}{4}
  \z{$\log_{3} 243$}{5}
  \z{$\log_{6} 216$}{3}
  \z{$\log_{8} 32\,768$}{5}
  \z{$\log_{9} 729$}{3}
  \z{$\log_{5} 125$}{3}
  \z{$\log_{10} 100\,000$}{5}
  \btnc
  \z{$\log_{5} 15\,625$}{6}
  \z{$\log_{5} 78\,125$}{7}
  \z{$\log_{7} 49$}{2}
  \z{$\log_{2} 16$}{4}
  \z{$\log_{9} 6\,561$}{4}
  \z{$\log_{4} 64$}{3}
  \z{$\log_{8} 64$}{2}
  \z{$\log_{7} 16\,807$}{5}
  \z{$\log_{6} 7\,776$}{5}
  \z{$\log_{2} 512$}{9}
  \z{$\log_{2} 32\,768$}{15}
  \z{$\log_{2} 65\,536$}{16}
  \z{$\log_{2} 8$}{3}
  \btnc
  \z{$\log_{3} 6\,561$}{8}
  \z{$\log_{2} 32$}{5}
  \z{$\log_{9} 81$}{2}
  \z{$\log_{4} 16$}{2}
  \z{$\log_{10} 10\,000$}{4}
  \z{$\log_{8} 4\,096$}{4}
  \z{$\log_{2} 16\,384$}{14}
  \z{$\log_{5} 3\,125$}{5}
  \z{$\log_{10} 1\,000$}{3}
  \z{$\log_{7} 343$}{3}
  \z{$\log_{2} 1\,024$}{10}
  \z{$\log_{7} 2\,401$}{4}
  \z{$\log_{2} 2\,048$}{11}
\end{zestaw}

\btexpect{60}% ======================================================
%  GENERATED by tools/gen_sets.py -- do not edit.
%  Seed 2060; regenerate with `make sets`.
% ======================================================

\begin{zestaw}{Potęgi ujemne i zero II}{3}{5:20}{3}
  \z{$6^{-4}$}{1/1296}
  \z{$2^{-1}$}{1/2}
  \z{$8^{-4}$}{1/4096}
  \z{$3^{-3}$}{1/27}
  \z{$5^{0}$}{1}
  \z{$10^{-2}$}{1/100}
  \z{$7^{-2}$}{1/49}
  \z{$8^{-1}$}{1/8}
  \z{$10^{-3}$}{1/1000}
  \z{$8^{-2}$}{1/64}
  \z{$9^{-1}$}{1/9}
  \z{$9^{-3}$}{1/729}
  \z{$4^{-3}$}{1/64}
  \z{$2^{-4}$}{1/16}
  \btnc
  \z{$3^{0}$}{1}
  \z{$3^{-4}$}{1/81}
  \z{$6^{-2}$}{1/36}
  \z{$3^{-1}$}{1/3}
  \z{$8^{-3}$}{1/512}
  \z{$5^{-4}$}{1/625}
  \z{$7^{-3}$}{1/343}
  \z{$7^{0}$}{1}
  \z{$6^{0}$}{1}
  \z{$5^{-1}$}{1/5}
  \z{$2^{0}$}{1}
  \z{$8^{0}$}{1}
  \z{$10^{0}$}{1}
  \btnc
  \z{$9^{0}$}{1}
  \z{$6^{-1}$}{1/6}
  \z{$7^{-4}$}{1/2401}
  \z{$9^{-2}$}{1/81}
  \z{$6^{-3}$}{1/216}
  \z{$4^{0}$}{1}
  \z{$2^{-2}$}{1/4}
  \z{$9^{-4}$}{1/6561}
  \z{$5^{-3}$}{1/125}
  \z{$7^{-1}$}{1/7}
  \z{$4^{-1}$}{1/4}
  \z{$4^{-2}$}{1/16}
  \z{$10^{-1}$}{1/10}
\end{zestaw}

\btexpect{61}% ======================================================
%  GENERATED by tools/gen_sets.py -- do not edit.
%  Seed 2061; regenerate with `make sets`.
% ======================================================

\begin{zestaw}{Ułamek na dziesiętny II}{3}{6:40}{3}
  \z{$10/20$ dziesiętnie}{0,5}
  \z{$8/40$ dziesiętnie}{0,2}
  \z{$3/20$ dziesiętnie}{0,15}
  \z{$24/25$ dziesiętnie}{0,96}
  \z{$2/4$ dziesiętnie}{0,5}
  \z{$1/8$ dziesiętnie}{0,125}
  \z{$7/25$ dziesiętnie}{0,28}
  \z{$3/8$ dziesiętnie}{0,375}
  \z{$1/2$ dziesiętnie}{0,5}
  \z{$12/25$ dziesiętnie}{0,48}
  \z{$7/8$ dziesiętnie}{0,875}
  \z{$4/10$ dziesiętnie}{0,4}
  \z{$7/20$ dziesiętnie}{0,35}
  \z{$9/10$ dziesiętnie}{0,9}
  \btnc
  \z{$2/20$ dziesiętnie}{0,1}
  \z{$20/25$ dziesiętnie}{0,8}
  \z{$20/50$ dziesiętnie}{0,4}
  \z{$1/5$ dziesiętnie}{0,2}
  \z{$5/20$ dziesiętnie}{0,25}
  \z{$40/50$ dziesiętnie}{0,8}
  \z{$16/25$ dziesiętnie}{0,64}
  \z{$3/4$ dziesiętnie}{0,75}
  \z{$11/40$ dziesiętnie}{0,275}
  \z{$1/40$ dziesiętnie}{0,025}
  \z{$1/4$ dziesiętnie}{0,25}
  \z{$2/25$ dziesiętnie}{0,08}
  \z{$6/10$ dziesiętnie}{0,6}
  \btnc
  \z{$8/50$ dziesiętnie}{0,16}
  \z{$7/50$ dziesiętnie}{0,14}
  \z{$2/10$ dziesiętnie}{0,2}
  \z{$11/25$ dziesiętnie}{0,44}
  \z{$3/5$ dziesiętnie}{0,6}
  \z{$18/25$ dziesiętnie}{0,72}
  \z{$14/20$ dziesiętnie}{0,7}
  \z{$2/5$ dziesiętnie}{0,4}
  \z{$32/40$ dziesiętnie}{0,8}
  \z{$17/50$ dziesiętnie}{0,34}
  \z{$25/40$ dziesiętnie}{0,625}
  \z{$1/10$ dziesiętnie}{0,1}
  \z{$5/8$ dziesiętnie}{0,625}
\end{zestaw}

\btexpect{62}% ======================================================
%  GENERATED by tools/gen_sets.py -- do not edit.
%  Seed 2062; regenerate with `make sets`.
% ======================================================

\begin{zestaw}{Dzielenie ułamków II}{3}{7:20}{3}
  \z{$5/8 \div 3/10$}{25/12}
  \z{$1/12 \div 9/10$}{5/54}
  \z{$6/12 \div 7/8$}{4/7}
  \z{$3/4 \div 1/10$}{15/2}
  \z{$1/2 \div 4/6$}{3/4}
  \z{$5/6 \div 2/5$}{25/12}
  \z{$3/8 \div 3/10$}{5/4}
  \z{$3/4 \div 1/6$}{9/2}
  \z{$1/2 \div 3/9$}{3/2}
  \z{$1/4 \div 3/5$}{5/12}
  \z{$1/12 \div 1/4$}{1/3}
  \z{$10/12 \div 6/9$}{5/4}
  \z{$2/3 \div 1/2$}{4/3}
  \z{$2/10 \div 1/5$}{1}
  \btnc
  \z{$9/12 \div 6/9$}{9/8}
  \z{$5/9 \div 5/9$}{1}
  \z{$1/3 \div 1/3$}{1}
  \z{$8/10 \div 2/10$}{4}
  \z{$2/4 \div 1/2$}{1}
  \z{$1/3 \div 3/10$}{10/9}
  \z{$3/4 \div 10/12$}{9/10}
  \z{$3/4 \div 4/5$}{15/16}
  \z{$1/2 \div 2/3$}{3/4}
  \z{$2/4 \div 6/9$}{3/4}
  \z{$7/12 \div 1/6$}{7/2}
  \z{$4/12 \div 3/4$}{4/9}
  \z{$1/6 \div 1/3$}{1/2}
  \btnc
  \z{$6/9 \div 1/3$}{2}
  \z{$6/10 \div 8/9$}{27/40}
  \z{$7/10 \div 1/4$}{14/5}
  \z{$3/4 \div 1/3$}{9/4}
  \z{$4/12 \div 1/3$}{1}
  \z{$5/9 \div 9/12$}{20/27}
  \z{$4/12 \div 6/8$}{4/9}
  \z{$1/2 \div 1/2$}{1}
  \z{$8/12 \div 2/6$}{2}
  \z{$3/6 \div 6/9$}{3/4}
  \z{$1/5 \div 1/12$}{12/5}
  \z{$1/2 \div 1/8$}{4}
  \z{$2/5 \div 3/5$}{2/3}
\end{zestaw}

\btexpect{63}% ======================================================
%  GENERATED by tools/gen_sets.py -- do not edit.
%  Seed 2063; regenerate with `make sets`.
% ======================================================

\begin{zestaw}{Notacja naukowa II}{3}{7:20}{2}
  \z{$4 \cdot 10^{6} \times 4 \cdot 10^{8}$}{$1{,}6 \cdot 10^{15}$}
  \z{$7 \cdot 10^{3} \times 3 \cdot 10^{2}$}{$2{,}1 \cdot 10^{6}$}
  \z{$5 \cdot 10^{4} \times 9 \cdot 10^{7}$}{$4{,}5 \cdot 10^{12}$}
  \z{$4 \cdot 10^{7} \times 5 \cdot 10^{8}$}{$2 \cdot 10^{16}$}
  \z{$9 \cdot 10^{2} \times 6 \cdot 10^{7}$}{$5{,}4 \cdot 10^{10}$}
  \z{$3 \cdot 10^{5} \times 9 \cdot 10^{4}$}{$2{,}7 \cdot 10^{10}$}
  \z{$9 \cdot 10^{5} \times 7 \cdot 10^{2}$}{$6{,}3 \cdot 10^{8}$}
  \z{$7 \cdot 10^{5} \times 5 \cdot 10^{5}$}{$3{,}5 \cdot 10^{11}$}
  \z{$7 \cdot 10^{3} \times 2 \cdot 10^{7}$}{$1{,}4 \cdot 10^{11}$}
  \z{$8 \cdot 10^{5} \times 6 \cdot 10^{5}$}{$4{,}8 \cdot 10^{11}$}
  \z{$9 \cdot 10^{3} \times 6 \cdot 10^{2}$}{$5{,}4 \cdot 10^{6}$}
  \z{$8 \cdot 10^{7} \times 5 \cdot 10^{7}$}{$4 \cdot 10^{15}$}
  \z{$9 \cdot 10^{5} \times 4 \cdot 10^{5}$}{$3{,}6 \cdot 10^{11}$}
  \z{$6 \cdot 10^{6} \times 8 \cdot 10^{8}$}{$4{,}8 \cdot 10^{15}$}
  \z{$3 \cdot 10^{7} \times 4 \cdot 10^{2}$}{$1{,}2 \cdot 10^{10}$}
  \z{$4 \cdot 10^{6} \times 6 \cdot 10^{6}$}{$2{,}4 \cdot 10^{13}$}
  \z{$5 \cdot 10^{5} \times 9 \cdot 10^{5}$}{$4{,}5 \cdot 10^{11}$}
  \z{$4 \cdot 10^{4} \times 2 \cdot 10^{5}$}{$8 \cdot 10^{9}$}
  \z{$2 \cdot 10^{4} \times 3 \cdot 10^{2}$}{$6 \cdot 10^{6}$}
  \z{$8 \cdot 10^{6} \times 2 \cdot 10^{6}$}{$1{,}6 \cdot 10^{13}$}
  \btnc
  \z{$2 \cdot 10^{7} \times 9 \cdot 10^{8}$}{$1{,}8 \cdot 10^{16}$}
  \z{$4 \cdot 10^{7} \times 3 \cdot 10^{3}$}{$1{,}2 \cdot 10^{11}$}
  \z{$9 \cdot 10^{4} \times 4 \cdot 10^{3}$}{$3{,}6 \cdot 10^{8}$}
  \z{$4 \cdot 10^{8} \times 4 \cdot 10^{8}$}{$1{,}6 \cdot 10^{17}$}
  \z{$3 \cdot 10^{4} \times 8 \cdot 10^{2}$}{$2{,}4 \cdot 10^{7}$}
  \z{$4 \cdot 10^{4} \times 9 \cdot 10^{8}$}{$3{,}6 \cdot 10^{13}$}
  \z{$6 \cdot 10^{8} \times 7 \cdot 10^{8}$}{$4{,}2 \cdot 10^{17}$}
  \z{$7 \cdot 10^{4} \times 5 \cdot 10^{8}$}{$3{,}5 \cdot 10^{13}$}
  \z{$6 \cdot 10^{7} \times 3 \cdot 10^{4}$}{$1{,}8 \cdot 10^{12}$}
  \z{$4 \cdot 10^{4} \times 9 \cdot 10^{4}$}{$3{,}6 \cdot 10^{9}$}
  \z{$6 \cdot 10^{6} \times 4 \cdot 10^{3}$}{$2{,}4 \cdot 10^{10}$}
  \z{$4 \cdot 10^{8} \times 3 \cdot 10^{3}$}{$1{,}2 \cdot 10^{12}$}
  \z{$7 \cdot 10^{4} \times 6 \cdot 10^{6}$}{$4{,}2 \cdot 10^{11}$}
  \z{$3 \cdot 10^{7} \times 5 \cdot 10^{3}$}{$1{,}5 \cdot 10^{11}$}
  \z{$3 \cdot 10^{8} \times 3 \cdot 10^{5}$}{$9 \cdot 10^{13}$}
  \z{$6 \cdot 10^{8} \times 7 \cdot 10^{5}$}{$4{,}2 \cdot 10^{14}$}
  \z{$4 \cdot 10^{6} \times 9 \cdot 10^{4}$}{$3{,}6 \cdot 10^{11}$}
  \z{$7 \cdot 10^{7} \times 2 \cdot 10^{3}$}{$1{,}4 \cdot 10^{11}$}
  \z{$9 \cdot 10^{7} \times 6 \cdot 10^{5}$}{$5{,}4 \cdot 10^{13}$}
  \z{$2 \cdot 10^{8} \times 7 \cdot 10^{8}$}{$1{,}4 \cdot 10^{17}$}
\end{zestaw}

\btexpect{64}% ======================================================
%  GENERATED by tools/gen_sets.py -- do not edit.
%  Seed 2064; regenerate with `make sets`.
% ======================================================

\begin{zestaw}{Wykładnik ułamkowy II}{3}{7:20}{3}
  \z{$49^{3/2}$}{343}
  \z{$9^{2/2}$}{9}
  \z{$121^{3/2}$}{1\,331}
  \z{$169^{2/2}$}{169}
  \z{$100^{2/2}$}{100}
  \z{$8^{1/3}$}{2}
  \z{$49^{2/2}$}{49}
  \z{$128^{3/7}$}{8}
  \z{$27^{2/3}$}{9}
  \z{$64^{1/3}$}{4}
  \z{$16^{1/4}$}{2}
  \z{$8^{2/3}$}{4}
  \z{$144^{1/2}$}{12}
  \z{$100^{3/2}$}{1\,000}
  \btnc
  \z{$16^{2/2}$}{16}
  \z{$4^{2/2}$}{4}
  \z{$81^{3/2}$}{729}
  \z{$64^{2/2}$}{64}
  \z{$121^{1/2}$}{11}
  \z{$144^{2/2}$}{144}
  \z{$64^{1/6}$}{2}
  \z{$81^{2/2}$}{81}
  \z{$100^{1/2}$}{10}
  \z{$9^{1/2}$}{3}
  \z{$169^{3/2}$}{2\,197}
  \z{$9^{3/2}$}{27}
  \z{$36^{1/2}$}{6}
  \btnc
  \z{$64^{1/2}$}{8}
  \z{$27^{1/3}$}{3}
  \z{$128^{1/7}$}{2}
  \z{$64^{3/3}$}{64}
  \z{$32^{1/5}$}{2}
  \z{$32^{3/5}$}{8}
  \z{$36^{3/2}$}{216}
  \z{$169^{1/2}$}{13}
  \z{$16^{1/2}$}{4}
  \z{$16^{2/4}$}{4}
  \z{$64^{2/3}$}{16}
  \z{$27^{3/3}$}{27}
  \z{$8^{3/3}$}{8}
\end{zestaw}

\btexpect{65}% ======================================================
%  GENERATED by tools/gen_sets.py -- do not edit.
%  Seed 2065; regenerate with `make sets`.
% ======================================================

\begin{zestaw}{Ułamek na procent II}{3}{6:40}{3}
  \z{$17/50$ na procent}{34\%}
  \z{$1/4$ na procent}{25\%}
  \z{$1/5$ na procent}{20\%}
  \z{$1/8$ na procent}{12,5\%}
  \z{$6/20$ na procent}{30\%}
  \z{$6/10$ na procent}{60\%}
  \z{$33/50$ na procent}{66\%}
  \z{$4/8$ na procent}{50\%}
  \z{$36/40$ na procent}{90\%}
  \z{$1/2$ na procent}{50\%}
  \z{$13/16$ na procent}{81,25\%}
  \z{$37/40$ na procent}{92,5\%}
  \z{$5/8$ na procent}{62,5\%}
  \z{$1/25$ na procent}{4\%}
  \btnc
  \z{$10/25$ na procent}{40\%}
  \z{$2/10$ na procent}{20\%}
  \z{$3/4$ na procent}{75\%}
  \z{$8/16$ na procent}{50\%}
  \z{$7/10$ na procent}{70\%}
  \z{$14/40$ na procent}{35\%}
  \z{$9/16$ na procent}{56,25\%}
  \z{$5/25$ na procent}{20\%}
  \z{$12/16$ na procent}{75\%}
  \z{$25/40$ na procent}{62,5\%}
  \z{$16/50$ na procent}{32\%}
  \z{$3/10$ na procent}{30\%}
  \z{$10/20$ na procent}{50\%}
  \btnc
  \z{$7/50$ na procent}{14\%}
  \z{$9/10$ na procent}{90\%}
  \z{$5/10$ na procent}{50\%}
  \z{$30/40$ na procent}{75\%}
  \z{$2/4$ na procent}{50\%}
  \z{$37/50$ na procent}{74\%}
  \z{$4/10$ na procent}{40\%}
  \z{$12/20$ na procent}{60\%}
  \z{$19/25$ na procent}{76\%}
  \z{$18/20$ na procent}{90\%}
  \z{$3/16$ na procent}{18,75\%}
  \z{$43/50$ na procent}{86\%}
  \z{$29/50$ na procent}{58\%}
\end{zestaw}

\btexpect{66}% ======================================================
%  GENERATED by tools/gen_sets.py -- do not edit.
%  Seed 2066; regenerate with `make sets`.
% ======================================================

\begin{zestaw}{Małe dziesiętne II}{3}{6:00}{3}
  \z{$0{,}003 \times 8\,000$}{24}
  \z{$0{,}08 \times 1\,500$}{120}
  \z{$0{,}009 \times 20\,000$}{180}
  \z{$0{,}002 \times 9\,000$}{18}
  \z{$0{,}007 \times 9\,000$}{63}
  \z{$0{,}005 \times 17\,000$}{85}
  \z{$0{,}05 \times 1\,700$}{85}
  \z{$0{,}06 \times 500$}{30}
  \z{$0{,}008 \times 14\,000$}{112}
  \z{$0{,}05 \times 700$}{35}
  \z{$0{,}008 \times 15\,000$}{120}
  \z{$0{,}002 \times 2\,000$}{4}
  \z{$0{,}006 \times 9\,000$}{54}
  \z{$0{,}04 \times 1\,600$}{64}
  \btnc
  \z{$0{,}03 \times 1\,500$}{45}
  \z{$0{,}008 \times 12\,000$}{96}
  \z{$0{,}02 \times 200$}{4}
  \z{$0{,}005 \times 19\,000$}{95}
  \z{$0{,}009 \times 17\,000$}{153}
  \z{$0{,}002 \times 8\,000$}{16}
  \z{$0{,}03 \times 300$}{9}
  \z{$0{,}006 \times 6\,000$}{36}
  \z{$0{,}07 \times 1\,500$}{105}
  \z{$0{,}09 \times 1\,000$}{90}
  \z{$0{,}09 \times 1\,900$}{171}
  \z{$0{,}07 \times 400$}{28}
  \z{$0{,}008 \times 7\,000$}{56}
  \btnc
  \z{$0{,}009 \times 10\,000$}{90}
  \z{$0{,}09 \times 900$}{81}
  \z{$0{,}07 \times 1\,300$}{91}
  \z{$0{,}002 \times 20\,000$}{40}
  \z{$0{,}04 \times 1\,300$}{52}
  \z{$0{,}005 \times 2\,000$}{10}
  \z{$0{,}009 \times 5\,000$}{45}
  \z{$0{,}04 \times 1\,200$}{48}
  \z{$0{,}005 \times 12\,000$}{60}
  \z{$0{,}004 \times 6\,000$}{24}
  \z{$0{,}009 \times 14\,000$}{126}
  \z{$0{,}03 \times 2\,000$}{60}
  \z{$0{,}003 \times 13\,000$}{39}
\end{zestaw}

\btexpect{67}% ======================================================
%  GENERATED by tools/gen_sets.py -- do not edit.
%  Seed 2067; regenerate with `make sets`.
% ======================================================

\begin{zestaw}{Mieszane — ułamki i potęgi}{3}{7:20}{3}
  \z{$1/3 \times 10/12$}{5/18}
  \z{$1/2 \times 9/12$}{3/8}
  \z{$4/5 \times 3/12$}{1/5}
  \z{$4/5 \times 1/3$}{4/15}
  \z{$1/3 \times 1/8$}{1/24}
  \z{$6/10 \times 1/10$}{3/50}
  \z{$1/5 \times 1/2$}{1/10}
  \z{$1/2 \times 5/12$}{5/24}
  \z{$1/9 \times 4/8$}{1/18}
  \z{$1/2 \times 1/2$}{1/4}
  \z{$5^{-2}$}{1/25}
  \z{$2^{-1}$}{1/2}
  \z{$10^{-4}$}{1/10000}
  \z{$3^{-3}$}{1/27}
  \btnc
  \z{$5^{-1}$}{1/5}
  \z{$6^{-3}$}{1/216}
  \z{$2^{-2}$}{1/4}
  \z{$6^{0}$}{1}
  \z{$5^{0}$}{1}
  \z{$8^{0}$}{1}
  \z{$\log_{3} 27$}{3}
  \z{$\log_{6} 46\,656$}{6}
  \z{$\log_{4} 65\,536$}{8}
  \z{$\log_{5} 78\,125$}{7}
  \z{$\log_{8} 32\,768$}{5}
  \z{$\log_{3} 9$}{2}
  \z{$\log_{2} 2\,048$}{11}
  \btnc
  \z{$\log_{10} 100\,000$}{5}
  \z{$\log_{7} 16\,807$}{5}
  \z{$\log_{4} 4\,096$}{6}
  \z{$47/50$ dziesiętnie}{0,94}
  \z{$3/10$ dziesiętnie}{0,3}
  \z{$5/8$ dziesiętnie}{0,625}
  \z{$1/4$ dziesiętnie}{0,25}
  \z{$7/10$ dziesiętnie}{0,7}
  \z{$4/8$ dziesiętnie}{0,5}
  \z{$30/40$ dziesiętnie}{0,75}
  \z{$1/2$ dziesiętnie}{0,5}
  \z{$5/10$ dziesiętnie}{0,5}
  \z{$46/50$ dziesiętnie}{0,92}
\end{zestaw}

\btexpect{68}% ======================================================
%  GENERATED by tools/gen_sets.py -- do not edit.
%  Seed 2068; regenerate with `make sets`.
% ======================================================

\begin{zestaw}{Mnożenie ułamków III}{3}{6:40}{3}
  \z{$2/3 \times 2/10$}{2/15}
  \z{$1/10 \times 2/9$}{1/45}
  \z{$2/5 \times 1/5$}{2/25}
  \z{$8/9 \times 2/12$}{4/27}
  \z{$1/2 \times 3/8$}{3/16}
  \z{$3/4 \times 3/5$}{9/20}
  \z{$6/12 \times 1/2$}{1/4}
  \z{$3/6 \times 6/10$}{3/10}
  \z{$3/6 \times 2/6$}{1/6}
  \z{$5/6 \times 1/6$}{5/36}
  \z{$6/8 \times 2/5$}{3/10}
  \z{$1/2 \times 2/5$}{1/5}
  \z{$1/10 \times 3/4$}{3/40}
  \z{$3/8 \times 3/9$}{1/8}
  \btnc
  \z{$1/5 \times 6/12$}{1/10}
  \z{$2/4 \times 11/12$}{11/24}
  \z{$3/8 \times 1/12$}{1/32}
  \z{$1/8 \times 11/12$}{11/96}
  \z{$1/5 \times 2/4$}{1/10}
  \z{$7/12 \times 1/3$}{7/36}
  \z{$3/6 \times 10/12$}{5/12}
  \z{$6/10 \times 1/3$}{1/5}
  \z{$1/2 \times 3/4$}{3/8}
  \z{$2/8 \times 2/3$}{1/6}
  \z{$1/3 \times 1/3$}{1/9}
  \z{$4/5 \times 3/6$}{2/5}
  \z{$1/9 \times 2/3$}{2/27}
  \btnc
  \z{$3/8 \times 1/8$}{3/64}
  \z{$8/9 \times 1/5$}{8/45}
  \z{$5/6 \times 7/10$}{7/12}
  \z{$2/5 \times 4/9$}{8/45}
  \z{$1/8 \times 2/9$}{1/36}
  \z{$7/10 \times 1/2$}{7/20}
  \z{$2/3 \times 4/5$}{8/15}
  \z{$5/10 \times 2/6$}{1/6}
  \z{$2/10 \times 1/4$}{1/20}
  \z{$2/10 \times 2/3$}{2/15}
  \z{$1/2 \times 1/2$}{1/4}
  \z{$1/3 \times 11/12$}{11/36}
  \z{$4/8 \times 5/6$}{5/12}
\end{zestaw}

\btexpect{69}% ======================================================
%  GENERATED by tools/gen_sets.py -- do not edit.
%  Seed 2069; regenerate with `make sets`.
% ======================================================

\begin{zestaw}{Logarytmy III}{3}{6:00}{3}
  \z{$\log_{7} 16\,807$}{5}
  \z{$\log_{6} 46\,656$}{6}
  \z{$\log_{3} 59\,049$}{10}
  \z{$\log_{5} 78\,125$}{7}
  \z{$\log_{10} 100\,000$}{5}
  \z{$\log_{4} 4\,096$}{6}
  \z{$\log_{3} 9$}{2}
  \z{$\log_{4} 65\,536$}{8}
  \z{$\log_{2} 4\,096$}{12}
  \z{$\log_{8} 32\,768$}{5}
  \z{$\log_{9} 59\,049$}{5}
  \z{$\log_{8} 64$}{2}
  \z{$\log_{3} 19\,683$}{9}
  \z{$\log_{7} 343$}{3}
  \btnc
  \z{$\log_{6} 7\,776$}{5}
  \z{$\log_{3} 27$}{3}
  \z{$\log_{5} 625$}{4}
  \z{$\log_{3} 6\,561$}{8}
  \z{$\log_{5} 125$}{3}
  \z{$\log_{6} 36$}{2}
  \z{$\log_{2} 16\,384$}{14}
  \z{$\log_{2} 128$}{7}
  \z{$\log_{4} 256$}{4}
  \z{$\log_{2} 16$}{4}
  \z{$\log_{4} 16$}{2}
  \z{$\log_{3} 729$}{6}
  \z{$\log_{2} 4$}{2}
  \btnc
  \z{$\log_{3} 2\,187$}{7}
  \z{$\log_{4} 16\,384$}{7}
  \z{$\log_{5} 15\,625$}{6}
  \z{$\log_{8} 4\,096$}{4}
  \z{$\log_{6} 1\,296$}{4}
  \z{$\log_{10} 100$}{2}
  \z{$\log_{10} 10\,000$}{4}
  \z{$\log_{2} 32$}{5}
  \z{$\log_{2} 8$}{3}
  \z{$\log_{2} 8\,192$}{13}
  \z{$\log_{7} 49$}{2}
  \z{$\log_{2} 32\,768$}{15}
  \z{$\log_{2} 1\,024$}{10}
\end{zestaw}

\btexpect{70}% ======================================================
%  GENERATED by tools/gen_sets.py -- do not edit.
%  Seed 2070; regenerate with `make sets`.
% ======================================================

\begin{zestaw}{Potęgi ujemne i zero III}{3}{5:20}{3}
  \z{$8^{-1}$}{1/8}
  \z{$3^{-1}$}{1/3}
  \z{$10^{-1}$}{1/10}
  \z{$2^{-1}$}{1/2}
  \z{$9^{0}$}{1}
  \z{$2^{0}$}{1}
  \z{$7^{-4}$}{1/2401}
  \z{$5^{-4}$}{1/625}
  \z{$6^{-3}$}{1/216}
  \z{$8^{-3}$}{1/512}
  \z{$8^{0}$}{1}
  \z{$2^{-3}$}{1/8}
  \z{$8^{-4}$}{1/4096}
  \z{$7^{-3}$}{1/343}
  \btnc
  \z{$4^{-2}$}{1/16}
  \z{$3^{0}$}{1}
  \z{$10^{-4}$}{1/10000}
  \z{$4^{-4}$}{1/256}
  \z{$7^{-1}$}{1/7}
  \z{$2^{-4}$}{1/16}
  \z{$7^{0}$}{1}
  \z{$5^{-2}$}{1/25}
  \z{$3^{-2}$}{1/9}
  \z{$6^{0}$}{1}
  \z{$10^{0}$}{1}
  \z{$8^{-2}$}{1/64}
  \z{$4^{-1}$}{1/4}
  \btnc
  \z{$10^{-2}$}{1/100}
  \z{$3^{-3}$}{1/27}
  \z{$7^{-2}$}{1/49}
  \z{$4^{0}$}{1}
  \z{$9^{-3}$}{1/729}
  \z{$3^{-4}$}{1/81}
  \z{$9^{-2}$}{1/81}
  \z{$6^{-2}$}{1/36}
  \z{$6^{-4}$}{1/1296}
  \z{$9^{-1}$}{1/9}
  \z{$4^{-3}$}{1/64}
  \z{$9^{-4}$}{1/6561}
  \z{$5^{0}$}{1}
\end{zestaw}

\btexpect{71}% ======================================================
%  GENERATED by tools/gen_sets.py -- do not edit.
%  Seed 2071; regenerate with `make sets`.
% ======================================================

\begin{zestaw}{Dzielenie ułamków III}{3}{7:20}{3}
  \z{$1/2 \div 1/4$}{2}
  \z{$2/3 \div 11/12$}{8/11}
  \z{$7/8 \div 1/3$}{21/8}
  \z{$3/10 \div 4/6$}{9/20}
  \z{$6/8 \div 3/5$}{5/4}
  \z{$3/6 \div 2/8$}{2}
  \z{$1/2 \div 1/2$}{1}
  \z{$1/2 \div 4/6$}{3/4}
  \z{$1/3 \div 3/4$}{4/9}
  \z{$1/10 \div 4/12$}{3/10}
  \z{$4/6 \div 1/2$}{4/3}
  \z{$7/9 \div 7/12$}{4/3}
  \z{$1/8 \div 1/12$}{3/2}
  \z{$3/6 \div 2/5$}{5/4}
  \btnc
  \z{$6/8 \div 7/12$}{9/7}
  \z{$2/4 \div 1/5$}{5/2}
  \z{$7/8 \div 4/6$}{21/16}
  \z{$5/12 \div 2/9$}{15/8}
  \z{$1/3 \div 2/3$}{1/2}
  \z{$3/9 \div 7/8$}{8/21}
  \z{$8/12 \div 1/2$}{4/3}
  \z{$3/5 \div 1/2$}{6/5}
  \z{$1/4 \div 6/10$}{5/12}
  \z{$4/9 \div 8/9$}{1/2}
  \z{$7/12 \div 5/8$}{14/15}
  \z{$4/8 \div 3/4$}{2/3}
  \z{$2/4 \div 1/2$}{1}
  \btnc
  \z{$4/6 \div 5/6$}{4/5}
  \z{$4/9 \div 2/12$}{8/3}
  \z{$5/8 \div 1/2$}{5/4}
  \z{$2/10 \div 3/12$}{4/5}
  \z{$1/2 \div 3/5$}{5/6}
  \z{$6/10 \div 1/2$}{6/5}
  \z{$3/9 \div 5/6$}{2/5}
  \z{$2/6 \div 3/12$}{4/3}
  \z{$2/9 \div 8/12$}{1/3}
  \z{$2/8 \div 1/6$}{3/2}
  \z{$3/5 \div 6/9$}{9/10}
  \z{$6/12 \div 11/12$}{6/11}
  \z{$2/12 \div 1/3$}{1/2}
\end{zestaw}

\btexpect{72}% ======================================================
%  GENERATED by tools/gen_sets.py -- do not edit.
%  Seed 2072; regenerate with `make sets`.
% ======================================================

\begin{zestaw}{Wykładnik ułamkowy III}{3}{7:20}{3}
  \z{$64^{1/3}$}{4}
  \z{$64^{3/6}$}{8}
  \z{$169^{2/2}$}{169}
  \z{$36^{3/2}$}{216}
  \z{$121^{3/2}$}{1\,331}
  \z{$64^{2/6}$}{4}
  \z{$125^{2/3}$}{25}
  \z{$81^{2/2}$}{81}
  \z{$27^{1/3}$}{3}
  \z{$64^{3/2}$}{512}
  \z{$64^{2/2}$}{64}
  \z{$64^{2/3}$}{16}
  \z{$128^{3/7}$}{8}
  \z{$64^{1/6}$}{2}
  \btnc
  \z{$64^{3/3}$}{64}
  \z{$16^{2/2}$}{16}
  \z{$144^{1/2}$}{12}
  \z{$32^{2/5}$}{4}
  \z{$49^{1/2}$}{7}
  \z{$8^{2/3}$}{4}
  \z{$27^{3/3}$}{27}
  \z{$9^{1/2}$}{3}
  \z{$25^{1/2}$}{5}
  \z{$25^{3/2}$}{125}
  \z{$100^{1/2}$}{10}
  \z{$144^{2/2}$}{144}
  \z{$64^{1/2}$}{8}
  \btnc
  \z{$25^{2/2}$}{25}
  \z{$4^{1/2}$}{2}
  \z{$49^{2/2}$}{49}
  \z{$8^{1/3}$}{2}
  \z{$125^{3/3}$}{125}
  \z{$16^{3/4}$}{8}
  \z{$121^{1/2}$}{11}
  \z{$81^{1/2}$}{9}
  \z{$9^{3/2}$}{27}
  \z{$36^{2/2}$}{36}
  \z{$144^{3/2}$}{1\,728}
  \z{$9^{2/2}$}{9}
  \z{$16^{1/2}$}{4}
\end{zestaw}



\chapter{Blok IV — Zastosowania}

Rachunki, w których intuicyjna odpowiedź jest regularnie zła. Średnia prędkość
tam i z powrotem nie jest średnią z dwóch prędkości. Procent przez trzy lata
nie jest procentem razy trzy. Do tego kombinatoryka, skala mapy i dzień
tygodnia z daty — jedyne zadanie w obu tomach, które wygląda na niemożliwe do
zrobienia w pamięci i nie jest.

% GENERATED by tools/gen_sets.py -- do not edit.
\btexpect{73}% ======================================================
%  GENERATED by tools/gen_sets.py -- do not edit.
%  Seed 2073; regenerate with `make sets`.
% ======================================================

\begin{zestaw}{Kombinatoryka}{3}{8:00}{3}
  \z{$3!$}{6}
  \z{$4!$}{24}
  \z{$C(10, 4)$}{210}
  \z{$10!$}{3\,628\,800}
  \z{$6!$}{720}
  \z{$C(4, 3)$}{4}
  \z{$C(9, 4)$}{126}
  \z{$C(6, 4)$}{15}
  \z{$C(7, 2)$}{21}
  \z{$P(12, 2)$}{132}
  \z{$P(10, 3)$}{720}
  \z{$C(8, 5)$}{56}
  \z{$P(5, 4)$}{120}
  \z{$7!$}{5\,040}
  \btnc
  \z{$5!$}{120}
  \z{$C(13, 3)$}{286}
  \z{$C(13, 5)$}{1\,287}
  \z{$C(13, 2)$}{78}
  \z{$C(8, 4)$}{70}
  \z{$P(6, 3)$}{120}
  \z{$C(9, 3)$}{84}
  \z{$P(12, 3)$}{1\,320}
  \z{$P(4, 2)$}{12}
  \z{$C(11, 2)$}{55}
  \z{$P(8, 4)$}{1\,680}
  \z{$P(11, 4)$}{7\,920}
  \z{$9!$}{362\,880}
  \btnc
  \z{$C(5, 3)$}{10}
  \z{$P(5, 3)$}{60}
  \z{$P(13, 4)$}{17\,160}
  \z{$P(11, 2)$}{110}
  \z{$P(10, 2)$}{90}
  \z{$P(7, 4)$}{840}
  \z{$P(4, 3)$}{24}
  \z{$C(12, 3)$}{220}
  \z{$C(12, 4)$}{495}
  \z{$C(12, 2)$}{66}
  \z{$P(8, 2)$}{56}
  \z{$P(13, 3)$}{1\,716}
  \z{$P(6, 2)$}{30}
\end{zestaw}

\btexpect{74}% ======================================================
%  GENERATED by tools/gen_sets.py -- do not edit.
%  Seed 2074; regenerate with `make sets`.
% ======================================================

\begin{zestaw}{km/h i m/s}{3}{6:00}{2}
  \z{$198$ km/h $\rightarrow$ m/s}{55}
  \z{$162$ km/h $\rightarrow$ m/s}{45}
  \z{$57,5$ m/s $\rightarrow$ km/h}{207}
  \z{$32,5$ m/s $\rightarrow$ km/h}{117}
  \z{$54$ km/h $\rightarrow$ m/s}{15}
  \z{$18$ km/h $\rightarrow$ m/s}{5}
  \z{$36$ km/h $\rightarrow$ m/s}{10}
  \z{$90$ km/h $\rightarrow$ m/s}{25}
  \z{$171$ km/h $\rightarrow$ m/s}{47,5}
  \z{$12,5$ m/s $\rightarrow$ km/h}{45}
  \z{$135$ km/h $\rightarrow$ m/s}{37,5}
  \z{$45$ m/s $\rightarrow$ km/h}{162}
  \z{$42,5$ m/s $\rightarrow$ km/h}{153}
  \z{$17,5$ m/s $\rightarrow$ km/h}{63}
  \z{$72$ km/h $\rightarrow$ m/s}{20}
  \z{$20$ m/s $\rightarrow$ km/h}{72}
  \z{$108$ km/h $\rightarrow$ m/s}{30}
  \z{$30$ m/s $\rightarrow$ km/h}{108}
  \z{$189$ km/h $\rightarrow$ m/s}{52,5}
  \z{$47,5$ m/s $\rightarrow$ km/h}{171}
  \btnc
  \z{$27,5$ m/s $\rightarrow$ km/h}{99}
  \z{$10$ m/s $\rightarrow$ km/h}{36}
  \z{$25$ m/s $\rightarrow$ km/h}{90}
  \z{$37,5$ m/s $\rightarrow$ km/h}{135}
  \z{$207$ km/h $\rightarrow$ m/s}{57,5}
  \z{$55$ m/s $\rightarrow$ km/h}{198}
  \z{$117$ km/h $\rightarrow$ m/s}{32,5}
  \z{$22,5$ m/s $\rightarrow$ km/h}{81}
  \z{$52,5$ m/s $\rightarrow$ km/h}{189}
  \z{$50$ m/s $\rightarrow$ km/h}{180}
  \z{$144$ km/h $\rightarrow$ m/s}{40}
  \z{$63$ km/h $\rightarrow$ m/s}{17,5}
  \z{$216$ km/h $\rightarrow$ m/s}{60}
  \z{$180$ km/h $\rightarrow$ m/s}{50}
  \z{$40$ m/s $\rightarrow$ km/h}{144}
  \z{$35$ m/s $\rightarrow$ km/h}{126}
  \z{$15$ m/s $\rightarrow$ km/h}{54}
  \z{$81$ km/h $\rightarrow$ m/s}{22,5}
  \z{$27$ km/h $\rightarrow$ m/s}{7,5}
  \z{$45$ km/h $\rightarrow$ m/s}{12,5}
\end{zestaw}

\btexpect{75}% ======================================================
%  GENERATED by tools/gen_sets.py -- do not edit.
%  Seed 2075; regenerate with `make sets`.
% ======================================================

\begin{zestaw}{Prędkość średnia}{3}{8:40}{2}
  \zz{$240$ km w $2$ h --- km/h}{120}
  \zz{$200$ km w $2$ h $30$ min --- km/h}{80}
  \zz{$180$ km w $1$ h $30$ min --- km/h}{120}
  \zz{$60$ km w $1$ h $30$ min --- km/h}{40}
  \zz{$75$ km w $1$ h $30$ min --- km/h}{50}
  \zz{$270$ km w $2$ h $15$ min --- km/h}{120}
  \zz{$175$ km w $1$ h $45$ min --- km/h}{100}
  \zz{$50$ km w $1$ h $15$ min --- km/h}{40}
  \zz{$70$ km w $1$ h $45$ min --- km/h}{40}
  \zz{$80$ km w $2$ h --- km/h}{40}
  \zz{$125$ km w $1$ h $15$ min --- km/h}{100}
  \zz{$180$ km w $2$ h --- km/h}{90}
  \zz{$225$ km w $2$ h $15$ min --- km/h}{100}
  \zz{$75$ km w $45$ min --- km/h}{100}
  \zz{$60$ km w $45$ min --- km/h}{80}
  \zz{$45$ km w $30$ min --- km/h}{90}
  \zz{$105$ km w $1$ h $45$ min --- km/h}{60}
  \zz{$90$ km w $2$ h --- km/h}{45}
  \zz{$45$ km w $45$ min --- km/h}{60}
  \zz{$105$ km w $1$ h $30$ min --- km/h}{70}
  \btnc
  \zz{$30$ km w $45$ min --- km/h}{40}
  \zz{$150$ km w $1$ h $15$ min --- km/h}{120}
  \zz{$120$ km w $1$ h $30$ min --- km/h}{80}
  \zz{$120$ km w $2$ h --- km/h}{60}
  \zz{$225$ km w $2$ h $30$ min --- km/h}{90}
  \zz{$125$ km w $2$ h $30$ min --- km/h}{50}
  \zz{$40$ km w $30$ min --- km/h}{80}
  \zz{$140$ km w $2$ h --- km/h}{70}
  \zz{$300$ km w $2$ h $30$ min --- km/h}{120}
  \zz{$175$ km w $2$ h $30$ min --- km/h}{70}
  \zz{$180$ km w $2$ h $15$ min --- km/h}{80}
  \zz{$100$ km w $2$ h --- km/h}{50}
  \zz{$90$ km w $2$ h $15$ min --- km/h}{40}
  \zz{$150$ km w $2$ h $30$ min --- km/h}{60}
  \zz{$60$ km w $30$ min --- km/h}{120}
  \zz{$210$ km w $1$ h $45$ min --- km/h}{120}
  \zz{$200$ km w $2$ h --- km/h}{100}
  \zz{$35$ km w $30$ min --- km/h}{70}
  \zz{$90$ km w $1$ h $30$ min --- km/h}{60}
  \zz{$90$ km w $45$ min --- km/h}{120}
\end{zestaw}

\btexpect{76}% ======================================================
%  GENERATED by tools/gen_sets.py -- do not edit.
%  Seed 2076; regenerate with `make sets`.
% ======================================================

\begin{zestaw}{Skala mapy}{3}{8:00}{2}
  \zz{skala $1:25\,000$, $10$ cm --- ile km}{2,5}
  \zz{skala $1:10\,000$, $11$ cm --- ile km}{1,1}
  \zz{skala $1:50\,000$, $7$ cm --- ile km}{3,5}
  \zz{skala $1:200\,000$, $10$ cm --- ile km}{20}
  \zz{skala $1:200\,000$, $4$ cm --- ile km}{8}
  \zz{skala $1:100\,000$, $12$ cm --- ile km}{12}
  \zz{skala $1:10\,000$, $2$ cm --- ile km}{0,2}
  \zz{skala $1:10\,000$, $12$ cm --- ile km}{1,2}
  \zz{skala $1:100\,000$, $4$ cm --- ile km}{4}
  \zz{skala $1:200\,000$, $6$ cm --- ile km}{12}
  \zz{skala $1:200\,000$, $2$ cm --- ile km}{4}
  \zz{skala $1:50\,000$, $11$ cm --- ile km}{5,5}
  \zz{skala $1:50\,000$, $10$ cm --- ile km}{5}
  \zz{skala $1:25\,000$, $9$ cm --- ile km}{2,25}
  \zz{skala $1:50\,000$, $5$ cm --- ile km}{2,5}
  \zz{skala $1:200\,000$, $3$ cm --- ile km}{6}
  \zz{skala $1:25\,000$, $6$ cm --- ile km}{1,5}
  \zz{skala $1:200\,000$, $8$ cm --- ile km}{16}
  \zz{skala $1:200\,000$, $7$ cm --- ile km}{14}
  \zz{skala $1:25\,000$, $4$ cm --- ile km}{1}
  \btnc
  \zz{skala $1:200\,000$, $12$ cm --- ile km}{24}
  \zz{skala $1:200\,000$, $9$ cm --- ile km}{18}
  \zz{skala $1:50\,000$, $12$ cm --- ile km}{6}
  \zz{skala $1:25\,000$, $7$ cm --- ile km}{1,75}
  \zz{skala $1:10\,000$, $4$ cm --- ile km}{0,4}
  \zz{skala $1:50\,000$, $4$ cm --- ile km}{2}
  \zz{skala $1:50\,000$, $3$ cm --- ile km}{1,5}
  \zz{skala $1:10\,000$, $3$ cm --- ile km}{0,3}
  \zz{skala $1:10\,000$, $7$ cm --- ile km}{0,7}
  \zz{skala $1:50\,000$, $8$ cm --- ile km}{4}
  \zz{skala $1:200\,000$, $11$ cm --- ile km}{22}
  \zz{skala $1:100\,000$, $7$ cm --- ile km}{7}
  \zz{skala $1:200\,000$, $5$ cm --- ile km}{10}
  \zz{skala $1:100\,000$, $9$ cm --- ile km}{9}
  \zz{skala $1:25\,000$, $3$ cm --- ile km}{0,75}
  \zz{skala $1:10\,000$, $6$ cm --- ile km}{0,6}
  \zz{skala $1:100\,000$, $3$ cm --- ile km}{3}
  \zz{skala $1:100\,000$, $5$ cm --- ile km}{5}
  \zz{skala $1:100\,000$, $11$ cm --- ile km}{11}
  \zz{skala $1:100\,000$, $6$ cm --- ile km}{6}
\end{zestaw}

\btexpect{77}% ======================================================
%  GENERATED by tools/gen_sets.py -- do not edit.
%  Seed 2077; regenerate with `make sets`.
% ======================================================

\begin{zestaw}{Tam i z powrotem}{3}{9:20}{2}
  \zz{tam $30$ km/h, z powrotem $20$ km/h --- średnia}{24}
  \zz{tam $72$ km/h, z powrotem $120$ km/h --- średnia}{90}
  \zz{tam $120$ km/h, z powrotem $240$ km/h --- średnia}{160}
  \zz{tam $72$ km/h, z powrotem $90$ km/h --- średnia}{80}
  \zz{tam $60$ km/h, z powrotem $40$ km/h --- średnia}{48}
  \zz{tam $60$ km/h, z powrotem $30$ km/h --- średnia}{40}
  \zz{tam $90$ km/h, z powrotem $60$ km/h --- średnia}{72}
  \zz{tam $240$ km/h, z powrotem $60$ km/h --- średnia}{96}
  \zz{tam $24$ km/h, z powrotem $120$ km/h --- średnia}{40}
  \zz{tam $72$ km/h, z powrotem $36$ km/h --- średnia}{48}
  \zz{tam $36$ km/h, z powrotem $45$ km/h --- średnia}{40}
  \zz{tam $120$ km/h, z powrotem $24$ km/h --- średnia}{40}
  \zz{tam $20$ km/h, z powrotem $80$ km/h --- średnia}{32}
  \zz{tam $120$ km/h, z powrotem $72$ km/h --- średnia}{90}
  \zz{tam $80$ km/h, z powrotem $48$ km/h --- średnia}{60}
  \zz{tam $240$ km/h, z powrotem $80$ km/h --- średnia}{120}
  \zz{tam $60$ km/h, z powrotem $20$ km/h --- średnia}{30}
  \zz{tam $50$ km/h, z powrotem $200$ km/h --- średnia}{80}
  \zz{tam $24$ km/h, z powrotem $40$ km/h --- średnia}{30}
  \zz{tam $48$ km/h, z powrotem $144$ km/h --- średnia}{72}
  \btnc
  \zz{tam $180$ km/h, z powrotem $36$ km/h --- średnia}{60}
  \zz{tam $144$ km/h, z powrotem $48$ km/h --- średnia}{72}
  \zz{tam $50$ km/h, z powrotem $75$ km/h --- średnia}{60}
  \zz{tam $80$ km/h, z powrotem $120$ km/h --- średnia}{96}
  \zz{tam $60$ km/h, z powrotem $120$ km/h --- średnia}{80}
  \zz{tam $120$ km/h, z powrotem $30$ km/h --- średnia}{48}
  \zz{tam $90$ km/h, z powrotem $180$ km/h --- średnia}{120}
  \zz{tam $45$ km/h, z powrotem $180$ km/h --- średnia}{72}
  \zz{tam $60$ km/h, z powrotem $100$ km/h --- średnia}{75}
  \zz{tam $20$ km/h, z powrotem $60$ km/h --- średnia}{30}
  \zz{tam $40$ km/h, z powrotem $120$ km/h --- średnia}{60}
  \zz{tam $60$ km/h, z powrotem $240$ km/h --- średnia}{96}
  \zz{tam $36$ km/h, z powrotem $72$ km/h --- średnia}{48}
  \zz{tam $200$ km/h, z powrotem $50$ km/h --- średnia}{80}
  \zz{tam $36$ km/h, z powrotem $60$ km/h --- średnia}{45}
  \zz{tam $180$ km/h, z powrotem $144$ km/h --- średnia}{160}
  \zz{tam $144$ km/h, z powrotem $240$ km/h --- średnia}{180}
  \zz{tam $30$ km/h, z powrotem $60$ km/h --- średnia}{40}
  \zz{tam $50$ km/h, z powrotem $150$ km/h --- średnia}{75}
  \zz{tam $60$ km/h, z powrotem $36$ km/h --- średnia}{45}
\end{zestaw}

\btexpect{78}% ======================================================
%  GENERATED by tools/gen_sets.py -- do not edit.
%  Seed 2078; regenerate with `make sets`.
% ======================================================

\begin{zestaw}{Dzień tygodnia}{3}{10:40}{2}
  \zz{2.04.2002 --- jaki dzień tygodnia}{wtorek}
  \zz{22.06.2030 --- jaki dzień tygodnia}{sobota}
  \zz{27.12.2046 --- jaki dzień tygodnia}{czwartek}
  \zz{26.10.2023 --- jaki dzień tygodnia}{czwartek}
  \zz{25.06.2043 --- jaki dzień tygodnia}{czwartek}
  \zz{12.10.2005 --- jaki dzień tygodnia}{środa}
  \zz{5.06.2034 --- jaki dzień tygodnia}{poniedziałek}
  \zz{16.11.2003 --- jaki dzień tygodnia}{niedziela}
  \zz{12.03.2036 --- jaki dzień tygodnia}{środa}
  \zz{12.05.2021 --- jaki dzień tygodnia}{środa}
  \zz{13.03.2011 --- jaki dzień tygodnia}{niedziela}
  \zz{16.06.2020 --- jaki dzień tygodnia}{wtorek}
  \zz{20.02.2012 --- jaki dzień tygodnia}{poniedziałek}
  \zz{9.01.2020 --- jaki dzień tygodnia}{czwartek}
  \zz{12.04.2042 --- jaki dzień tygodnia}{sobota}
  \zz{17.02.2015 --- jaki dzień tygodnia}{wtorek}
  \zz{24.06.2025 --- jaki dzień tygodnia}{wtorek}
  \zz{15.01.2032 --- jaki dzień tygodnia}{czwartek}
  \zz{19.12.2015 --- jaki dzień tygodnia}{sobota}
  \zz{20.08.2036 --- jaki dzień tygodnia}{środa}
  \btnc
  \zz{8.11.2024 --- jaki dzień tygodnia}{piątek}
  \zz{19.01.2015 --- jaki dzień tygodnia}{poniedziałek}
  \zz{21.02.2041 --- jaki dzień tygodnia}{czwartek}
  \zz{22.07.2019 --- jaki dzień tygodnia}{poniedziałek}
  \zz{17.11.2036 --- jaki dzień tygodnia}{poniedziałek}
  \zz{4.11.2042 --- jaki dzień tygodnia}{wtorek}
  \zz{4.09.2039 --- jaki dzień tygodnia}{niedziela}
  \zz{21.06.2040 --- jaki dzień tygodnia}{czwartek}
  \zz{2.02.2034 --- jaki dzień tygodnia}{czwartek}
  \zz{24.01.2007 --- jaki dzień tygodnia}{środa}
  \zz{27.09.2040 --- jaki dzień tygodnia}{czwartek}
  \zz{25.09.2031 --- jaki dzień tygodnia}{czwartek}
  \zz{8.07.2036 --- jaki dzień tygodnia}{wtorek}
  \zz{15.04.2042 --- jaki dzień tygodnia}{wtorek}
  \zz{24.10.2036 --- jaki dzień tygodnia}{piątek}
  \zz{8.02.2016 --- jaki dzień tygodnia}{poniedziałek}
  \zz{9.04.2048 --- jaki dzień tygodnia}{czwartek}
  \zz{4.01.2037 --- jaki dzień tygodnia}{niedziela}
  \zz{16.07.2002 --- jaki dzień tygodnia}{wtorek}
  \zz{24.09.2032 --- jaki dzień tygodnia}{piątek}
\end{zestaw}

\btexpect{79}% ======================================================
%  GENERATED by tools/gen_sets.py -- do not edit.
%  Seed 2079; regenerate with `make sets`.
% ======================================================

\begin{zestaw}{Procent przez lata}{3}{9:20}{2}
  \zz{$5\,000$ przy $20\%$ przez $2$ lata}{7\,200}
  \zz{$2\,000$ przy $50\%$ przez $2$ lata}{4\,500}
  \zz{$10\,000$ przy $20\%$ przez $2$ lata}{14\,400}
  \zz{$2\,000$ przy $10\%$ przez $3$ lata}{2\,662}
  \zz{$1\,000$ przy $50\%$ przez $2$ lata}{2\,250}
  \zz{$2\,000$ przy $20\%$ przez $2$ lata}{2\,880}
  \zz{$8\,000$ przy $50\%$ przez $3$ lata}{27\,000}
  \zz{$4\,000$ przy $50\%$ przez $2$ lata}{9\,000}
  \zz{$5\,000$ przy $50\%$ przez $2$ lata}{11\,250}
  \zz{$5\,000$ przy $25\%$ przez $2$ lata}{7\,812}
  \zz{$5\,000$ przy $50\%$ przez $3$ lata}{16\,875}
  \zz{$2\,000$ przy $20\%$ przez $3$ lata}{3\,456}
  \zz{$4\,000$ przy $25\%$ przez $3$ lata}{7\,812}
  \zz{$10\,000$ przy $25\%$ przez $3$ lata}{19\,531}
  \zz{$1\,000$ przy $10\%$ przez $2$ lata}{1\,210}
  \zz{$10\,000$ przy $20\%$ przez $3$ lata}{17\,280}
  \zz{$1\,000$ przy $20\%$ przez $2$ lata}{1\,440}
  \zz{$1\,000$ przy $10\%$ przez $3$ lata}{1\,331}
  \zz{$1\,000$ przy $50\%$ przez $3$ lata}{3\,375}
  \zz{$5\,000$ przy $25\%$ przez $3$ lata}{9\,765}
  \btnc
  \zz{$4\,000$ przy $20\%$ przez $3$ lata}{6\,912}
  \zz{$10\,000$ przy $25\%$ przez $2$ lata}{15\,625}
  \zz{$4\,000$ przy $25\%$ przez $2$ lata}{6\,250}
  \zz{$5\,000$ przy $10\%$ przez $3$ lata}{6\,655}
  \zz{$10\,000$ przy $10\%$ przez $2$ lata}{12\,100}
  \zz{$4\,000$ przy $10\%$ przez $2$ lata}{4\,840}
  \zz{$1\,000$ przy $25\%$ przez $3$ lata}{1\,952}
  \zz{$8\,000$ przy $20\%$ przez $2$ lata}{11\,520}
  \zz{$2\,000$ przy $10\%$ przez $2$ lata}{2\,420}
  \zz{$4\,000$ przy $50\%$ przez $3$ lata}{13\,500}
  \zz{$10\,000$ przy $50\%$ przez $2$ lata}{22\,500}
  \zz{$5\,000$ przy $20\%$ przez $3$ lata}{8\,640}
  \zz{$8\,000$ przy $10\%$ przez $2$ lata}{9\,680}
  \zz{$10\,000$ przy $50\%$ przez $3$ lata}{33\,750}
  \zz{$2\,000$ przy $25\%$ przez $3$ lata}{3\,906}
  \zz{$2\,000$ przy $50\%$ przez $3$ lata}{6\,750}
  \zz{$8\,000$ przy $20\%$ przez $3$ lata}{13\,824}
  \zz{$8\,000$ przy $50\%$ przez $2$ lata}{18\,000}
  \zz{$1\,000$ przy $20\%$ przez $3$ lata}{1\,728}
  \zz{$8\,000$ przy $25\%$ przez $3$ lata}{15\,625}
\end{zestaw}

\btexpect{80}% ======================================================
%  GENERATED by tools/gen_sets.py -- do not edit.
%  Seed 2080; regenerate with `make sets`.
% ======================================================

\begin{zestaw}{Kombinatoryka II}{3}{8:00}{3}
  \z{$9!$}{362\,880}
  \z{$C(13, 5)$}{1\,287}
  \z{$C(6, 5)$}{6}
  \z{$P(13, 5)$}{154\,440}
  \z{$P(14, 3)$}{2\,184}
  \z{$7!$}{5\,040}
  \z{$C(7, 4)$}{35}
  \z{$P(8, 3)$}{336}
  \z{$C(14, 2)$}{91}
  \z{$8!$}{40\,320}
  \z{$10!$}{3\,628\,800}
  \z{$P(13, 4)$}{17\,160}
  \z{$P(8, 2)$}{56}
  \z{$C(13, 3)$}{286}
  \btnc
  \z{$P(14, 4)$}{24\,024}
  \z{$C(8, 3)$}{56}
  \z{$C(6, 4)$}{15}
  \z{$6!$}{720}
  \z{$P(13, 2)$}{156}
  \z{$4!$}{24}
  \z{$C(7, 5)$}{21}
  \z{$5!$}{120}
  \z{$P(4, 2)$}{12}
  \z{$C(7, 2)$}{21}
  \z{$C(13, 4)$}{715}
  \z{$P(9, 3)$}{504}
  \z{$C(9, 3)$}{84}
  \btnc
  \z{$3!$}{6}
  \z{$C(14, 3)$}{364}
  \z{$P(4, 3)$}{24}
  \z{$C(13, 2)$}{78}
  \z{$C(8, 2)$}{28}
  \z{$P(7, 5)$}{2\,520}
  \z{$C(11, 2)$}{55}
  \z{$P(5, 4)$}{120}
  \z{$C(8, 4)$}{70}
  \z{$C(4, 2)$}{6}
  \z{$C(11, 5)$}{462}
  \z{$P(5, 3)$}{60}
  \z{$P(6, 3)$}{120}
\end{zestaw}

\btexpect{81}% ======================================================
%  GENERATED by tools/gen_sets.py -- do not edit.
%  Seed 2081; regenerate with `make sets`.
% ======================================================

\begin{zestaw}{km/h i m/s II}{3}{6:00}{2}
  \z{$153$ km/h $\rightarrow$ m/s}{42,5}
  \z{$216$ km/h $\rightarrow$ m/s}{60}
  \z{$144$ km/h $\rightarrow$ m/s}{40}
  \z{$180$ km/h $\rightarrow$ m/s}{50}
  \z{$55$ m/s $\rightarrow$ km/h}{198}
  \z{$72$ km/h $\rightarrow$ m/s}{20}
  \z{$32,5$ m/s $\rightarrow$ km/h}{117}
  \z{$99$ km/h $\rightarrow$ m/s}{27,5}
  \z{$27,5$ m/s $\rightarrow$ km/h}{99}
  \z{$45$ m/s $\rightarrow$ km/h}{162}
  \z{$126$ km/h $\rightarrow$ m/s}{35}
  \z{$54$ km/h $\rightarrow$ m/s}{15}
  \z{$40$ m/s $\rightarrow$ km/h}{144}
  \z{$42,5$ m/s $\rightarrow$ km/h}{153}
  \z{$5$ m/s $\rightarrow$ km/h}{18}
  \z{$35$ m/s $\rightarrow$ km/h}{126}
  \z{$81$ km/h $\rightarrow$ m/s}{22,5}
  \z{$22,5$ m/s $\rightarrow$ km/h}{81}
  \z{$198$ km/h $\rightarrow$ m/s}{55}
  \z{$162$ km/h $\rightarrow$ m/s}{45}
  \btnc
  \z{$37,5$ m/s $\rightarrow$ km/h}{135}
  \z{$27$ km/h $\rightarrow$ m/s}{7,5}
  \z{$207$ km/h $\rightarrow$ m/s}{57,5}
  \z{$7,5$ m/s $\rightarrow$ km/h}{27}
  \z{$30$ m/s $\rightarrow$ km/h}{108}
  \z{$189$ km/h $\rightarrow$ m/s}{52,5}
  \z{$63$ km/h $\rightarrow$ m/s}{17,5}
  \z{$171$ km/h $\rightarrow$ m/s}{47,5}
  \z{$20$ m/s $\rightarrow$ km/h}{72}
  \z{$12,5$ m/s $\rightarrow$ km/h}{45}
  \z{$60$ m/s $\rightarrow$ km/h}{216}
  \z{$117$ km/h $\rightarrow$ m/s}{32,5}
  \z{$52,5$ m/s $\rightarrow$ km/h}{189}
  \z{$36$ km/h $\rightarrow$ m/s}{10}
  \z{$15$ m/s $\rightarrow$ km/h}{54}
  \z{$108$ km/h $\rightarrow$ m/s}{30}
  \z{$90$ km/h $\rightarrow$ m/s}{25}
  \z{$18$ km/h $\rightarrow$ m/s}{5}
  \z{$57,5$ m/s $\rightarrow$ km/h}{207}
  \z{$50$ m/s $\rightarrow$ km/h}{180}
\end{zestaw}

\btexpect{82}% ======================================================
%  GENERATED by tools/gen_sets.py -- do not edit.
%  Seed 2082; regenerate with `make sets`.
% ======================================================

\begin{zestaw}{Prędkość średnia II}{3}{8:40}{2}
  \zz{$40$ km w $30$ min --- km/h}{80}
  \zz{$100$ km w $2$ h --- km/h}{50}
  \zz{$105$ km w $1$ h $30$ min --- km/h}{70}
  \zz{$200$ km w $2$ h $30$ min --- km/h}{80}
  \zz{$75$ km w $1$ h $30$ min --- km/h}{50}
  \zz{$175$ km w $2$ h $30$ min --- km/h}{70}
  \zz{$150$ km w $1$ h $15$ min --- km/h}{120}
  \zz{$150$ km w $1$ h $30$ min --- km/h}{100}
  \zz{$100$ km w $2$ h $30$ min --- km/h}{40}
  \zz{$240$ km w $2$ h --- km/h}{120}
  \zz{$90$ km w $1$ h $30$ min --- km/h}{60}
  \zz{$60$ km w $30$ min --- km/h}{120}
  \zz{$300$ km w $2$ h $30$ min --- km/h}{120}
  \zz{$90$ km w $2$ h $15$ min --- km/h}{40}
  \zz{$135$ km w $1$ h $30$ min --- km/h}{90}
  \zz{$25$ km w $30$ min --- km/h}{50}
  \zz{$45$ km w $30$ min --- km/h}{90}
  \zz{$140$ km w $2$ h --- km/h}{70}
  \zz{$125$ km w $1$ h $15$ min --- km/h}{100}
  \zz{$250$ km w $2$ h $30$ min --- km/h}{100}
  \btnc
  \zz{$135$ km w $2$ h $15$ min --- km/h}{60}
  \zz{$180$ km w $2$ h --- km/h}{90}
  \zz{$150$ km w $2$ h $30$ min --- km/h}{60}
  \zz{$80$ km w $2$ h --- km/h}{40}
  \zz{$20$ km w $30$ min --- km/h}{40}
  \zz{$90$ km w $2$ h --- km/h}{45}
  \zz{$60$ km w $45$ min --- km/h}{80}
  \zz{$100$ km w $1$ h $15$ min --- km/h}{80}
  \zz{$125$ km w $2$ h $30$ min --- km/h}{50}
  \zz{$120$ km w $2$ h --- km/h}{60}
  \zz{$180$ km w $1$ h $30$ min --- km/h}{120}
  \zz{$30$ km w $30$ min --- km/h}{60}
  \zz{$175$ km w $1$ h $45$ min --- km/h}{100}
  \zz{$270$ km w $2$ h $15$ min --- km/h}{120}
  \zz{$50$ km w $1$ h $15$ min --- km/h}{40}
  \zz{$180$ km w $2$ h $15$ min --- km/h}{80}
  \zz{$120$ km w $1$ h $30$ min --- km/h}{80}
  \zz{$30$ km w $45$ min --- km/h}{40}
  \zz{$160$ km w $2$ h --- km/h}{80}
  \zz{$35$ km w $30$ min --- km/h}{70}
\end{zestaw}

\btexpect{83}\input{sets-adv/generated/83-skala-b}
\btexpect{84}\input{sets-adv/generated/84-harmoniczna-b}
\btexpect{85}% ======================================================
%  GENERATED by tools/gen_sets.py -- do not edit.
%  Seed 2085; regenerate with `make sets`.
% ======================================================

\begin{zestaw}{Dzień tygodnia II}{3}{10:40}{2}
  \zz{24.03.2005 --- jaki dzień tygodnia}{czwartek}
  \zz{7.12.2014 --- jaki dzień tygodnia}{niedziela}
  \zz{2.01.2002 --- jaki dzień tygodnia}{środa}
  \zz{6.11.2038 --- jaki dzień tygodnia}{sobota}
  \zz{20.02.2011 --- jaki dzień tygodnia}{niedziela}
  \zz{7.07.2016 --- jaki dzień tygodnia}{czwartek}
  \zz{4.03.2041 --- jaki dzień tygodnia}{poniedziałek}
  \zz{3.08.2018 --- jaki dzień tygodnia}{piątek}
  \zz{4.10.2019 --- jaki dzień tygodnia}{piątek}
  \zz{23.11.2033 --- jaki dzień tygodnia}{środa}
  \zz{16.02.2047 --- jaki dzień tygodnia}{sobota}
  \zz{17.09.2037 --- jaki dzień tygodnia}{czwartek}
  \zz{27.10.2013 --- jaki dzień tygodnia}{niedziela}
  \zz{1.04.2018 --- jaki dzień tygodnia}{niedziela}
  \zz{1.04.2035 --- jaki dzień tygodnia}{niedziela}
  \zz{24.09.2002 --- jaki dzień tygodnia}{wtorek}
  \zz{22.07.2039 --- jaki dzień tygodnia}{piątek}
  \zz{23.06.2023 --- jaki dzień tygodnia}{piątek}
  \zz{4.07.2047 --- jaki dzień tygodnia}{czwartek}
  \zz{7.02.2002 --- jaki dzień tygodnia}{czwartek}
  \btnc
  \zz{13.10.2010 --- jaki dzień tygodnia}{środa}
  \zz{16.01.2045 --- jaki dzień tygodnia}{poniedziałek}
  \zz{10.08.2035 --- jaki dzień tygodnia}{piątek}
  \zz{12.04.2043 --- jaki dzień tygodnia}{niedziela}
  \zz{10.09.2049 --- jaki dzień tygodnia}{piątek}
  \zz{18.11.2010 --- jaki dzień tygodnia}{czwartek}
  \zz{14.02.2027 --- jaki dzień tygodnia}{niedziela}
  \zz{2.07.2028 --- jaki dzień tygodnia}{niedziela}
  \zz{5.04.2031 --- jaki dzień tygodnia}{sobota}
  \zz{4.11.2031 --- jaki dzień tygodnia}{wtorek}
  \zz{20.12.2019 --- jaki dzień tygodnia}{piątek}
  \zz{21.01.2012 --- jaki dzień tygodnia}{sobota}
  \zz{26.08.2040 --- jaki dzień tygodnia}{niedziela}
  \zz{12.09.2020 --- jaki dzień tygodnia}{sobota}
  \zz{18.04.2041 --- jaki dzień tygodnia}{czwartek}
  \zz{12.03.2011 --- jaki dzień tygodnia}{sobota}
  \zz{9.12.2005 --- jaki dzień tygodnia}{piątek}
  \zz{7.12.2015 --- jaki dzień tygodnia}{poniedziałek}
  \zz{6.09.2004 --- jaki dzień tygodnia}{poniedziałek}
  \zz{1.01.2025 --- jaki dzień tygodnia}{środa}
\end{zestaw}

\btexpect{86}\input{sets-adv/generated/86-procent-lata-b}
\btexpect{87}% ======================================================
%  GENERATED by tools/gen_sets.py -- do not edit.
%  Seed 2087; regenerate with `make sets`.
% ======================================================

\begin{zestaw}{Kombinatoryka III}{3}{8:00}{3}
  \z{$6!$}{720}
  \z{$10!$}{3\,628\,800}
  \z{$C(12, 2)$}{66}
  \z{$4!$}{24}
  \z{$P(13, 2)$}{156}
  \z{$9!$}{362\,880}
  \z{$C(4, 2)$}{6}
  \z{$C(12, 4)$}{495}
  \z{$P(8, 5)$}{6\,720}
  \z{$8!$}{40\,320}
  \z{$P(14, 2)$}{182}
  \z{$P(6, 2)$}{30}
  \z{$P(7, 4)$}{840}
  \z{$C(7, 4)$}{35}
  \btnc
  \z{$P(4, 3)$}{24}
  \z{$C(5, 2)$}{10}
  \z{$C(5, 3)$}{10}
  \z{$P(14, 5)$}{240\,240}
  \z{$C(14, 3)$}{364}
  \z{$P(10, 4)$}{5\,040}
  \z{$C(9, 4)$}{126}
  \z{$3!$}{6}
  \z{$P(13, 3)$}{1\,716}
  \z{$C(10, 3)$}{120}
  \z{$P(8, 4)$}{1\,680}
  \z{$C(8, 4)$}{70}
  \z{$P(11, 3)$}{990}
  \btnc
  \z{$C(13, 4)$}{715}
  \z{$P(8, 3)$}{336}
  \z{$P(11, 2)$}{110}
  \z{$C(7, 3)$}{35}
  \z{$C(10, 2)$}{45}
  \z{$P(9, 5)$}{15\,120}
  \z{$P(12, 5)$}{95\,040}
  \z{$P(10, 2)$}{90}
  \z{$P(13, 5)$}{154\,440}
  \z{$P(7, 3)$}{210}
  \z{$P(13, 4)$}{17\,160}
  \z{$C(11, 4)$}{330}
  \z{$7!$}{5\,040}
\end{zestaw}

\btexpect{88}% ======================================================
%  GENERATED by tools/gen_sets.py -- do not edit.
%  Seed 2088; regenerate with `make sets`.
% ======================================================

\begin{zestaw}{km/h i m/s III}{3}{6:00}{2}
  \z{$180$ km/h $\rightarrow$ m/s}{50}
  \z{$20$ m/s $\rightarrow$ km/h}{72}
  \z{$35$ m/s $\rightarrow$ km/h}{126}
  \z{$171$ km/h $\rightarrow$ m/s}{47,5}
  \z{$45$ km/h $\rightarrow$ m/s}{12,5}
  \z{$153$ km/h $\rightarrow$ m/s}{42,5}
  \z{$12,5$ m/s $\rightarrow$ km/h}{45}
  \z{$117$ km/h $\rightarrow$ m/s}{32,5}
  \z{$57,5$ m/s $\rightarrow$ km/h}{207}
  \z{$45$ m/s $\rightarrow$ km/h}{162}
  \z{$27$ km/h $\rightarrow$ m/s}{7,5}
  \z{$189$ km/h $\rightarrow$ m/s}{52,5}
  \z{$10$ m/s $\rightarrow$ km/h}{36}
  \z{$63$ km/h $\rightarrow$ m/s}{17,5}
  \z{$30$ m/s $\rightarrow$ km/h}{108}
  \z{$37,5$ m/s $\rightarrow$ km/h}{135}
  \z{$54$ km/h $\rightarrow$ m/s}{15}
  \z{$5$ m/s $\rightarrow$ km/h}{18}
  \z{$216$ km/h $\rightarrow$ m/s}{60}
  \z{$25$ m/s $\rightarrow$ km/h}{90}
  \btnc
  \z{$17,5$ m/s $\rightarrow$ km/h}{63}
  \z{$81$ km/h $\rightarrow$ m/s}{22,5}
  \z{$36$ km/h $\rightarrow$ m/s}{10}
  \z{$40$ m/s $\rightarrow$ km/h}{144}
  \z{$135$ km/h $\rightarrow$ m/s}{37,5}
  \z{$162$ km/h $\rightarrow$ m/s}{45}
  \z{$72$ km/h $\rightarrow$ m/s}{20}
  \z{$22,5$ m/s $\rightarrow$ km/h}{81}
  \z{$32,5$ m/s $\rightarrow$ km/h}{117}
  \z{$52,5$ m/s $\rightarrow$ km/h}{189}
  \z{$207$ km/h $\rightarrow$ m/s}{57,5}
  \z{$90$ km/h $\rightarrow$ m/s}{25}
  \z{$15$ m/s $\rightarrow$ km/h}{54}
  \z{$47,5$ m/s $\rightarrow$ km/h}{171}
  \z{$126$ km/h $\rightarrow$ m/s}{35}
  \z{$55$ m/s $\rightarrow$ km/h}{198}
  \z{$60$ m/s $\rightarrow$ km/h}{216}
  \z{$27,5$ m/s $\rightarrow$ km/h}{99}
  \z{$108$ km/h $\rightarrow$ m/s}{30}
  \z{$50$ m/s $\rightarrow$ km/h}{180}
\end{zestaw}

\btexpect{89}% ======================================================
%  GENERATED by tools/gen_sets.py -- do not edit.
%  Seed 2089; regenerate with `make sets`.
% ======================================================

\begin{zestaw}{Prędkość średnia III}{3}{8:40}{2}
  \zz{$60$ km w $45$ min --- km/h}{80}
  \zz{$125$ km w $1$ h $15$ min --- km/h}{100}
  \zz{$90$ km w $2$ h --- km/h}{45}
  \zz{$180$ km w $1$ h $30$ min --- km/h}{120}
  \zz{$125$ km w $2$ h $30$ min --- km/h}{50}
  \zz{$250$ km w $2$ h $30$ min --- km/h}{100}
  \zz{$100$ km w $1$ h $15$ min --- km/h}{80}
  \zz{$180$ km w $2$ h $15$ min --- km/h}{80}
  \zz{$150$ km w $1$ h $15$ min --- km/h}{120}
  \zz{$75$ km w $1$ h $15$ min --- km/h}{60}
  \zz{$20$ km w $30$ min --- km/h}{40}
  \zz{$150$ km w $2$ h --- km/h}{75}
  \zz{$160$ km w $2$ h --- km/h}{80}
  \zz{$45$ km w $45$ min --- km/h}{60}
  \zz{$150$ km w $1$ h $30$ min --- km/h}{100}
  \zz{$135$ km w $2$ h $15$ min --- km/h}{60}
  \zz{$175$ km w $1$ h $45$ min --- km/h}{100}
  \zz{$210$ km w $1$ h $45$ min --- km/h}{120}
  \zz{$200$ km w $2$ h $30$ min --- km/h}{80}
  \zz{$300$ km w $2$ h $30$ min --- km/h}{120}
  \btnc
  \zz{$30$ km w $45$ min --- km/h}{40}
  \zz{$100$ km w $2$ h $30$ min --- km/h}{40}
  \zz{$40$ km w $30$ min --- km/h}{80}
  \zz{$70$ km w $1$ h $45$ min --- km/h}{40}
  \zz{$240$ km w $2$ h --- km/h}{120}
  \zz{$35$ km w $30$ min --- km/h}{70}
  \zz{$105$ km w $1$ h $45$ min --- km/h}{60}
  \zz{$30$ km w $30$ min --- km/h}{60}
  \zz{$200$ km w $2$ h --- km/h}{100}
  \zz{$60$ km w $30$ min --- km/h}{120}
  \zz{$180$ km w $2$ h --- km/h}{90}
  \zz{$50$ km w $1$ h $15$ min --- km/h}{40}
  \zz{$270$ km w $2$ h $15$ min --- km/h}{120}
  \zz{$75$ km w $1$ h $30$ min --- km/h}{50}
  \zz{$120$ km w $1$ h $30$ min --- km/h}{80}
  \zz{$105$ km w $1$ h $30$ min --- km/h}{70}
  \zz{$225$ km w $2$ h $15$ min --- km/h}{100}
  \zz{$90$ km w $45$ min --- km/h}{120}
  \zz{$140$ km w $2$ h --- km/h}{70}
  \zz{$45$ km w $30$ min --- km/h}{90}
\end{zestaw}

\btexpect{90}% ======================================================
%  GENERATED by tools/gen_sets.py -- do not edit.
%  Seed 2090; regenerate with `make sets`.
% ======================================================

\begin{zestaw}{Skala mapy III}{3}{8:00}{2}
  \zz{skala $1:25\,000$, $11$ cm --- ile km}{2,75}
  \zz{skala $1:50\,000$, $9$ cm --- ile km}{4,5}
  \zz{skala $1:200\,000$, $2$ cm --- ile km}{4}
  \zz{skala $1:25\,000$, $8$ cm --- ile km}{2}
  \zz{skala $1:10\,000$, $5$ cm --- ile km}{0,5}
  \zz{skala $1:25\,000$, $9$ cm --- ile km}{2,25}
  \zz{skala $1:200\,000$, $10$ cm --- ile km}{20}
  \zz{skala $1:50\,000$, $4$ cm --- ile km}{2}
  \zz{skala $1:100\,000$, $3$ cm --- ile km}{3}
  \zz{skala $1:100\,000$, $7$ cm --- ile km}{7}
  \zz{skala $1:200\,000$, $9$ cm --- ile km}{18}
  \zz{skala $1:200\,000$, $8$ cm --- ile km}{16}
  \zz{skala $1:25\,000$, $10$ cm --- ile km}{2,5}
  \zz{skala $1:10\,000$, $12$ cm --- ile km}{1,2}
  \zz{skala $1:10\,000$, $8$ cm --- ile km}{0,8}
  \zz{skala $1:100\,000$, $2$ cm --- ile km}{2}
  \zz{skala $1:25\,000$, $2$ cm --- ile km}{0,5}
  \zz{skala $1:50\,000$, $12$ cm --- ile km}{6}
  \zz{skala $1:100\,000$, $10$ cm --- ile km}{10}
  \zz{skala $1:200\,000$, $3$ cm --- ile km}{6}
  \btnc
  \zz{skala $1:200\,000$, $7$ cm --- ile km}{14}
  \zz{skala $1:50\,000$, $6$ cm --- ile km}{3}
  \zz{skala $1:10\,000$, $10$ cm --- ile km}{1}
  \zz{skala $1:25\,000$, $12$ cm --- ile km}{3}
  \zz{skala $1:100\,000$, $5$ cm --- ile km}{5}
  \zz{skala $1:100\,000$, $8$ cm --- ile km}{8}
  \zz{skala $1:100\,000$, $4$ cm --- ile km}{4}
  \zz{skala $1:10\,000$, $3$ cm --- ile km}{0,3}
  \zz{skala $1:10\,000$, $6$ cm --- ile km}{0,6}
  \zz{skala $1:50\,000$, $11$ cm --- ile km}{5,5}
  \zz{skala $1:10\,000$, $7$ cm --- ile km}{0,7}
  \zz{skala $1:50\,000$, $2$ cm --- ile km}{1}
  \zz{skala $1:100\,000$, $6$ cm --- ile km}{6}
  \zz{skala $1:100\,000$, $9$ cm --- ile km}{9}
  \zz{skala $1:50\,000$, $7$ cm --- ile km}{3,5}
  \zz{skala $1:25\,000$, $7$ cm --- ile km}{1,75}
  \zz{skala $1:200\,000$, $6$ cm --- ile km}{12}
  \zz{skala $1:10\,000$, $4$ cm --- ile km}{0,4}
  \zz{skala $1:10\,000$, $11$ cm --- ile km}{1,1}
  \zz{skala $1:10\,000$, $2$ cm --- ile km}{0,2}
\end{zestaw}

\btexpect{91}% ======================================================
%  GENERATED by tools/gen_sets.py -- do not edit.
%  Seed 2091; regenerate with `make sets`.
% ======================================================

\begin{zestaw}{Mieszane — zastosowania}{3}{8:40}{2}
  \z{$C(10, 2)$}{45}
  \z{$4!$}{24}
  \z{$C(11, 3)$}{165}
  \z{$7!$}{5\,040}
  \z{$P(12, 4)$}{11\,880}
  \z{$P(6, 5)$}{720}
  \z{$3!$}{6}
  \z{$C(11, 4)$}{330}
  \z{$8!$}{40\,320}
  \z{$C(4, 2)$}{6}
  \z{$117$ km/h $\rightarrow$ m/s}{32,5}
  \z{$54$ km/h $\rightarrow$ m/s}{15}
  \z{$5$ m/s $\rightarrow$ km/h}{18}
  \z{$99$ km/h $\rightarrow$ m/s}{27,5}
  \z{$207$ km/h $\rightarrow$ m/s}{57,5}
  \z{$20$ m/s $\rightarrow$ km/h}{72}
  \z{$72$ km/h $\rightarrow$ m/s}{20}
  \z{$144$ km/h $\rightarrow$ m/s}{40}
  \z{$50$ m/s $\rightarrow$ km/h}{180}
  \z{$22,5$ m/s $\rightarrow$ km/h}{81}
  \btnc
  \z{$7^{5} \bmod 13$}{11}
  \z{$8^{5} \bmod 11$}{10}
  \z{$9^{6} \bmod 13$}{1}
  \z{$9^{3} \bmod 7$}{1}
  \z{$3^{3} \bmod 7$}{6}
  \z{$2^{4} \bmod 7$}{2}
  \z{$7^{4} \bmod 5$}{1}
  \z{$7^{6} \bmod 11$}{4}
  \z{$5^{6} \bmod 5$}{0}
  \z{$2^{6} \bmod 9$}{1}
  \z{$1/3 \div 2/3$}{1/2}
  \z{$8/12 \div 2/4$}{4/3}
  \z{$8/10 \div 4/8$}{8/5}
  \z{$1/9 \div 7/12$}{4/21}
  \z{$1/5 \div 5/10$}{2/5}
  \z{$2/4 \div 7/8$}{4/7}
  \z{$1/2 \div 5/6$}{3/5}
  \z{$7/10 \div 8/10$}{7/8}
  \z{$1/2 \div 1/2$}{1}
  \z{$1/2 \div 4/6$}{3/4}
\end{zestaw}

\btexpect{92}% ======================================================
%  GENERATED by tools/gen_sets.py -- do not edit.
%  Seed 2092; regenerate with `make sets`.
% ======================================================

\begin{zestaw}{Tam i z powrotem III}{3}{9:20}{2}
  \zz{tam $30$ km/h, z powrotem $150$ km/h --- średnia}{50}
  \zz{tam $144$ km/h, z powrotem $72$ km/h --- średnia}{96}
  \zz{tam $30$ km/h, z powrotem $60$ km/h --- średnia}{40}
  \zz{tam $60$ km/h, z powrotem $36$ km/h --- średnia}{45}
  \zz{tam $120$ km/h, z powrotem $60$ km/h --- średnia}{80}
  \zz{tam $120$ km/h, z powrotem $80$ km/h --- średnia}{96}
  \zz{tam $36$ km/h, z powrotem $45$ km/h --- średnia}{40}
  \zz{tam $36$ km/h, z powrotem $180$ km/h --- średnia}{60}
  \zz{tam $60$ km/h, z powrotem $180$ km/h --- średnia}{90}
  \zz{tam $200$ km/h, z powrotem $50$ km/h --- średnia}{80}
  \zz{tam $240$ km/h, z powrotem $48$ km/h --- średnia}{80}
  \zz{tam $80$ km/h, z powrotem $20$ km/h --- średnia}{32}
  \zz{tam $72$ km/h, z powrotem $90$ km/h --- średnia}{80}
  \zz{tam $120$ km/h, z powrotem $240$ km/h --- średnia}{160}
  \zz{tam $180$ km/h, z powrotem $90$ km/h --- średnia}{120}
  \zz{tam $80$ km/h, z powrotem $120$ km/h --- średnia}{96}
  \zz{tam $20$ km/h, z powrotem $60$ km/h --- średnia}{30}
  \zz{tam $72$ km/h, z powrotem $120$ km/h --- średnia}{90}
  \zz{tam $200$ km/h, z powrotem $120$ km/h --- średnia}{150}
  \zz{tam $24$ km/h, z powrotem $48$ km/h --- średnia}{32}
  \btnc
  \zz{tam $80$ km/h, z powrotem $48$ km/h --- średnia}{60}
  \zz{tam $90$ km/h, z powrotem $180$ km/h --- średnia}{120}
  \zz{tam $75$ km/h, z powrotem $50$ km/h --- średnia}{60}
  \zz{tam $120$ km/h, z powrotem $200$ km/h --- średnia}{150}
  \zz{tam $144$ km/h, z powrotem $48$ km/h --- średnia}{72}
  \zz{tam $144$ km/h, z powrotem $240$ km/h --- średnia}{180}
  \zz{tam $20$ km/h, z powrotem $180$ km/h --- średnia}{36}
  \zz{tam $20$ km/h, z powrotem $30$ km/h --- średnia}{24}
  \zz{tam $24$ km/h, z powrotem $72$ km/h --- średnia}{36}
  \zz{tam $120$ km/h, z powrotem $72$ km/h --- średnia}{90}
  \zz{tam $48$ km/h, z powrotem $144$ km/h --- średnia}{72}
  \zz{tam $36$ km/h, z powrotem $72$ km/h --- średnia}{48}
  \zz{tam $240$ km/h, z powrotem $120$ km/h --- średnia}{160}
  \zz{tam $180$ km/h, z powrotem $60$ km/h --- średnia}{90}
  \zz{tam $48$ km/h, z powrotem $240$ km/h --- średnia}{80}
  \zz{tam $30$ km/h, z powrotem $120$ km/h --- średnia}{48}
  \zz{tam $30$ km/h, z powrotem $20$ km/h --- średnia}{24}
  \zz{tam $50$ km/h, z powrotem $75$ km/h --- średnia}{60}
  \zz{tam $80$ km/h, z powrotem $240$ km/h --- średnia}{120}
  \zz{tam $30$ km/h, z powrotem $90$ km/h --- średnia}{45}
\end{zestaw}

\btexpect{93}% ======================================================
%  GENERATED by tools/gen_sets.py -- do not edit.
%  Seed 2093; regenerate with `make sets`.
% ======================================================

\begin{zestaw}{Dzień tygodnia III}{3}{10:40}{2}
  \zz{12.09.2027 --- jaki dzień tygodnia}{niedziela}
  \zz{21.03.2001 --- jaki dzień tygodnia}{środa}
  \zz{24.04.2025 --- jaki dzień tygodnia}{czwartek}
  \zz{27.03.2021 --- jaki dzień tygodnia}{sobota}
  \zz{18.03.2026 --- jaki dzień tygodnia}{środa}
  \zz{19.08.2034 --- jaki dzień tygodnia}{sobota}
  \zz{27.10.2012 --- jaki dzień tygodnia}{sobota}
  \zz{5.05.2000 --- jaki dzień tygodnia}{piątek}
  \zz{22.04.2032 --- jaki dzień tygodnia}{czwartek}
  \zz{2.09.2041 --- jaki dzień tygodnia}{poniedziałek}
  \zz{25.03.2049 --- jaki dzień tygodnia}{czwartek}
  \zz{23.01.2016 --- jaki dzień tygodnia}{sobota}
  \zz{25.12.2035 --- jaki dzień tygodnia}{wtorek}
  \zz{22.10.2019 --- jaki dzień tygodnia}{wtorek}
  \zz{20.11.2036 --- jaki dzień tygodnia}{czwartek}
  \zz{12.11.2028 --- jaki dzień tygodnia}{niedziela}
  \zz{27.04.2033 --- jaki dzień tygodnia}{środa}
  \zz{23.06.2017 --- jaki dzień tygodnia}{piątek}
  \zz{23.10.2026 --- jaki dzień tygodnia}{piątek}
  \zz{23.04.2032 --- jaki dzień tygodnia}{piątek}
  \btnc
  \zz{10.07.2033 --- jaki dzień tygodnia}{niedziela}
  \zz{8.04.2011 --- jaki dzień tygodnia}{piątek}
  \zz{3.08.2007 --- jaki dzień tygodnia}{piątek}
  \zz{10.05.2029 --- jaki dzień tygodnia}{czwartek}
  \zz{26.04.2008 --- jaki dzień tygodnia}{sobota}
  \zz{27.03.2013 --- jaki dzień tygodnia}{środa}
  \zz{14.07.2017 --- jaki dzień tygodnia}{piątek}
  \zz{18.11.2047 --- jaki dzień tygodnia}{poniedziałek}
  \zz{11.06.2047 --- jaki dzień tygodnia}{wtorek}
  \zz{2.10.2002 --- jaki dzień tygodnia}{środa}
  \zz{21.09.2037 --- jaki dzień tygodnia}{poniedziałek}
  \zz{24.04.2049 --- jaki dzień tygodnia}{sobota}
  \zz{26.10.2019 --- jaki dzień tygodnia}{sobota}
  \zz{17.09.2044 --- jaki dzień tygodnia}{sobota}
  \zz{3.09.2013 --- jaki dzień tygodnia}{wtorek}
  \zz{16.01.2047 --- jaki dzień tygodnia}{środa}
  \zz{3.09.2049 --- jaki dzień tygodnia}{piątek}
  \zz{9.03.2026 --- jaki dzień tygodnia}{poniedziałek}
  \zz{12.06.2045 --- jaki dzień tygodnia}{poniedziałek}
  \zz{10.04.2031 --- jaki dzień tygodnia}{czwartek}
\end{zestaw}

\btexpect{94}% ======================================================
%  GENERATED by tools/gen_sets.py -- do not edit.
%  Seed 2094; regenerate with `make sets`.
% ======================================================

\begin{zestaw}{Procent przez lata III}{3}{9:20}{2}
  \zz{$1\,000$ przy $25\%$ przez $2$ lata}{1\,562}
  \zz{$5\,000$ przy $20\%$ przez $3$ lata}{8\,640}
  \zz{$10\,000$ przy $10\%$ przez $3$ lata}{13\,310}
  \zz{$4\,000$ przy $50\%$ przez $2$ lata}{9\,000}
  \zz{$4\,000$ przy $20\%$ przez $2$ lata}{5\,760}
  \zz{$4\,000$ przy $25\%$ przez $3$ lata}{7\,812}
  \zz{$4\,000$ przy $20\%$ przez $3$ lata}{6\,912}
  \zz{$4\,000$ przy $50\%$ przez $3$ lata}{13\,500}
  \zz{$2\,000$ przy $20\%$ przez $3$ lata}{3\,456}
  \zz{$8\,000$ przy $10\%$ przez $2$ lata}{9\,680}
  \zz{$8\,000$ przy $25\%$ przez $2$ lata}{12\,500}
  \zz{$10\,000$ przy $50\%$ przez $2$ lata}{22\,500}
  \zz{$5\,000$ przy $10\%$ przez $3$ lata}{6\,655}
  \zz{$5\,000$ przy $25\%$ przez $3$ lata}{9\,765}
  \zz{$2\,000$ przy $10\%$ przez $2$ lata}{2\,420}
  \zz{$10\,000$ przy $50\%$ przez $3$ lata}{33\,750}
  \zz{$10\,000$ przy $10\%$ przez $2$ lata}{12\,100}
  \zz{$1\,000$ przy $25\%$ przez $3$ lata}{1\,952}
  \zz{$2\,000$ przy $50\%$ przez $2$ lata}{4\,500}
  \zz{$10\,000$ przy $25\%$ przez $3$ lata}{19\,531}
  \btnc
  \zz{$5\,000$ przy $10\%$ przez $2$ lata}{6\,050}
  \zz{$8\,000$ przy $20\%$ przez $2$ lata}{11\,520}
  \zz{$5\,000$ przy $25\%$ przez $2$ lata}{7\,812}
  \zz{$1\,000$ przy $10\%$ przez $3$ lata}{1\,331}
  \zz{$8\,000$ przy $50\%$ przez $2$ lata}{18\,000}
  \zz{$8\,000$ przy $50\%$ przez $3$ lata}{27\,000}
  \zz{$10\,000$ przy $20\%$ przez $2$ lata}{14\,400}
  \zz{$8\,000$ przy $20\%$ przez $3$ lata}{13\,824}
  \zz{$5\,000$ przy $50\%$ przez $2$ lata}{11\,250}
  \zz{$4\,000$ przy $10\%$ przez $3$ lata}{5\,324}
  \zz{$2\,000$ przy $20\%$ przez $2$ lata}{2\,880}
  \zz{$10\,000$ przy $20\%$ przez $3$ lata}{17\,280}
  \zz{$2\,000$ przy $10\%$ przez $3$ lata}{2\,662}
  \zz{$8\,000$ przy $25\%$ przez $3$ lata}{15\,625}
  \zz{$5\,000$ przy $20\%$ przez $2$ lata}{7\,200}
  \zz{$5\,000$ przy $50\%$ przez $3$ lata}{16\,875}
  \zz{$8\,000$ przy $10\%$ przez $3$ lata}{10\,648}
  \zz{$4\,000$ przy $25\%$ przez $2$ lata}{6\,250}
  \zz{$1\,000$ przy $50\%$ przez $3$ lata}{3\,375}
  \zz{$2\,000$ przy $50\%$ przez $3$ lata}{6\,750}
\end{zestaw}

\btexpect{95}% ======================================================
%  GENERATED by tools/gen_sets.py -- do not edit.
%  Seed 2095; regenerate with `make sets`.
% ======================================================

\begin{zestaw}{Kombinatoryka IV}{3}{8:00}{3}
  \z{$P(8, 5)$}{6\,720}
  \z{$6!$}{720}
  \z{$8!$}{40\,320}
  \z{$P(14, 4)$}{24\,024}
  \z{$P(11, 5)$}{55\,440}
  \z{$C(13, 3)$}{286}
  \z{$C(11, 3)$}{165}
  \z{$C(12, 4)$}{495}
  \z{$C(5, 3)$}{10}
  \z{$9!$}{362\,880}
  \z{$C(14, 4)$}{1\,001}
  \z{$P(8, 4)$}{1\,680}
  \z{$C(12, 3)$}{220}
  \z{$10!$}{3\,628\,800}
  \btnc
  \z{$5!$}{120}
  \z{$P(6, 3)$}{120}
  \z{$C(8, 2)$}{28}
  \z{$P(5, 2)$}{20}
  \z{$C(9, 2)$}{36}
  \z{$4!$}{24}
  \z{$C(14, 3)$}{364}
  \z{$P(12, 4)$}{11\,880}
  \z{$P(13, 4)$}{17\,160}
  \z{$C(8, 5)$}{56}
  \z{$P(7, 2)$}{42}
  \z{$P(10, 3)$}{720}
  \z{$7!$}{5\,040}
  \btnc
  \z{$P(5, 4)$}{120}
  \z{$C(13, 5)$}{1\,287}
  \z{$P(6, 4)$}{360}
  \z{$P(9, 4)$}{3\,024}
  \z{$P(8, 3)$}{336}
  \z{$C(7, 5)$}{21}
  \z{$P(8, 2)$}{56}
  \z{$C(4, 3)$}{4}
  \z{$P(14, 2)$}{182}
  \z{$C(5, 2)$}{10}
  \z{$P(14, 3)$}{2\,184}
  \z{$3!$}{6}
  \z{$P(10, 2)$}{90}
\end{zestaw}

\btexpect{96}% ======================================================
%  GENERATED by tools/gen_sets.py -- do not edit.
%  Seed 2096; regenerate with `make sets`.
% ======================================================

\begin{zestaw}{km/h i m/s IV}{3}{6:00}{2}
  \z{$5$ m/s $\rightarrow$ km/h}{18}
  \z{$55$ m/s $\rightarrow$ km/h}{198}
  \z{$99$ km/h $\rightarrow$ m/s}{27,5}
  \z{$171$ km/h $\rightarrow$ m/s}{47,5}
  \z{$162$ km/h $\rightarrow$ m/s}{45}
  \z{$30$ m/s $\rightarrow$ km/h}{108}
  \z{$27,5$ m/s $\rightarrow$ km/h}{99}
  \z{$7,5$ m/s $\rightarrow$ km/h}{27}
  \z{$189$ km/h $\rightarrow$ m/s}{52,5}
  \z{$207$ km/h $\rightarrow$ m/s}{57,5}
  \z{$54$ km/h $\rightarrow$ m/s}{15}
  \z{$50$ m/s $\rightarrow$ km/h}{180}
  \z{$72$ km/h $\rightarrow$ m/s}{20}
  \z{$135$ km/h $\rightarrow$ m/s}{37,5}
  \z{$15$ m/s $\rightarrow$ km/h}{54}
  \z{$25$ m/s $\rightarrow$ km/h}{90}
  \z{$10$ m/s $\rightarrow$ km/h}{36}
  \z{$81$ km/h $\rightarrow$ m/s}{22,5}
  \z{$198$ km/h $\rightarrow$ m/s}{55}
  \z{$90$ km/h $\rightarrow$ m/s}{25}
  \btnc
  \z{$27$ km/h $\rightarrow$ m/s}{7,5}
  \z{$45$ m/s $\rightarrow$ km/h}{162}
  \z{$180$ km/h $\rightarrow$ m/s}{50}
  \z{$52,5$ m/s $\rightarrow$ km/h}{189}
  \z{$108$ km/h $\rightarrow$ m/s}{30}
  \z{$117$ km/h $\rightarrow$ m/s}{32,5}
  \z{$36$ km/h $\rightarrow$ m/s}{10}
  \z{$35$ m/s $\rightarrow$ km/h}{126}
  \z{$37,5$ m/s $\rightarrow$ km/h}{135}
  \z{$63$ km/h $\rightarrow$ m/s}{17,5}
  \z{$47,5$ m/s $\rightarrow$ km/h}{171}
  \z{$126$ km/h $\rightarrow$ m/s}{35}
  \z{$42,5$ m/s $\rightarrow$ km/h}{153}
  \z{$144$ km/h $\rightarrow$ m/s}{40}
  \z{$57,5$ m/s $\rightarrow$ km/h}{207}
  \z{$40$ m/s $\rightarrow$ km/h}{144}
  \z{$45$ km/h $\rightarrow$ m/s}{12,5}
  \z{$22,5$ m/s $\rightarrow$ km/h}{81}
  \z{$32,5$ m/s $\rightarrow$ km/h}{117}
  \z{$18$ km/h $\rightarrow$ m/s}{5}
\end{zestaw}



\chapter{Łamigłówki}

Dwa zestawy pisane ręcznie. W pierwszym każde zadanie ma krótką odpowiedź,
którą łatwo sprawdzić, i drogę do niej, która nie jest pierwsza z brzegu.
W drugim \emph{podpowiedź jest treścią} — wynik sprawdzisz w dziesięć sekund,
chodzi o skrót.

% GENERATED by tools/gen_sets.py -- do not edit.
\btexpect{97}% ======================================================
%  Zagadki -- poziom zaawansowany. Hand-written, for the reason the basic
%  volume's are: a word trap is a joke and a joke has an author.
%
%  These are harder than the basic volume's in a specific way rather than a
%  general one -- each has a short answer that is easy to check and a route to
%  it that is not the first one anybody tries.
% ======================================================

\begin{zestaw}{Zagadki — trudniejsze}{3}{20:00}{1}
  \zz{Trzy skrzynki: „jabłka”, „pomarańcze”, „mieszane”. \;/\; Każda etykieta jest błędna. \;/\; Ile owoców trzeba wyjąć, żeby ustalić zawartość wszystkich trzech?}{1 \hfill {\itshape\footnotesize Ze skrzynki „mieszane” — skoro etykieta kłamie, jest jednorodna.}}
  \zz{Dwa sznurki, każdy pali się dokładnie godzinę, \;/\; ale nierównomiernie. \;/\; Jak odmierzyć 45 minut?}{Zapal I z obu końców i II z jednego \hfill {\itshape\footnotesize I gaśnie po 30 min; wtedy podpal drugi koniec II — zostanie mu 15.}}
  \zz{Osiem monet, jedna cięższa, waga szalkowa bez odważników. \;/\; Ile ważeń wystarczy w najgorszym razie?}{2 \hfill {\itshape\footnotesize 3 na 3, potem 1 na 1 — trzy grupy, nie dwie.}}
  \zz{Dwanaście monet, jedna inna — nie wiadomo, cięższa czy lżejsza. \;/\; Ile ważeń wystarczy?}{3 \hfill {\itshape\footnotesize Trzy ważenia rozróżniają $3^3 = 27$ przypadków; tu jest ich 24.}}
  \zz{Cztery osoby przechodzą nocą przez most, jedna latarka, \;/\; idą po dwie, w tempie wolniejszego. \;/\; Czasy: 1, 2, 5, 10 minut. Najkrótszy łączny czas?}{17 minut \hfill {\itshape\footnotesize Dwaj najwolniejsi idą razem — raz.}}
  \zz{Dzbanki na 3 i 5 litrów, kran bez podziałki. \;/\; Jak odmierzyć dokładnie 4 litry?}{Napełnij 5, przelej do 3, zostaną 2 \hfill {\itshape\footnotesize Opróżnij trójkę, wlej te 2, dopełnij piątkę i dolej 1 do trójki.}}
  \zz{Ciąg: $1$, $11$, $21$, $1211$, $111221$, \ldots \;/\; Jaki jest następny wyraz?}{312211 \hfill {\itshape\footnotesize Każdy wyraz opisuje poprzedni: „trzy jedynki, dwie dwójki, jedna jedynka”.}}
  \zz{Iloma zerami kończy się $100!$?}{24 \hfill {\itshape\footnotesize Tyle, ile piątek w rozkładzie: $20 + 4$.}}
  \zz{Ile liczb od 1 do 1000 zawiera cyfrę 7?}{271 \hfill {\itshape\footnotesize Łatwiej policzyć te bez siódemki: $9 \times 9 \times 9 = 729$.}}
  \zz{Ile razy na dobę wskazówki zegara \;/\; tworzą kąt prosty?}{44 \hfill {\itshape\footnotesize 22 razy na 12 godzin — nie 24, bo nakładanie się „zjada” dwa.}}
  \zz{Mam tyle samo braci co sióstr. \;/\; Moja siostra ma dwa razy więcej braci niż sióstr. \;/\; Ilu nas jest?}{7 \hfill {\itshape\footnotesize 4 braci i 3 siostry; ja jestem bratem.}}
  \zz{Sznur opina Ziemię po równiku. \;/\; Dokładasz do niego 1 metr. \;/\; Jak wysoko uniesie się nad powierzchnię?}{Około 16 cm \hfill {\itshape\footnotesize $1/(2\pi)$ — promień Ziemi się skraca.}}
  \zz{Siatka $3 \times 3$ kwadratów. \;/\; Ile jest dróg z lewego górnego rogu do prawego dolnego, \;/\; jeśli wolno iść tylko w prawo i w dół?}{20 \hfill {\itshape\footnotesize Sześć kroków, z czego trzy w prawo: $C(6,3)$.}}
  \zz{Pierwszą połowę trasy jedziesz 60 km/h. \;/\; Jak szybko musisz jechać drugą połowę, \;/\; żeby średnia wyszła 120 km/h?}{Nie da się \hfill {\itshape\footnotesize Na samą pierwszą połowę zużyłeś już cały czas przewidziany na całość.}}
  \zz{W pokoju jest $n$ osób. \;/\; Przy jakim najmniejszym $n$ szansa, że dwie mają urodziny tego samego dnia, \;/\; przekracza połowę?}{23 \hfill {\itshape\footnotesize Liczy się liczba par, a nie osób: $C(23,2) = 253$.}}
  \zz{Kartka złożona na pół 50 razy \;/\; miałaby grubość rzędu?}{Odległości do Słońca \hfill {\itshape\footnotesize $2^{50}$ razy $0{,}1$ mm to około $10^{8}$ km.}}
\end{zestaw}

\btexpect{98}% ======================================================
%  Triki zaawansowane -- hand-written.
%
%  Here the HINT is the content: the answer takes ten seconds to verify and the
%  shortcut is the thing worth owning. That is the same rule the basic volume's
%  trick set states, one level of difficulty further on.
% ======================================================

\begin{zestaw}{Triki zaawansowane}{3}{12:00}{2}
  \zz{$997 \times 998$}{995\,006 \hfill {\itshape\footnotesize Obie blisko 1000: $997-2$ \,|\, $3 \times 2$ na trzech miejscach.}}
  \zz{$87 \times 93$}{8\,091 \hfill {\itshape\footnotesize $90^2 - 3^2$.}}
  \zz{$105^2$}{11\,025 \hfill {\itshape\footnotesize $10 \times 11$, dopisz $25$.}}
  \zz{$76 \times 74$}{5\,624 \hfill {\itshape\footnotesize Te same dziesiątki, jedności sumują się do 10: $7 \times 8$ \,|\, $6 \times 4$.}}
  \zz{$1\,001 \times 437$}{437\,437 \hfill {\itshape\footnotesize $1001 = 7 \times 11 \times 13$ — powtarza trzycyfrową liczbę.}}
  \zz{$999 \times 999$}{998\,001 \hfill {\itshape\footnotesize $(1000-1)^2 = 10^6 - 2000 + 1$.}}
  \zz{$25 \times 32$}{800 \hfill {\itshape\footnotesize $\times 100$, potem $:4$.}}
  \zz{$64 \times 125$}{8\,000 \hfill {\itshape\footnotesize $125 = 1000:8$.}}
  \btnc
  \zz{Pierwiastek z $5\,329$}{73 \hfill {\itshape\footnotesize Kończy się na 9, więc pierwiastek na 3 lub 7; leży między 70 a 80.}}
  \zz{$19 \times 21$}{399 \hfill {\itshape\footnotesize $20^2 - 1$.}}
  \zz{$111 \times 111$}{12\,321 \hfill {\itshape\footnotesize Palindrom — rośnie do liczby jedynek i wraca.}}
  \zz{$1/7$ dziesiętnie, sześć miejsc}{0,142857 \hfill {\itshape\footnotesize Ten sam sześciocyfrowy cykl obraca się dla $2/7$, $3/7$ i reszty.}}
  \zz{$7/8$ na procent}{87,5\% \hfill {\itshape\footnotesize $1/8 = 12{,}5\%$, więc $7/8$ to $100 - 12{,}5$.}}
  \zz{$2^{16}$}{65\,536 \hfill {\itshape\footnotesize $2^{10} \times 2^{6}$, czyli $1024 \times 64$.}}
  \zz{Czy $4\,321 \times 1\,234 = 5\,332\,124$? \;/\; Sprawdź przez dziewiątki.}{Nie \hfill {\itshape\footnotesize Cyfry kontrolne: $1 \times 1 = 1$, a wynik daje 2.}}
  \zz{$48 \times 15$}{720 \hfill {\itshape\footnotesize $\times 10$ plus połowa tego: $480 + 240$.}}
\end{zestaw}



\chapter{Pomiary kontrolne}

Cztery zestawy, po jednym na blok, i jedyne w tej książce zrobione po to, żeby
wracać do nich wielokrotnie. Każdy ma pod spodem pięć wierszy na pomiar.

Plan wywołuje je co siódmy dzień, przez cztery tygodnie każdy. Poza planem: rób
ten ze swojego bloku raz w tygodniu i nie porównuj wyniku z niczym innym.

% GENERATED by tools/gen_sets.py -- do not edit.
\btexpect{99}% ======================================================
%  GENERATED by tools/gen_sets.py -- do not edit.
%  Seed 2097; regenerate with `make sets`.
% ======================================================

\begin{zestaw}[5]{Pomiar kontrolny I}{2}{6:40}{3}
  \z{$11010_2$}{26}
  \z{$104$ dwójkowo}{1101000}
  \z{$1001001_2$}{73}
  \z{$100$ dwójkowo}{1100100}
  \z{$125$ dwójkowo}{1111101}
  \z{$209$ dwójkowo}{11010001}
  \z{$1111001_2$}{121}
  \z{$165$ dwójkowo}{10100101}
  \z{$828 \bmod 11$}{3}
  \z{$411 \bmod 4$}{3}
  \z{$316 \bmod 6$}{4}
  \z{$935 \bmod 3$}{2}
  \z{$540 \bmod 4$}{0}
  \z{$389 \bmod 6$}{5}
  \btnc
  \z{$218 \bmod 4$}{2}
  \z{$370 \bmod 11$}{7}
  \z{$203$ szesnastkowo}{CB}
  \z{$244$ szesnastkowo}{F4}
  \z{$144$ szesnastkowo}{90}
  \z{$\mathrm{D1}_{16}$}{209}
  \z{$58$ szesnastkowo}{3A}
  \z{$68$ szesnastkowo}{44}
  \z{$204$ szesnastkowo}{CC}
  \z{$112$ szesnastkowo}{70}
  \z{$226\,395$ --- cyfra kontrolna}{9}
  \z{$324\,240$ --- cyfra kontrolna}{6}
  \z{$245\,039$ --- cyfra kontrolna}{5}
  \btnc
  \z{$786\,490$ --- cyfra kontrolna}{7}
  \z{$206\,791$ --- cyfra kontrolna}{7}
  \z{$740\,400$ --- cyfra kontrolna}{6}
  \z{$7\,595$ --- cyfra kontrolna}{8}
  \z{$381\,066$ --- cyfra kontrolna}{6}
  \z{$2^{4} \bmod 9$}{7}
  \z{$2^{2} \bmod 13$}{4}
  \z{$8^{4} \bmod 5$}{1}
  \z{$3^{6} \bmod 5$}{4}
  \z{$7^{5} \bmod 13$}{11}
  \z{$6^{5} \bmod 7$}{6}
  \z{$7^{3} \bmod 11$}{2}
  \z{$4^{4} \bmod 9$}{4}
\end{zestaw}

\btexpect{100}% ======================================================
%  GENERATED by tools/gen_sets.py -- do not edit.
%  Seed 2098; regenerate with `make sets`.
% ======================================================

\begin{zestaw}[5]{Pomiar kontrolny II}{2}{7:20}{3}
  \z{$144$ na czynniki}{$2^{4} \cdot 3^{2}$}
  \z{$400$ na czynniki}{$2^{4} \cdot 5^{2}$}
  \z{$335$ na czynniki}{$5 \cdot 67$}
  \z{$327$ na czynniki}{$3 \cdot 109$}
  \z{$150$ na czynniki}{$2 \cdot 3 \cdot 5^{2}$}
  \z{$386$ na czynniki}{$2 \cdot 193$}
  \z{$220$ na czynniki}{$2^{2} \cdot 5 \cdot 11$}
  \z{$300$ na czynniki}{$2^{2} \cdot 3 \cdot 5^{2}$}
  \z{$129$ na czynniki}{$3 \cdot 43$}
  \z{$200$ na czynniki}{$2^{3} \cdot 5^{2}$}
  \z{NWD$(261, 75)$}{3}
  \z{NWD$(738, 243)$}{9}
  \z{NWD$(299, 949)$}{13}
  \z{NWD$(112, 511)$}{7}
  \btnc
  \z{NWD$(684, 600)$}{12}
  \z{NWD$(282, 240)$}{6}
  \z{NWD$(741, 1\,634)$}{19}
  \z{NWD$(33, 138)$}{3}
  \z{NWD$(517, 616)$}{11}
  \z{NWD$(418, 649)$}{11}
  \z{ile dzielników ma $96$}{12}
  \z{ile dzielników ma $294$}{12}
  \z{ile dzielników ma $352$}{12}
  \z{ile dzielników ma $213$}{4}
  \z{ile dzielników ma $132$}{12}
  \z{ile dzielników ma $301$}{4}
  \z{ile dzielników ma $382$}{4}
  \btnc
  \z{ile dzielników ma $112$}{10}
  \z{ile dzielników ma $312$}{16}
  \z{ile dzielników ma $187$}{4}
  \z{$2 + 4 + \ldots + 102$}{2\,652}
  \z{$1 + 3 + \ldots + 45$}{529}
  \z{$1 + 2 + \ldots + 47$}{1\,128}
  \z{$1 + 3 + \ldots + 107$}{2\,916}
  \z{$2 + 4 + \ldots + 68$}{1\,190}
  \z{$1 + 3 + \ldots + 17$}{81}
  \z{$1 + 3 + \ldots + 69$}{1\,225}
  \z{$2 + 4 + \ldots + 60$}{930}
  \z{$2 + 4 + \ldots + 74$}{1\,406}
  \z{$2 + 4 + \ldots + 118$}{3\,540}
\end{zestaw}

\btexpect{101}% ======================================================
%  GENERATED by tools/gen_sets.py -- do not edit.
%  Seed 2099; regenerate with `make sets`.
% ======================================================

\begin{zestaw}[5]{Pomiar kontrolny III}{3}{7:20}{3}
  \z{$7/9 \times 3/5$}{7/15}
  \z{$1/3 \times 1/2$}{1/6}
  \z{$4/9 \times 10/12$}{10/27}
  \z{$1/6 \times 7/9$}{7/54}
  \z{$11/12 \times 4/6$}{11/18}
  \z{$2/12 \times 4/6$}{1/9}
  \z{$6/8 \times 8/9$}{2/3}
  \z{$1/2 \times 3/9$}{1/6}
  \z{$4/5 \div 1/3$}{12/5}
  \z{$1/3 \div 1/4$}{4/3}
  \z{$7/12 \div 2/3$}{7/8}
  \z{$7/8 \div 3/6$}{7/4}
  \z{$1/2 \div 4/9$}{9/8}
  \z{$1/5 \div 4/12$}{3/5}
  \btnc
  \z{$3/10 \div 7/10$}{3/7}
  \z{$3/8 \div 1/6$}{9/4}
  \z{$2^{-2}$}{1/4}
  \z{$8^{-4}$}{1/4096}
  \z{$5^{-3}$}{1/125}
  \z{$3^{-4}$}{1/81}
  \z{$10^{-1}$}{1/10}
  \z{$10^{0}$}{1}
  \z{$6^{-2}$}{1/36}
  \z{$3^{-3}$}{1/27}
  \z{$\log_{7} 16\,807$}{5}
  \z{$\log_{8} 32\,768$}{5}
  \z{$\log_{6} 7\,776$}{5}
  \btnc
  \z{$\log_{9} 59\,049$}{5}
  \z{$\log_{7} 49$}{2}
  \z{$\log_{5} 15\,625$}{6}
  \z{$\log_{8} 512$}{3}
  \z{$\log_{3} 59\,049$}{10}
  \z{$7/8$ dziesiętnie}{0,875}
  \z{$7/10$ dziesiętnie}{0,7}
  \z{$1/2$ dziesiętnie}{0,5}
  \z{$33/50$ dziesiętnie}{0,66}
  \z{$7/20$ dziesiętnie}{0,35}
  \z{$3/4$ dziesiętnie}{0,75}
  \z{$1/4$ dziesiętnie}{0,25}
  \z{$34/50$ dziesiętnie}{0,68}
\end{zestaw}

\btexpect{102}% ======================================================
%  GENERATED by tools/gen_sets.py -- do not edit.
%  Seed 2100; regenerate with `make sets`.
% ======================================================

\begin{zestaw}[5]{Pomiar kontrolny IV}{3}{8:40}{2}
  \z{$C(11, 2)$}{55}
  \z{$C(5, 4)$}{5}
  \z{$C(9, 2)$}{36}
  \z{$6!$}{720}
  \z{$9!$}{362\,880}
  \z{$C(8, 2)$}{28}
  \z{$P(14, 4)$}{24\,024}
  \z{$C(14, 4)$}{1\,001}
  \z{$C(5, 3)$}{10}
  \z{$P(11, 4)$}{7\,920}
  \z{$27,5$ m/s $\rightarrow$ km/h}{99}
  \z{$42,5$ m/s $\rightarrow$ km/h}{153}
  \z{$50$ m/s $\rightarrow$ km/h}{180}
  \z{$40$ m/s $\rightarrow$ km/h}{144}
  \z{$90$ km/h $\rightarrow$ m/s}{25}
  \z{$36$ km/h $\rightarrow$ m/s}{10}
  \z{$162$ km/h $\rightarrow$ m/s}{45}
  \z{$135$ km/h $\rightarrow$ m/s}{37,5}
  \z{$126$ km/h $\rightarrow$ m/s}{35}
  \z{$117$ km/h $\rightarrow$ m/s}{32,5}
  \btnc
  \z{$128^{2/7}$}{4}
  \z{$64^{3/6}$}{8}
  \z{$27^{1/3}$}{3}
  \z{$4^{2/2}$}{4}
  \z{$144^{1/2}$}{12}
  \z{$8^{3/3}$}{8}
  \z{$25^{2/2}$}{25}
  \z{$81^{3/2}$}{729}
  \z{$27^{2/3}$}{9}
  \z{$64^{2/6}$}{4}
  \z{$2 \cdot 10^{5} \times 6 \cdot 10^{2}$}{$1{,}2 \cdot 10^{8}$}
  \z{$2 \cdot 10^{3} \times 8 \cdot 10^{3}$}{$1{,}6 \cdot 10^{7}$}
  \z{$8 \cdot 10^{8} \times 2 \cdot 10^{5}$}{$1{,}6 \cdot 10^{14}$}
  \z{$4 \cdot 10^{7} \times 7 \cdot 10^{3}$}{$2{,}8 \cdot 10^{11}$}
  \z{$6 \cdot 10^{2} \times 6 \cdot 10^{5}$}{$3{,}6 \cdot 10^{8}$}
  \z{$3 \cdot 10^{6} \times 2 \cdot 10^{5}$}{$6 \cdot 10^{11}$}
  \z{$3 \cdot 10^{3} \times 9 \cdot 10^{4}$}{$2{,}7 \cdot 10^{8}$}
  \z{$4 \cdot 10^{2} \times 4 \cdot 10^{2}$}{$1{,}6 \cdot 10^{5}$}
  \z{$9 \cdot 10^{6} \times 8 \cdot 10^{7}$}{$7{,}2 \cdot 10^{14}$}
  \z{$6 \cdot 10^{2} \times 3 \cdot 10^{6}$}{$1{,}8 \cdot 10^{9}$}
\end{zestaw}


